% =============================================================================
% paper/lit_review.tex
% Section 2: Related Literature  (+ Expected Dynamics, + Forward Research Agenda)
%
% Stand-alone section file — % =============================================================================
% paper/lit_review.tex
% Section 2: Related Literature  (+ Expected Dynamics, + Forward Research Agenda)
%
% Stand-alone section file — % =============================================================================
% paper/lit_review.tex
% Section 2: Related Literature  (+ Expected Dynamics, + Forward Research Agenda)
%
% Stand-alone section file — % =============================================================================
% paper/lit_review.tex
% Section 2: Related Literature  (+ Expected Dynamics, + Forward Research Agenda)
%
% Stand-alone section file — \input{lit_review} into the main paper.
% Cite keys follow the convention author(s)year; matching @article / @book /
% @techreport entries live in Bibliography_base.bib (86 entries as of 2026-06-05).
%
% STRUCTURE
%   2   Related Literature
%     2.1  Economic complexity, capability structure, and diversification theory
%     2.2  Capability theory, corporate coherence, and the direction of adjustment
%     2.3  Trade policy shocks and firm adjustment
%     2.4  Strategic dependencies and geopolitical trade shocks
%     2.5  Identification
%     2.6  Contribution
%   3   Expected Dynamics: What the Literature Predicts We Should Observe
%     3.1  The out-of-block margin (d_out)
%     3.2  The in-block margin (d_in) and retrenchment
%     3.3  Share_core and the direction of adjustment
%     3.4  Heterogeneity and the strategic-dependence moderator
%     3.5  Synthesis: hypotheses, signs, and magnitude benchmarks  [2 tables]
%   4   Open Questions and a Forward Research Agenda
%     4.1  The gaps this paper fills
%     4.2  Beyond this paper: spin-off research questions
%
% NOTE: tables use plain \hline tabular (no booktabs dependency) so the file
% compiles stand-alone via lit_review_test.tex and when \input into the paper.
% =============================================================================

\section{Related Literature}
\label{sec:lit}

This paper sits at the intersection of four bodies of work: the economic-complexity
literature on firm-level capability structure and the theory of diversification; the
capability-based theory of the firm and evolutionary economics; the empirical trade
literature on the effects of tariff shocks on firm outcomes; and the emerging literature
on strategic dependencies, geopolitical risk, and the weaponisation of economic
interdependence. A fifth, methodological strand covers identification in shift-share
designs with heterogeneous treatment effects. The unifying thread across these
literatures is a single dynamic: an external shock that degrades the value of a firm's
existing activities forces a \emph{reconfiguration} of its productive knowledge, and
the literatures disagree on whether that reconfiguration deepens a coherent core or
pushes outward into new capability bundles. We discuss each stream in turn, emphasising
what each predicts about the \emph{direction} and \emph{persistence} of adjustment, and
close with a statement of our contribution.

% ---------------------------------------------------------------------------
\subsection{Economic Complexity, Capability Structure, and Diversification Theory}
\label{subsec:lit_efc}
% ---------------------------------------------------------------------------

The product-space framework of \citet{hausmann2007} and \citet{hidalgo2009}
established that countries export products in patterns that are far from random:
goods that require similar underlying capabilities tend to be co-exported, generating
a proximity structure in the product space that conditions the path of structural
transformation. \citet{hidalgo2011} formalise this intuition through a bipartite
capabilities model in which countries and products are represented as the two node
sets of a bipartite graph, and the adjacency matrix $M_{cp}=1$ whenever country $c$
has revealed comparative advantage (RCA) in product $p$. The key result is that
diversification is convex in the capability stock: countries with few capabilities
earn near-zero marginal returns from acquiring any single new one, generating the
low-income quiescence trap that has been documented empirically across the
income distribution \citep{hausmann2007}. This non-linearity is fundamental to the
agenda of using complexity tools in a causal setting: treatment-effect heterogeneity
by capability endowment is an architectural feature of the framework, not an
add-on.

The Fitness-Complexity (FC) algorithm of \citet{cristelli2013} provides the
operational complexity metric adopted in the subsequent firm-level literature. Rather
than the linear iterative procedure of \citet{hidalgo2009}, FC defines country fitness
as the sum of product complexities and product complexity as the harmonic mean of
exporter fitness---a bounded operator that uniquely rewards diversification and
assigns complexity to a product according to the least capable exporter that ships it,
giving the algorithm its asymmetric censoring property. The algorithm converges to a
unique fixed point and, as shown by \citet{cristelli2017}, generates a selective
predictability scheme in which high-fitness countries follow laminar, predictable
growth trajectories while low-fitness countries exhibit chaotic dynamics---a property
that has direct implications for heterogeneous treatment responses in our setting:
firms embedded in high-complexity capability clusters should display more systematic
responses to tariff shocks than firms operating at the periphery of the product
space. \citet{angelini2017} characterise the long-run dynamics of products in the complexity
plane, showing that they drift toward a stationary configuration so that
shock-induced dislocations are transient and mean-reverting---which is why we treat
product complexity as a well-defined pre-shock control and read post-shock movements
against a fixed baseline.
\citet{stojkoski2025} extend the framework to an optimisation context,
showing that the ECI-growth relationship is stronger at lower income levels and that
relatedness is one of several predictors of product-entry, a finding that validates
the multi-dimensional block structure we exploit over simpler relatedness-only
measures. \citet{koch2021} sounds a useful caution about what such metrics measure: the
complexity--growth link is in fact \emph{stronger} when complexity is computed on
value-added rather than gross exports ($R^2 = 0.63$ versus $0.31$), so gross-flow
complexity blurs the connection to realised value capture. Since our firm-customs
measures are necessarily built on gross flows, we read them as capturing capability
\emph{structure} rather than realised value capture, which is precisely why we treat
the block partition as a measure of capability architecture rather than of performance.

Three strands sharpen what ``relatedness'' and ``diversification'' actually mean for
the direction of adjustment. First, the related-variety tradition of
\citet{frenken2007} distinguishes \emph{related} variety (entry into activities that
share capabilities, which drives regional \emph{employment} growth through
within-sector Jacobs spillovers---though, tellingly, not productivity growth) from
\emph{unrelated} variety (which provides a portfolio hedge that lowers unemployment
under asymmetric sectoral shocks). This is the regional-economics antecedent of our
$d_{\rm in}$ versus $d_{\rm out}$ distinction. \citet{boschma2015} show that the relative payoff to
related versus unrelated diversification is institutionally contingent: coordinated
market economies favour incremental related diversification, while liberal market
economies support the unrelated, exploratory kind, a result that motivates treating
the regional production system, not just the firm, as a unit of heterogeneity.
\citet{freire2019} provides a structural model of endogenous diversification in which
capability accumulation and product entry co-evolve, giving microfoundations to the
empirical relatedness regularities. Second, \citet{mealy2022} demonstrate that the
complexity toolkit transfers to any policy-relevant subset of products---their
green-complexity indices show strong path dependence in capability accumulation
(a $0.92$ correlation between green-complexity potential and the green-complexity
index, with beginning-of-period potential predicting subsequent green-export
growth), a template we adopt when restricting attention to strategically dependent
product categories. Third, \citet{nomaler2023} issue a confounding warning that
directly disciplines our outcome construction: density-based relatedness metrics are
$\sim$93 percent explained by diversity and ubiquity alone, so a measured ``relatedness''
effect can be an artefact of how many products a firm makes and how common those
products are. Their symmetric treatment of RCA \emph{gains} and \emph{losses}, together with
the finding that lost specialisations are often \emph{more} complex than gained
ones, reframes diversification as a two-sided entry-and-exit process and is one of the
closest conceptual antecedents to our reconfiguration question.

The application of complexity tools to firm-level data is recent. \citet{bruno2018, bruno2023} construct firm-product bipartite
networks from Customs microdata and show that RCA-based binarisation
(with BiCM entropy-null correction as a robustness) yields block-modular structures
that are stable across years and meaningfully predict firm growth, innovation, and
survival. \citet{pugliese2019} show that firms diversify coherently within their
existing technological portfolio, so that a firm's position within its capability
cluster is informative about its future mobility---a finding that motivates our
within-block heterogeneity analysis; the country-level analogue is \citet{pugliese2017},
who show that complex economies escape the low-income trap by moving laterally into
related capabilities. \citet{coad2021} push
this furthest, applying the FC algorithm to a firm$\times$capability matrix for $\sim$40{,}000
Indian firms and recovering a nested ``capabilities ladder'' whose rungs predict growth
and survival but---tellingly---not profitability, a decoupling we return to in
Section~\ref{subsec:lit_cap}. \citet{destefano2025}
provide the most direct methodological antecedent for our outcome measures: they
decompose firm diversification into in-block ($d_{\rm in}$) and out-of-block
($d_{\rm out}$) counts using the block partition produced by the BRIM algorithm of
\citet{barber2007}, define $\text{Share\_core} = d_{\rm in}/(d_{\rm in}+d_{\rm out})$,
and show that out-of-block diversification is the growth-positive margin while
retrenchment toward the core---the $d_{\rm in}$ direction---is the
growth-negative, lock-in-reinforcing margin. We adopt their exact outcome measures
and employ the same partition algorithm.

The network-method foundations underpinning the partition are well established.
\citet{newman2006} introduced the modularity objective and its spectral optimisation;
\citet{blondel2008} provide the Louvain heuristic that makes community detection
tractable at the scale of full-universe firm-product matrices; \citet{ducharenas2005}
supply the extremal-optimisation alternative we use as a partition robustness check;
and \citet{fortunato2016} provide the authoritative critical survey of the field's
validation pitfalls. \citet{barber2007} adapts modularity to the bipartite case (BRIM),
and \citet{soleribalta2018}
show that the BRIM objective is insensitive to within-block
nested structure and can therefore miss the hierarchical ranking of firms inside a
capability block; their in-block nestedness quality function $I$ jointly optimises the
block partition and within-block NODF \citep[the consistent nestedness metric of][]{almeidaneto2008},
providing the basis for our within-block rank ordering of firms. \citet{mariani2019} prove that the FC algorithm is a nestedness
maximiser and clarify the relationship between modularity and nestedness in firm-level
networks: the two are positively correlated at low network densities (the regime typical
of firm-product matrices) and can be jointly estimated. Validation of the BRIM
partition against entropy null models---using the Bipartite Configuration Model of
\citet{squartini2011} and \citet{saracco2015}, with the projection reconciliation of
\citet{cimini2022}---is a required pre-processing step
to ensure that block structure reflects genuine capability clustering rather than
degree-sequence artefacts, and we implement it throughout. A wider set of Economic Fitness and
Complexity (EFC) studies---broader applications of economic-complexity methods---rounds
out the methodological base: global-value-chain specialisation \citep{bontadini2021, bontadini2025},
the complexity of structural change and the sustainability transition
\citep{caldarola2024a, caldarola2024b}, machine-learning reconstruction of trade flows
\citep{gnecco2023}, and relatedness-based forecasting of industrial upgrading
\citep{albora2023a, albora2023b}.

% ---------------------------------------------------------------------------
\subsection{Capability Theory, Corporate Coherence, and the Direction of Adjustment}
\label{subsec:lit_cap}
% ---------------------------------------------------------------------------

The theoretical framework connecting firm capability stocks to the direction of
diversification draws from the resource-based view of the firm, evolutionary economics,
and the learning-and-adaptation literature. \citet{penrose1959} establishes that
firms are bundles of productive resources whose services can be redeployed across
activities at varying cost: the excess services that routinely arise from indivisible
resources create an organic pressure toward related diversification that is internal
to the firm, not driven by market structure. \citet{teece1982} formalises this in a
scope-economy model showing that the boundaries of the multi-product firm are
determined by the jointness of underlying knowledge inputs; firms expand into activities
that share tacit knowledge with existing lines, generating the coherence property
that the EFC literature quantifies through block-modular partitions.
\citet{teece1994} make this measurable through the survivor-principle relatedness
metric---sectors are ``related'' when more firms combine them than chance predicts---which
is the product-space analogue of the bipartite co-occurrence logic at the heart of the
EFC programme. The corporate-coherence literature that builds on it is directly
informative about the \emph{dynamics} we expect to observe. \citet{piscitello2000}
finds that large firms reduced product coherence while \emph{maintaining}
technological coherence over 1977--1995: capabilities are the stable organising
principle even as product portfolios churn, which is exactly why a capability-block
partition should track adjustment more faithfully than a raw product count.
\citet{valvano2003} provide the most pointed antecedent for a trade-shock prediction:
in Italian manufacturing, firms in sectors more exposed to single-market competition
(``sensitive'' sectors) became \emph{more} coherent over time, and entrants into the
leadership ranks were more coherent than incumbents---direct evidence that intensified
competition pushes surviving firms toward their core ($d_{\rm in}$). \citet{delios2008}
add the institutional contingency: within-country product diversification is more
valuable where market institutions are weak, because diversification substitutes for
missing capital, labour, and information markets---a moderation channel that warns
against reading every diversification response as pure capability exploitation.

The corporate-finance literature on the ``diversification discount'' supplies the
causal-direction caution that disciplines our interpretation. \citet{langstulz1994}
document a robust industry-adjusted discount in Tobin's $q$ of 0.35--0.50 for
diversified US firms throughout 1978--1990, but show that firms diversify \emph{after}
their core industry weakens---the discount precedes diversification rather than being
caused by it. \citet{matsusaka2001} rationalises this as capability-match search: a
firm whose current match deteriorates maintains its existing lines while experimenting
with new ones (the ``$b$-strategy''), liquidating and restarting only when the match
falls below a lower threshold. \citet{bernardo2002} recast the same fact as
real-option value: single-segment firms are worth more because they retain unexercised
expansion options. Translated to our setting, these models deliver a sharp,
non-obvious prediction: a \emph{moderate} tariff shock should push firms toward
exploratory $d_{\rm out}$ search (the $b$-strategy), whereas a \emph{severe} shock that
drives match quality below the liquidation threshold collapses the option value of
organisational capabilities and forces retrenchment or exit---so the sign of the
average response is genuinely ambiguous and should depend on shock intensity, exactly
the non-linearity our event study is designed to detect.

The evolutionary perspective of \citet{dosi1988, dosi1995} adds a regime dimension:
the rate and direction of capability accumulation depend on both the firm's own
endowment and the technological regime of the sector---the opportunity conditions,
appropriability, and cumulativeness that structure search. Under Schumpeterian Mark
II conditions (cumulative, large-firm-dominated regimes), incumbents sustain
technological advantages through internal learning, generating the technological
lock-in that maps directly onto high $\text{Share\_core}$. The implication for our
setting is that firms in complex, cumulative technological regimes should display
stronger $d_{\rm in}$ responses to adverse shocks---retrenchment to the core is both
the natural adjustment and the organisationally feasible one when search capabilities
are already embedded in existing product clusters.

\citet{march1991} provides the ambiguity at the heart of our hypotheses. The
exploitation--exploration tension implies that the direction of capability adjustment
under stress is not a priori signed: firms that face an adverse demand or competitive
shock in their existing core can either intensify exploitation of known competencies
(retrenchment) or accelerate exploration beyond the core (diversification-as-escape).
March's formal result is that organisations systematically over-exploit under most
plausible parameter configurations, predicting the retrenchment (H1) hypothesis as
the modal response but leaving open the exploration (H1$'$) case for firms with
sufficient slack and absorptive capacity. \citet{patel1997} document the empirical
content of this tension at the firm level, showing that large innovative firms spread
R\&D effort across fields but in highly persistent patterns---path dependence is the
rule, not the exception---while \citet{bell1993} establish that technological
capability accumulation in developing-country firms is even more path-dependent,
with explicit investments in learning required to escape existing trajectories.
\citet{storper1992} connects this to trade: product-based technological learning
generates agglomeration economies in particular product-territory combinations, so
that an adverse trade shock removes both demand and the collective learning
infrastructure that sustains capability accumulation---a channel that implies
persistent, not merely transitory, capability effects of tariff shocks.

The firm-growth literature, summarised in \citet{coad2009} and documented empirically
in \citet{bottazzi2003, bottazzi2005, bottazzi2006}, shows that firm-level growth is
Laplacian---fat-tailed, near-unit-root---and that the variance of growth rates scales
inversely with firm size at a rate explained entirely by diversification across product
lines. The scope-economy mechanism \citep{bottazzi2006} implies that firms which
diversify within a related cluster (high $d_{\rm in}$) reduce the co-movement of
their individual product revenue streams more effectively than firms that scatter across
unrelated products. Under tariff shock conditions, firms that lose access to existing
export markets face a portfolio rebalancing problem whose solution---from a pure
variance-minimisation standpoint---is to deepen within-block exposure rather than
scatter to unfamiliar destinations, reinforcing the retrenchment prediction. A
recurring caveat from this literature is that performance dimensions decouple:
\citet{grazzi2012} shows that Italian exporters carry a large productivity premium but
\emph{no} profitability premium, and \citet{coad2021} that capability-ladder position
predicts growth and survival but not profit margins; \citet{syverson2011} surveys the
enormous, persistent productivity dispersion that these patterns sit inside. The
collective lesson---that capability advantage need not show up in profits or even in
growth---is why we measure the \emph{architecture} of the export basket
($d_{\rm in}$, $d_{\rm out}$, $\text{Share\_core}$) rather than a performance scalar.
\citet{aghion2001, aghion2005} add the innovation margin: neck-and-neck competitors
facing increased competitive pressure innovate to escape competition, while laggards
reduce innovation effort. Translating to our setting, firms with high within-block
fitness ranks face incentives to deepen core capabilities under tariff pressure (the
escape-competition channel), while low-rank firms---unable to sustain the cost of
maintaining their current product portfolio---are more likely to prune $d_{\rm out}$
and consolidate around the most defensible core products. This generates the
within-block heterogeneity in treatment effects that we exploit empirically.

% ---------------------------------------------------------------------------
\subsection{Trade Policy Shocks and Firm Adjustment}
\label{subsec:lit_trade}
% ---------------------------------------------------------------------------

A large reduced-form literature establishes that tariff changes cause significant
adjustments in firm outcomes. \citet{pavcnik2002} and \citet{topalova2011} document
that trade liberalisation raises TFP through competitive selection and technology
imports in Chilean and Indian manufacturing, respectively. \citet{amiti2007} show,
for Indonesian firms, that input tariff reductions deliver TFP gains at least twice as
large as those from output tariff reductions---a finding driven by access to high-quality
imported intermediates---and construct the shift-share input-tariff exposure variable that is
the methodological template for our treatment. \citet{tomasi2022} replicate this
design for Vietnamese firms during the 2007 WTO accession, finding that governance
quality moderates the tariff-TFP effect, providing the closest precedent for our
moderated shift-share design with ExpDep as the heterogeneity variable.

What this literature consistently measures is performance: productivity, employment,
markups, survival. It is essentially silent on the \emph{composition} of a firm's
product portfolio---which activities are retained, dropped, or newly entered, and
whether the reconfiguration deepens a coherent core or spreads capabilities into
new domains. The EFC literature reviewed in Section~\ref{subsec:lit_efc} supplies
the tools to measure that composition, but has applied them almost exclusively in
descriptive, non-causal settings. This paper occupies the intersection left empty
by both strands. The theoretical channel linking trade exposure to capability
adjustment is supplied by the trade-and-innovation literature: \citet{akcigit2022}
unify the market-size, business-stealing, and escape-competition effects of trade on
innovation incentives, showing that the net effect on a firm's capability investment
is sign-ambiguous and depends on its competitive position---the firm-level analogue of
the exploitation--exploration tension. \citet{riverabatiz1991} establish the deeper
point that integration affects growth through flows of \emph{ideas} as well as goods,
so that a tariff shock that severs a firm from a technological frontier market can have
permanent (growth-rate), not merely level, effects on its capability trajectory---reinforcing the
persistence prediction from \citet{storper1992}.

The most direct evidence that trade shocks restructure capability portfolios rather
than merely rescaling performance comes from \citet{fontagne2018}, who show that
Chinese export quotas on textiles and clothing generated a significant \emph{Typical
Product Vector} consolidation among French exporting firms: products at the core of
a firm's stable portfolio were retained at a 7.7 percentage-point higher rate than
peripheral products in treated sectors. This is, in effect, a $d_{\rm in}$ response
to a trade-policy shock measured with product-vector tools---the closest antecedent
in the literature for our main outcome. \citet{fontagne2022} provide the
product-level trade elasticities that we use as Bartik weights: 5,050 HS6-digit
elasticities estimated via PPML gravity, averaging $-5.3$ (median $-3.5$), with
HS-section means spanning $-18.5$ (minerals), $-6.1$ (machinery), and $-3.6$
(footwear), capturing the heterogeneous disruption that a uniform tariff schedule
imposes on firms with different product portfolios.

The dynamic and firm-level dimensions of portfolio adjustment are examined by
\citet{crozet2021}, who study the effect of EU and US sanctions on the extensive
margin of French exporters to Iran and Russia. Using monthly customs data and a
dynamic probit model with three-way fixed effects and analytical bias correction,
they find that sanctions lower the probability of serving a sanctioned market by 39
percent (Iran) and 23 percent (Russia), with strong state dependence: the entry-cost
factor is 19--23 percent higher for firms not active in the previous year. Crucially,
the lifting of sanctions on Myanmar generates only a small and gradual positive response,
establishing an \emph{asymmetry} between imposition and removal that implies permanent,
not transient, $d_{\rm out}$ contraction even after the policy shock is reversed.
This hysteresis result---consistent with the Roberts--Tybout sunk-cost model---means
that treatment effects in our event study should be interpreted as reflecting a
permanent restructuring of the capability portfolio, not a temporary demand disruption.
The firm-level heterogeneity documented by \citet{crozet2021}---that trade-finance-dependent
firms, inexperienced firms, and non-specialist firms exit more sharply---maps
directly onto the moderation hypotheses we test: private Chinese firms (documented
by \citet{guariglia2011} as the most financially constrained in the NBS panel) should
display amplified $d_{\rm out}$ reductions, while SOEs with soft budget constraints
should display more muted responses.

A complementary strand examines the export portfolio diversification strategies that
firms deploy to manage demand volatility. \citet{vannoorenberghe2016} show, for
Chinese firm-level data, that destination diversification reduces volatility for large
exporters but \emph{increases} it for small ones, with the switching point determined
by a composition effect: occasional exports to marginal destinations dominate the
diversification gain for small firms, while stable core-market exports dominate for
large ones. \citet{crozet2025} formalise this puzzle through a phyloeconomic diversity
index $\Phi$ that adds a dissimilarity dimension to the standard richness-evenness
decomposition. Their central finding---that conditional on the number of destinations
and the Herfindahl concentration of export sales, higher $\Phi$ (more dissimilar
markets) is associated with \emph{higher} export sales volatility---reveals that
hedging through geographic diversification has limits: distant, structurally different
destination markets carry higher idiosyncratic demand variance, so that firms forced
to substitute familiar Western markets with unfamiliar alternatives do not achieve
the volatility reduction that naive diversification logic predicts. For our setting,
this implies that Chinese firms which maintain $d_{\rm out}$ by pivoting to non-Western
destinations face not only lower export levels but higher volatility---a double cost
that reinforces the variance-minimising case for $d_{\rm in}$ consolidation.
\citet{barbieri2024} provide a parallel decomposition at the territorial level:
\emph{sectoral} diversification (a proxy for productive capabilities, the conceptual
sibling of $d_{\rm in}$) raises export intensity, while \emph{geographic}
diversification (a proxy for geoeconomic capabilities, loosely analogous to the
resilience role we attribute to $d_{\rm out}$) is associated with greater export
stability; institutional quality independently improves both margins.\footnote{The
analogy is functional rather than literal. Barbieri's second axis is
\emph{destinations}, whereas our $d_{\rm out}$ remains a \emph{product}-space
measure (out-of-block products), following \citet{destefano2025}, who works in pure
firm--product space with no destination dimension. The roles align---one margin scales
output, the other buys resilience---but the underlying axes differ, and in their data
the dominant stabiliser is EU market exposure rather than geographic diversification
per se. We therefore treat the destination dimension as a second-stage extension
(Section~\ref{subsec:spinoffs}) rather than as a component of $d_{\rm out}$.}

We estimate treatment using the Global Tariff Database (GTD) of \citet{teti2024},
the first applied-tariff panel to cover bilateral HS6 tariffs continuously from 1988
to 2021 while correcting for the false interpolation and selection artefacts of UNCTAD
TRAINS. \citet{teti2019} establish, using structural gravity, that Rules-of-Origin
compliance is prohibitively costly for the machinery and electronics sectors at the
core of Chinese exports, confirming that MFN tariffs---the rate captured by the
GTD---are the economically relevant wedge for our firms. The welfare magnitudes from
\citet{teti2025} calibrate our expected effect sizes: a GE simulation of the
2018--2025 US--China tariff war estimates cumulative real-income losses of roughly
$-1.6$ percent for Canada and Mexico---economies with high export exposure to the US
market---against only $-0.3$ percent for China in the aggregate. We use the
$-1.6$ percent figure as an external benchmark for the most exposed Chinese firms:
by analogy with the highly exposed Canadian and Mexican case, the uneven distribution
of tariff exposure across product portfolios should push firms concentrated in
high-regulatory-dependence categories toward magnitudes well above the national
average. The mapping from these country-level losses to firm-level exposure is our
own; \citet{teti2025} simulate countries and US states, not firms. The distortionary
role of size-contingent and sector-specific regulation documented by
\citet{garicano2016}---who find welfare costs of 1.3--3.5 percent of GDP from French
size thresholds---is the domestic-policy analogue of this exposure heterogeneity and
motivates controlling for regulatory environment in the moderation analysis.

% ---------------------------------------------------------------------------
\subsection{Strategic Dependencies and Geopolitical Trade Shocks}
\label{subsec:lit_geo}
% ---------------------------------------------------------------------------

The design of the heterogeneity variable ExpDep draws from a separate and newer
literature on strategic economic dependencies. \citet{farrell2019} identify the
conditions under which the topology of global economic networks generates geopolitical
leverage: hub-and-spoke structures create chokepoint nodes through which flows can be
controlled (flow restriction) and from which information can be extracted (panopticon
effects). Crucially, weaponisation requires two joint conditions---jurisdictional
control over the chokepoint and domestic institutions capable of deploying that control
as a foreign policy instrument---which limits the set of products and bilateral
relationships that qualify as strategic in the relevant sense. The Farrell--Newman
framework grounds our moderation hypothesis: Chinese firms whose export baskets are
concentrated in products where Western jurisdictions have identified chokepoint control
face not a neutral tariff shock but one accompanied by active policy interest in
reducing dependence---generating an amplified and persistent treatment effect.

The quantitative operationalisation of this insight is provided by \citet{euc2021},
who develop the Core Dependency Indicators (CDI) methodology to identify EU strategic
dependencies at the HS6-product level. The methodology applies three thresholds:
supplier concentration (HHI of extra-EU import shares exceeding 0.4, implying fewer
than 2.5 effective suppliers), import reliance (extra-EU imports exceeding 50 percent
of total EU imports of the product), and substitutability (extra-EU import volume
exceeding total EU export volume, so that domestic production cannot cover the import
volume even in an emergency). Of approximately 5,000 HS6 products, 137 pass all three
thresholds in the EU's most sensitive ecosystems, with China accounting for roughly
52 percent of their import value. \citet{arjona2023} refine the methodology with a
dynamic four-year threshold and a re-export correction: accounting for transit trade
\emph{raises} the count of dependent products by roughly 27 percent, since ignoring
re-exports understates true extra-EU reliance, yielding a dependent-product set of
204 items that provides the product classification we use to construct ExpDep.
The forward-looking scenario work of \citet{cagnin2021} situates these indicators in
the EU's Open Strategic Autonomy agenda, projecting an escalation of dependency-reducing
policy through 2040 that makes the persistence of any treatment effect we estimate a
policy-relevant, not merely statistical, question.

The network-science analogue of the CDI methodology is provided by
\citet{korniyenko2017}, who combine three structural indicators of product trade networks---
the standard deviation of weighted outdegree centrality (capturing hub dominance), the
product of clustering coefficient and network diameter (capturing geographic
compartmentalisation), and a Revealed Factor Intensity measure of substitutability---to
classify 3,578 intermediate and capital goods into four risk tiers via k-median
clustering. Their persistently risky group comprises 421 products over 2003--2014,
concentrated in machinery and mechanical appliances (HS 84--85), transport equipment,
pharmaceuticals, and precision instruments. The causal validity of the classification
is demonstrated by two out-of-sample cases: the 2011 Japan earthquake and Thai floods,
where products subsequently in short supply were correctly identified as Group~4 one
year prior. The overlap between the Korniyenko Group~4 products and the
Arjona CDI-dependent set defines the most extreme ExpDep exposure: products
that are simultaneously structurally concentrated, topologically fragile, and actively
targeted by EU and US reshoring policy.

\citet{evenett2024} provide real-time evidence on the policy escalation dynamic. Their
New Industrial Policy Observatory (NIPO), covering 2,500$+$ new industrial policy
measures across 75 jurisdictions in 2023, documents a 73.8 percent probability that a
subsidy by one major economy (China, EU, or US) is matched by a subsidy in the same
product by another within 12 months. PPML regressions confirm that IP usage correlates
positively with RCA---governments target sectors where they already have comparative
advantage, consistent with political-economy capture rather than market-failure
correction. For our event study, this finding implies that the Teti GTD tariff
shock at time $t$ is not a clean one-off treatment but the first signal in an
escalating policy sequence: treatment intensity grows after the event-window, pushing
treatment effects toward permanence.

\citet{fontana2024} embed these observations in a political-economy theory of
strategic-assets policy, drawing on the three logics of capability blockade formalised
by Ding and Dafoe: cumulative-strategic blocking (targeting technological RCA through
standards exclusion and export controls), infrastructure-strategic blocking (targeting
critical material supply chains through reshoring subsidies), and dependency-strategic
blocking (exploiting existing dependencies through sanctions and FDI screening). The
taxonomy provides a structural basis for splitting the Teti tariff sample by motive:
tariff shocks in sectors matching the cumulative-strategic logic (advanced
manufacturing, semiconductors) should generate stronger $d_{\rm in}$ responses than
shocks in sectors matching the infrastructure-strategic logic (raw materials, chemicals),
where substitution is constrained by the CDI3 condition. We use this taxonomy in a
sample-split robustness exercise.

% ---------------------------------------------------------------------------
\subsection{Identification}
\label{subsec:lit_id}
% ---------------------------------------------------------------------------

Our identification strategy rests on a single organising principle: \emph{everything
that could mechanically respond to the shock must be predetermined relative to it.}
The estimand is the within-firm response of an export-capability portfolio to a
trade-policy shock, and the credibility of that estimand depends on three objects
being fixed before the shock---the exposure weights, the disruption elasticities, and
the capability map against which the portfolio is read. We discuss each, and then
state precisely which of three admissible network constructions we adopt and why.

\paragraph{A continuous, time-varying treatment.}
Following the input-tariff shift-share of \citet{amiti2007}, \citet{topalova2011}, and
\citet{tomasi2022}, we construct the tariff exposure variable as
\[
  \text{TariffShock}_{it}
  = \sum_p s_{ip,\text{pre}}\cdot \varepsilon_p \cdot \Delta\ln(1+\tau_{pt}),
\]
where $s_{ip,\text{pre}}$ are firm--product export shares in a strictly pre-shock
baseline window, $\Delta\ln(1+\tau_{pt})$ are changes in bilateral applied tariffs
from the Teti GTD, and $\varepsilon_p$ are the product-level trade elasticities of
\citet{fontagne2022}. We treat $\text{TariffShock}_{it}$ as a \emph{continuous,
time-varying, non-absorbing} dose rather than a binary staggered treatment: the dose
grows and recedes as tariffs move, and discretising it into a treated/control
indicator would discard precisely the cross-firm variation in exposure intensity that
identifies the dose--response. This choice is deliberate and consequential, as we make
clear below. The elasticities $\varepsilon_p$ are estimated by \citet{fontagne2022} on
the universe of 152 importers and 189 exporters---they are a property of the
\emph{product}, external to our firm sample, and hence cannot be a function of the
behaviour we study; products with a non-significant estimate inherit their HS-section
mean. The predetermined shares satisfy the Bartik exogeneity condition, which we defend
on \emph{both} of its routes: from the shares, following \citet{goldsmithpinkham2020},
and from the shocks, following \citet{borusyakhull2022}, reporting Rotemberg-weight
diagnostics and a recentred-instrument specification.

\paragraph{Dynamics and inference for a continuous dose.}
Because the treatment is continuous and non-absorbing, the cohort$\times$relative-time
machinery of \citet{sunabraham2021} is not the natural primary specification---its
cohorts are defined by the timing of a binary switch. We therefore estimate the
dynamics through a distributed-lag event study, $Y_{it}=\sum_k \beta_k\,
\text{TariffShock}_{i,t-k}+\alpha_i+\delta_t+\varepsilon_{it}$, in which the lead
coefficients ($k<0$) constitute the parallel-trends test and the lag coefficients
($k>0$) trace the persistence (hysteresis) prediction. Inference follows the
shift-share-robust standard errors of \citet{adao2019}, which are native to continuous
exposure designs and which we report as primary against firm- and sector-clustered
alternatives. As robustness we re-estimate the dynamics with the continuous,
non-absorbing estimator of \citet{dechaisemartin2020}, the continuous-treatment
difference-in-differences framework of \citet{callawaygoodmanbacon2024}, and---after
discretising exposure into terciles---the heterogeneity-robust estimators of
\citet{sunabraham2021}, \citet{callawaysantanna2021}, and \citet{borusyak2024}. We
flag the assumption that the continuous design carries a \emph{stronger} parallel-trends
requirement than the binary case---\citet{callawaygoodmanbacon2024} show that
identifying the dose--response slope requires the absence of selection into dose level,
which we probe but cannot fully test.

\paragraph{The capability map: three admissible constructions, and the one we adopt.}
The outcomes $d_{\rm in}$, $d_{\rm out}$, and $\text{Share\_core}$ are read off a
block partition of the firm--product export network \citep[the export, not production,
basis of][]{destefano2025}, obtained by bipartite modularity maximisation
\citep{barber2007} on an RCA-binarised incidence matrix validated against the Bipartite
Configuration Model \citep{squartini2011,saracco2015,cimini2022}. The non-trivial
question is \emph{when} the partition is estimated, because the map is itself the
measuring instrument for the outcome, and a map estimated on post-shock data would
co-move with the treatment. Three constructions are admissible:
\begin{enumerate}
  \renewcommand{\labelenumi}{(\roman{enumi})}
  \item \textbf{Frozen baseline partition.} The partition $b(p)$ and each firm's core
        block $c(f)$ are estimated once, on the pre-shock baseline window, and held
        fixed; only the firm's RCA incidence $M_{fpt}$ varies over time. The
        panel variation in $d_{\rm in}/d_{\rm out}$ is then unambiguously the firm
        relocating against a frozen, predetermined grid.
  \item \textbf{Predetermined evolving partition.} The map is allowed to evolve, but
        estimated only on \emph{pre-shock} slices through the temporal multilayer
        modularity of \citet{mucha2010}---with inter-slice coupling and resolution
        chosen by partition stability \citep{bazzi2016}---and then extrapolated forward.
        This nests (i) as the strong-coupling limit and captures any pre-existing drift
        in the capability space, at the cost of measurement error that grows with the
        post-shock horizon.
  \item \textbf{Contemporaneous multilayer.} The map evolves using all years including
        post-shock. Under our continuous treatment this construction is \emph{not}
        causally identified: with no never-treated firms, there is no
        treatment-independent control group from which to anchor the map, so the
        partition---and hence the $(d_{\rm in}+d_{\rm out})$ denominator---remains a
        function of the dose. We therefore use it only descriptively, to document how
        the capability space reorganises, and never as the causal outcome.
\end{enumerate}
\textbf{We adopt (i) as the headline specification}, precisely because it has the
smallest researcher degrees of freedom and the cleanest exogeneity argument: the map
literally precedes the shock and cannot respond to it. Its one weakness---that a frozen
map grows stale---generates classical measurement error and therefore attenuation
\emph{against} our hypotheses, an honest bias rather than a confounding one. We
neutralise even that concern by reporting (ii) as robustness, which lets the map evolve
on pre-shock information and leaves the estimates intact, and by re-estimating (i) with
a leave-one-out partition that excludes each focal firm from the map that classifies
it, breaking the reflexive channel. Two construction rules are stated explicitly and
checked: products with no baseline block are assigned by product-space proximity (or
dropped, with the affected export share reported), and firms born after the baseline
are handled within the exit-selection analysis rather than assigned a post-baseline
core. We compute RCA with the standard contemporaneous, \emph{within-China} Balassa
denominator: because the capability blocks are themselves built from the Chinese
firm--product network, a China-relative benchmark---specialisation relative to the
Chinese export economy---is the internally consistent and economically natural object,
and we retain it as the headline measure. The denominator carries a second-order
\emph{inference} concern rather than an economic one: under a sector-wide shock the
aggregate Chinese export of a treated product falls for all exposed firms, mechanically
inflating the within-sample benchmark and potentially moving threshold crossings for
reasons unrelated to the firm's own portfolio choice. We treat this as robustness, not
as a change of measure. We re-estimate with a leave-one-out within-China denominator,
which preserves the China-relative interpretation while removing the firm's reflexive
contribution, and, as a more aggressive check, against a rest-of-world (ex-China) BACI
benchmark, which trades the China referent for full exogeneity to the Chinese firms'
treated response.

\paragraph{An elegant alternative under a modified treatment.}
We note one design under which the \emph{contemporaneous} evolving map (iii) would
become causally admissible: were we to trade the continuous dose for a binary,
staggered treatment, the not-yet-treated firms would furnish a treatment-independent
control group from which the time-evolving partition could be estimated at each date
\citep[in the spirit of the control-group construction of][]{callawaysantanna2021},
yielding a map that is simultaneously evolving and exogenous. This is a genuine
optimisation of elegance over the frozen design, but it is purchased by discarding the
exposure-intensity information that motivates our continuous specification, and it
relies on a credible never-/not-yet-treated set under the diffuse Teti tariff
schedule. We therefore retain the continuous treatment with the frozen map as our
primary design, and present the binary control-group-map construction as a deliberate
robustness fork rather than the main specification.

\paragraph{Two concessions we make at the outset.}
Neither threat is specific to the network construction, and we address both head-on
rather than conceal them. First, the shock induces firm exit \citep{crozet2021}; since
an exited firm has no $d_{\rm in}/d_{\rm out}$, survivors are selected, and we therefore
estimate the effect on the probability of exit, bound the survivor effect, and read our
portfolio estimates as conservative. Second, a tariff on one product reallocates demand
to substitutes held by other firms in the same block, so the stable-unit-treatment
assumption is strained; the \citet{adao2019} correction absorbs the resulting
cross-firm correlation in the standard errors but not the causal interference, which we
discuss and, where feasible, probe with leave-out-block exposure.

\paragraph{Firm-level data and the ownership moderator.}
The panel of Chinese exporting firms comes from the NBS Annual Survey of Industrial
Enterprises, used in the same form by \citet{guariglia2011}. That paper establishes
three features of the NBS data that shape our design: SOEs operate under soft budget
constraints and face negligible financial frictions, private domestic firms are
severely financially constrained (near-unit-root self-financing), and coastal firms
without foreign capital face the binding constraint. Because firms switch ownership
status over the panel, we assign type by a consistent across-years rule following
\citet{guariglia2011}. The ownership heterogeneity motivates our SOE--private split:
privately owned firms facing financial constraints respond more sharply to export market
closures because they cannot subsidise loss-making product lines from retained earnings,
making the $d_{\rm out}$ contraction more immediate and the $d_{\rm in}$ consolidation
more forced.

% ---------------------------------------------------------------------------
\subsection{Our Contribution}
\label{subsec:contribution}
% ---------------------------------------------------------------------------

This paper contributes to all five literatures described above. To the economic-complexity
literature, we provide the first causal identification of a firm-level capability
response---$d_{\rm in}$, $d_{\rm out}$, and $\text{Share\_core}$---to an exogenous external shock,
demonstrating that the block-modular firm--product partition is not merely a descriptive
device but a meaningful predictor of differential adjustment paths. To the capability-theory
and evolutionary-economics literatures, we offer empirical content for the
exploitation--exploration tension formalised by \citet{march1991}: the direction of
adjustment under tariff shocks is shown to be systematically moderated by regional
strategic dependence, consistent with the Penrose--Teece prediction that capability
redeployment is constrained by the jointness of underlying knowledge inputs, while the
magnitude is consistent with the permanent-effect predictions of \citet{storper1992}
and the lock-in mechanism documented in the Schumpeterian Mark~II literature.

To the trade literature, we shift the outcome from performance (the productivity,
markup, and employment measures that dominate the Pavcnik--Topalova--Amiti generation) to
capability architecture, providing evidence that the same tariff shocks that raise
aggregate TFP through competitive selection simultaneously compress the capability
frontier of exposed firms in ways that the performance lens does not reveal. To
the strategic-dependency literature, we take the Farrell--Newman chokepoint
hypothesis to firm-level data: firms operating in CDI-dependent product
categories display systematically larger capability consolidation effects than firms
in non-dependent categories, and this heterogeneity is not explained by tariff
magnitudes alone, consistent with the dual treatment of higher tariffs and active
EU reshoring policy targeting the same product space. Finally, to the identification
literature, we demonstrate that a continuous-dose distributed-lag event study---with the
heterogeneity-robust continuous-treatment estimators of \citet{dechaisemartin2020} and
\citet{callawaygoodmanbacon2024} as robustness---is tractable at the scale of
full-universe firm--product panels, that the \citet{adao2019} shift-share-robust
correction materially affects standard errors in heterogeneous-exposure settings with
product-level clustering, and that reading a network-derived outcome against a
\emph{predetermined} (frozen, or pre-shock-estimated) capability partition is what keeps
the estimand exogenous to the shock.

% =============================================================================
\section{Expected Dynamics: What the Literature Predicts We Should Observe}
\label{sec:expected}
% =============================================================================

Because the empirical estimates are not yet in hand, it is worth making the priors
explicit. This section reads the literatures of Section~\ref{sec:lit} \emph{forward}:
for each outcome and each moderator, it states what dynamic the prior work leads us to
expect once the planned design is executed, what sign each hypothesis carries, and
what magnitude range is plausible. The aim is to specify in advance what a result
``consistent with the literature'' would look like---so that a surprising estimate can
be recognised as surprising---and to organise the diverse predictions into a single
testable structure. Nothing here is a finding; everything here is a prediction that
the data will confirm, refute, or qualify.

% ---------------------------------------------------------------------------
\subsection{The Out-of-Block Margin ($d_{\rm out}$)}
\label{subsec:exp_dout}
% ---------------------------------------------------------------------------

Three literatures converge on a contraction of $d_{\rm out}$ following an adverse
tariff shock. The sunk-cost extensive-margin mechanism of \citet{crozet2021} predicts
that exposed firms stop serving the products and markets that become tariff-burdened,
and---because re-entry costs are 19--23 percent higher for lapsed exporters---the
contraction is \emph{hysteretic}: it should not reverse within the event window even
if tariffs are later rolled back. The variance-cost mechanism of \citet{crozet2025}
and \citet{vannoorenberghe2016} predicts that firms which try to \emph{preserve}
$d_{\rm out}$ by pivoting to unfamiliar, structurally dissimilar destinations inherit
higher demand volatility, so the privately optimal response for all but the largest
firms is to let $d_{\rm out}$ fall rather than chase replacement markets. And the
diversification-discount search models of \citet{matsusaka2001} and \citet{bernardo2002}
predict that only firms whose core match remains above the liquidation threshold retain
the option to explore; below it, exploratory $d_{\rm out}$ search collapses. The net
expectation is a \emph{negative} average effect on $d_{\rm out}$, concentrated among
smaller, financially constrained, and inexperienced firms, and persistent over the
post-shock horizon. The one literature that cuts the other way is the escape-competition
margin of \citet{aghion2005} and \citet{akcigit2022}: a high-capability firm facing
competitive pressure may innovate \emph{outward}, lifting $d_{\rm out}$. Because that
channel is confined to frontier firms, the expectation is a negative average masking a
positive tail---visible only in the interaction with firm capability rank, not in the
mean.

% ---------------------------------------------------------------------------
\subsection{The In-Block Margin ($d_{\rm in}$) and Retrenchment}
\label{subsec:exp_din}
% ---------------------------------------------------------------------------

The prediction for $d_{\rm in}$ is the cleaner of the two. The Typical-Product-Vector
consolidation of \citet{fontagne2018}---a 7.7 percentage-point retention premium for
core products in treated sectors---is, mechanically, a $d_{\rm in}$ response to a
trade-policy shock, and it is the single closest empirical antecedent for our main
outcome. The scope-economy variance argument of \citet{bottazzi2006} predicts that a
firm rebalancing a disrupted portfolio minimises revenue covariance by deepening within
a related cluster rather than scattering, which raises $d_{\rm in}$ relative to
$d_{\rm out}$. The competition-to-coherence evidence of \citet{valvano2003}---firms in
single-market-exposed sectors became more coherent over time---predicts the same sign.
And the Schumpeterian Mark~II lock-in of \citet{dosi1995} predicts that the retrenchment
is strongest exactly where capabilities are most cumulative. The expectation is
therefore a \emph{positive} average effect on $d_{\rm in}$, larger for firms in
complex, cumulative technological regimes, and---per \citet{storper1992} and
\citet{riverabatiz1991}---persistent rather than transitory, because what is lost when
a market closes is more than demand: it is the collective learning infrastructure that
sustained the broader basket.

% ---------------------------------------------------------------------------
\subsection{Share\_core and the Direction of Adjustment}
\label{subsec:exp_share}
% ---------------------------------------------------------------------------

The summary statistic $\text{Share\_core} = d_{\rm in}/(d_{\rm in}+d_{\rm out})$ signs
the exploitation--exploration choice of \citet{march1991}. If the $d_{\rm out}$
contraction and the $d_{\rm in}$ consolidation predicted above both materialise,
$\text{Share\_core}$ rises---the retrenchment hypothesis (H1). March's own result that
organisations systematically over-exploit makes this the modal expectation. But the
literature is deliberately not unanimous: the dynamic-capabilities reading
(\citealp{penrose1959}; \citealp{teece1982}) and the moderate-shock search prediction
of \citet{matsusaka2001} leave open the exploration case (H1$'$), in which $\text{Share\_core}$
\emph{falls} because firms with slack treat the shock as a prompt to reallocate toward
new capability bundles. The empirical question is therefore which force dominates on
average, and, more informatively, where in the firm distribution the sign flips. The
strong prior, following \citet{destefano2025}, is that the average response is toward
the growth-\emph{negative} core (H1), because the growth-positive $d_{\rm out}$ margin
is precisely the one that sunk costs and variance costs make hardest to defend under
stress.

% ---------------------------------------------------------------------------
\subsection{Heterogeneity and the Strategic-Dependence Moderator}
\label{subsec:exp_het}
% ---------------------------------------------------------------------------

The sharpest predictions are conditional, not average. The Farrell--Newman chokepoint
logic \citep{farrell2019}, operationalised through the CDI product set of
\citet{arjona2023} and the network-fragility tiers of \citet{korniyenko2017}, predicts
that firms whose baskets concentrate in strategically dependent products face not a
neutral tariff but a tariff \emph{plus} active reshoring policy aimed at the same
products---so their retrenchment should be amplified and, given the escalation dynamic
of \citet{evenett2024}, more persistent. The financial-constraint heterogeneity of
\citet{guariglia2011} predicts that private Chinese firms retrench more sharply than
soft-budget SOEs. The size-and-volatility composition effect of
\citet{vannoorenberghe2016} predicts that small firms cannot afford to defend
$d_{\rm out}$ at all. The institutional-contingency result of \citet{delios2008} and
\citet{boschma2015} predicts that the related-versus-unrelated payoff itself shifts
with the regional institutional environment, so that ExpDep should be read as a
production-system attribute and not only a firm attribute. The blockade-logic taxonomy
of \citet{fontana2024} predicts that cumulative-strategic sectors retrench harder than
infrastructure-strategic ones. The unifying expectation is that \emph{heterogeneity
dominates the average}: a modest mean effect concealing a steep gradient in strategic
dependence, financial constraint, and capability rank.

% ---------------------------------------------------------------------------
\subsection{Synthesis: Hypotheses, Signs, and Magnitude Benchmarks}
\label{subsec:exp_synth}
% ---------------------------------------------------------------------------

Table~\ref{tab:hyp} collects the hypotheses, their literature sources, and the sign
each predicts. Table~\ref{tab:mag} brackets the plausible magnitude range using two
external benchmarks: a lower bound from pure network-disruption transmission
(\citealp{korniyenko2017}: a 1pp rise in risky imports from a disaster-struck origin
cuts the importer's own exports by $\sim$0.7 percent, with no policy response) and an upper
bound from the most exposed economies in a targeted tariff war (\citealp{teti2025}:
$-1.6$ percent cumulative real income for Canada and Mexico under the 2018--2025 US tariffs,
versus $-0.3$ percent for China in the aggregate---a country-level benchmark we map onto the
most exposed Chinese firms).

\begin{table}[ht]
\centering
\caption{Hypotheses, literature sources, and predicted signs.}
\label{tab:hyp}
\footnotesize
\setlength{\tabcolsep}{4pt}
\renewcommand{\arraystretch}{1.3}
\begin{tabular}{|c|p{4.4cm}|p{4.2cm}|p{2.5cm}|}
\hline
\textbf{Hyp.} & \textbf{Statement} & \textbf{Principal literature source} & \textbf{Predicted effect} \\
\hline
H1 & Adverse tariff shock raises $\text{Share\_core}$ / $d_{\rm in}$ (retrenchment to the coherent core). &
\citet{fontagne2018}; \citet{bottazzi2006}; \citet{valvano2003}; \citet{dosi1995}; \citet{march1991} &
$d_{\rm in}\uparrow$, $\text{Share\_core}\uparrow$ \\
\hline
H1$'$ & Adverse shock raises $d_{\rm out}$ / lowers $\text{Share\_core}$ (exploration beyond the core). &
\citet{penrose1959}; \citet{teece1982}; \citet{matsusaka2001}; \citet{aghion2005} &
$d_{\rm out}\uparrow$ (frontier-firm tail only) \\
\hline
H2 & The retrenchment response is amplified in more strategically dependent production systems. &
\citet{farrell2019}; \citet{arjona2023}; \citet{korniyenko2017}; \citet{evenett2024} &
$\beta_{2}>0$ on Tariff $\times$ ExpDep \\
\hline
H3 & The $d_{\rm out}$ contraction is hysteretic (does not reverse within the window). &
\citet{crozet2021}; \citet{storper1992}; \citet{riverabatiz1991} &
persistent, asymmetric \\
\hline
H4 & Responses are steeper for financially constrained / private / small firms. &
\citet{guariglia2011}; \citet{vannoorenberghe2016}; \citet{coad2021} &
steeper gradient \\
\hline
\end{tabular}
\end{table}

\begin{table}[ht]
\centering
\caption{Magnitude benchmarks for the $d_{\rm out}$ response (per s.d.\ of TariffShock).}
\label{tab:mag}
\footnotesize
\setlength{\tabcolsep}{4pt}
\renewcommand{\arraystretch}{1.3}
\begin{tabular}{|p{3.8cm}|c|p{6.2cm}|}
\hline
\textbf{Benchmark} & \textbf{Magnitude} & \textbf{Interpretation} \\
\hline
Lower bound --- pure network disruption \citep{korniyenko2017} & $\sim 0.7\%$ &
Supply-chain transmission of a localised shock, no policy intervention. \\
\hline
Upper bound --- most exposed economies, targeted tariff war \citep{teti2025} & $\sim 1.6\%$ &
Country-level loss for Canada/Mexico under the 2018--2025 US tariffs (vs.\ $-0.3\%$ for China), mapped onto the most exposed firms. \\
\hline
\end{tabular}
\end{table}

% =============================================================================
\section{Open Questions and a Forward Research Agenda}
\label{sec:agenda}
% =============================================================================

The review above leaves a small number of questions genuinely open and opens a larger
number of new ones. This section separates the two: first, the specific gaps in the
existing literature that \emph{this} paper is designed to close; second, the spin-off
questions that the same body of work raises but that lie beyond the present design and
would each support a distinct paper.

% ---------------------------------------------------------------------------
\subsection{The Gaps This Paper Fills}
\label{subsec:gaps_filled}
% ---------------------------------------------------------------------------

\begin{enumerate}
\item \textbf{Causal capability reconfiguration, not descriptive structure.} The EFC
literature (\citealp{bruno2023}; \citealp{pugliese2019}; \citealp{destefano2025})
measures block-modular capability structure but almost always descriptively. No prior
work links a \emph{predetermined, exogenous} trade-policy shock to the in/out
decomposition of firm diversification. This paper supplies that causal link.

\item \textbf{Composition, not performance.} The trade-and-productivity literature
(\citealp{pavcnik2002}; \citealp{topalova2011}; \citealp{amiti2007}; \citealp{tomasi2022})
measures how tariff shocks move productivity, markups, and survival. Because
performance dimensions decouple from capability (\citealp{grazzi2012};
\citealp{coad2021}; \citealp{syverson2011}), these outcomes are silent on \emph{which}
activities firms keep, drop, or enter. This paper moves the outcome to the architecture
of the export basket.

\item \textbf{Firm-level test of the chokepoint hypothesis.} The strategic-dependency
literature (\citealp{farrell2019}; \citealp{euc2021}; \citealp{arjona2023};
\citealp{korniyenko2017}) maps dependencies at the product and country level but has no
firm-level causal test. By interacting the tariff shock with a firm-specific
strategic-dependence exposure, this paper provides that test.

\item \textbf{A signed reading of exploration vs.\ exploitation under a real shock.}
March's (\citeyear{march1991}) tension has rich theory but little exogenous-shock
evidence at the firm level. The $\text{Share\_core}$ outcome under a predetermined tariff shock
delivers a directly signed estimate of which way firms move when their core is
threatened.

\item \textbf{Scalability of continuous-treatment event-study estimators.} We demonstrate
that a continuous-dose distributed-lag design with \citet{adao2019} shift-share inference---and
the heterogeneity-robust continuous-treatment estimators of \citet{dechaisemartin2020} and
\citet{callawaygoodmanbacon2024}---is tractable on a full-universe firm--product panel, a
methodological contribution in its own right for the heterogeneous-exposure shift-share
literature.
\end{enumerate}

% ---------------------------------------------------------------------------
\subsection{Beyond This Paper: Spin-Off Research Questions}
\label{subsec:spinoffs}
% ---------------------------------------------------------------------------

Each of the following is motivated by the reviewed literature but deliberately left
outside the present design; each would constitute a separate study.

\paragraph{(a) The firm--product--destination tensor.}
\citet{destefano2025} works in pure firm--product space; \citet{crozet2025} and
\citet{vannoorenberghe2016} show that the destination dimension carries its own
volatility structure. A natural extension promotes the bipartite firm--product matrix
to a firm--product--destination tensor and asks whether capability blocks reorganise
\emph{geographically}---whether retrenchment to the core is also a retreat to familiar
markets. This is the volatility-minimising hypothesis of $\Phi$-diversity made
firm-specific.

\paragraph{(b) Contemporaneous temporal-coupled partitions as a causal object.}
Our headline design freezes the baseline block, and our robustness lets it evolve on
\emph{pre-shock} information through the temporal multilayer machinery of \citet{mucha2010}
(with coupling and resolution chosen as in \citealp{bazzi2016}). What we deliberately leave
open is the \emph{contemporaneous} version: treating \emph{block reconfiguration itself}---the
birth, death, and merger of capability clusters measured with post-shock data---as the causal
outcome. Under our continuous treatment this is not identified, because there is no
never-treated control group from which to anchor an evolving map; a dedicated study could
recover it under a binary, staggered design by estimating the partition from not-yet-treated
firms, asking how trade shocks reshape the modular geometry of an entire industry rather than
the position of a single firm within a fixed geometry.

\paragraph{(c) The antidumping and safeguard stage.}
The Teti GTD \citep{teti2024} captures statutory MFN and preferential tariffs but
\emph{not} antidumping, safeguard, or trade-war tariffs. Pairing it with a temporary-trade-barriers
database would isolate the capability response to discretionary, firm-targeting
protection---plausibly a sharper and more anticipatory shock than scheduled tariff
changes, and a cleaner test of the escalation dynamic of \citet{evenett2024}.

\paragraph{(d) Instrumenting the policy-escalation component.}
\citet{evenett2024} document that industrial-policy intensity rises in election years
and that subsidies are matched tit-for-tat with 73.8 percent probability. An election-cycle
or tit-for-tat instrument for the policy-driven component of a firm's tariff shock
would separate the escalation channel from the mechanical tariff change---turning the
escalation caveat of this paper into an identification strategy of its own.

\paragraph{(e) Green-capability reconfiguration.}
\citet{mealy2022} show the complexity toolkit applies to any product subset and that
green-capability accumulation is strongly path dependent. Re-running the present design
with the green-product set as the treated category would ask whether decarbonisation
policy and trade policy push firms in the same or opposite capability directions---a
direct bridge between the strategic-autonomy agenda of \citet{cagnin2021} and the green
transition.

\paragraph{(f) Institutional voids as a moderator.}
\citet{delios2008} and \citet{boschma2015} show that the value of diversification, and
the related-versus-unrelated balance, depend on institutional quality. Exploiting
within-China regional institutional variation as a second moderator would test whether
capability retrenchment under trade shocks is steeper or shallower where market
institutions are weaker, distinguishing capability exploitation from
institutional substitution as motives for diversification.

\paragraph{(g) The gain/loss asymmetry of complexity.}
\citet{nomaler2023} show that lost specialisations are often \emph{more} complex than
gained ones, and that density-based relatedness is heavily confounded by diversity and
ubiquity. A study that scores firm-level product \emph{exits} by complexity---asking
whether tariff shocks make firms shed their most sophisticated products first---would
test whether trade shocks erode the capability frontier from the top, a sharper and
more policy-relevant question than net diversification counts can answer.
 into the main paper.
% Cite keys follow the convention author(s)year; matching @article / @book /
% @techreport entries live in Bibliography_base.bib (86 entries as of 2026-06-05).
%
% STRUCTURE
%   2   Related Literature
%     2.1  Economic complexity, capability structure, and diversification theory
%     2.2  Capability theory, corporate coherence, and the direction of adjustment
%     2.3  Trade policy shocks and firm adjustment
%     2.4  Strategic dependencies and geopolitical trade shocks
%     2.5  Identification
%     2.6  Contribution
%   3   Expected Dynamics: What the Literature Predicts We Should Observe
%     3.1  The out-of-block margin (d_out)
%     3.2  The in-block margin (d_in) and retrenchment
%     3.3  Share_core and the direction of adjustment
%     3.4  Heterogeneity and the strategic-dependence moderator
%     3.5  Synthesis: hypotheses, signs, and magnitude benchmarks  [2 tables]
%   4   Open Questions and a Forward Research Agenda
%     4.1  The gaps this paper fills
%     4.2  Beyond this paper: spin-off research questions
%
% NOTE: tables use plain \hline tabular (no booktabs dependency) so the file
% compiles stand-alone via lit_review_test.tex and when \input into the paper.
% =============================================================================

\section{Related Literature}
\label{sec:lit}

This paper sits at the intersection of four bodies of work: the economic-complexity
literature on firm-level capability structure and the theory of diversification; the
capability-based theory of the firm and evolutionary economics; the empirical trade
literature on the effects of tariff shocks on firm outcomes; and the emerging literature
on strategic dependencies, geopolitical risk, and the weaponisation of economic
interdependence. A fifth, methodological strand covers identification in shift-share
designs with heterogeneous treatment effects. The unifying thread across these
literatures is a single dynamic: an external shock that degrades the value of a firm's
existing activities forces a \emph{reconfiguration} of its productive knowledge, and
the literatures disagree on whether that reconfiguration deepens a coherent core or
pushes outward into new capability bundles. We discuss each stream in turn, emphasising
what each predicts about the \emph{direction} and \emph{persistence} of adjustment, and
close with a statement of our contribution.

% ---------------------------------------------------------------------------
\subsection{Economic Complexity, Capability Structure, and Diversification Theory}
\label{subsec:lit_efc}
% ---------------------------------------------------------------------------

The product-space framework of \citet{hausmann2007} and \citet{hidalgo2009}
established that countries export products in patterns that are far from random:
goods that require similar underlying capabilities tend to be co-exported, generating
a proximity structure in the product space that conditions the path of structural
transformation. \citet{hidalgo2011} formalise this intuition through a bipartite
capabilities model in which countries and products are represented as the two node
sets of a bipartite graph, and the adjacency matrix $M_{cp}=1$ whenever country $c$
has revealed comparative advantage (RCA) in product $p$. The key result is that
diversification is convex in the capability stock: countries with few capabilities
earn near-zero marginal returns from acquiring any single new one, generating the
low-income quiescence trap that has been documented empirically across the
income distribution \citep{hausmann2007}. This non-linearity is fundamental to the
agenda of using complexity tools in a causal setting: treatment-effect heterogeneity
by capability endowment is an architectural feature of the framework, not an
add-on.

The Fitness-Complexity (FC) algorithm of \citet{cristelli2013} provides the
operational complexity metric adopted in the subsequent firm-level literature. Rather
than the linear iterative procedure of \citet{hidalgo2009}, FC defines country fitness
as the sum of product complexities and product complexity as the harmonic mean of
exporter fitness---a bounded operator that uniquely rewards diversification and
assigns complexity to a product according to the least capable exporter that ships it,
giving the algorithm its asymmetric censoring property. The algorithm converges to a
unique fixed point and, as shown by \citet{cristelli2017}, generates a selective
predictability scheme in which high-fitness countries follow laminar, predictable
growth trajectories while low-fitness countries exhibit chaotic dynamics---a property
that has direct implications for heterogeneous treatment responses in our setting:
firms embedded in high-complexity capability clusters should display more systematic
responses to tariff shocks than firms operating at the periphery of the product
space. \citet{angelini2017} characterise the long-run dynamics of products in the complexity
plane, showing that they drift toward a stationary configuration so that
shock-induced dislocations are transient and mean-reverting---which is why we treat
product complexity as a well-defined pre-shock control and read post-shock movements
against a fixed baseline.
\citet{stojkoski2025} extend the framework to an optimisation context,
showing that the ECI-growth relationship is stronger at lower income levels and that
relatedness is one of several predictors of product-entry, a finding that validates
the multi-dimensional block structure we exploit over simpler relatedness-only
measures. \citet{koch2021} sounds a useful caution about what such metrics measure: the
complexity--growth link is in fact \emph{stronger} when complexity is computed on
value-added rather than gross exports ($R^2 = 0.63$ versus $0.31$), so gross-flow
complexity blurs the connection to realised value capture. Since our firm-customs
measures are necessarily built on gross flows, we read them as capturing capability
\emph{structure} rather than realised value capture, which is precisely why we treat
the block partition as a measure of capability architecture rather than of performance.

Three strands sharpen what ``relatedness'' and ``diversification'' actually mean for
the direction of adjustment. First, the related-variety tradition of
\citet{frenken2007} distinguishes \emph{related} variety (entry into activities that
share capabilities, which drives regional \emph{employment} growth through
within-sector Jacobs spillovers---though, tellingly, not productivity growth) from
\emph{unrelated} variety (which provides a portfolio hedge that lowers unemployment
under asymmetric sectoral shocks). This is the regional-economics antecedent of our
$d_{\rm in}$ versus $d_{\rm out}$ distinction. \citet{boschma2015} show that the relative payoff to
related versus unrelated diversification is institutionally contingent: coordinated
market economies favour incremental related diversification, while liberal market
economies support the unrelated, exploratory kind, a result that motivates treating
the regional production system, not just the firm, as a unit of heterogeneity.
\citet{freire2019} provides a structural model of endogenous diversification in which
capability accumulation and product entry co-evolve, giving microfoundations to the
empirical relatedness regularities. Second, \citet{mealy2022} demonstrate that the
complexity toolkit transfers to any policy-relevant subset of products---their
green-complexity indices show strong path dependence in capability accumulation
(a $0.92$ correlation between green-complexity potential and the green-complexity
index, with beginning-of-period potential predicting subsequent green-export
growth), a template we adopt when restricting attention to strategically dependent
product categories. Third, \citet{nomaler2023} issue a confounding warning that
directly disciplines our outcome construction: density-based relatedness metrics are
$\sim$93 percent explained by diversity and ubiquity alone, so a measured ``relatedness''
effect can be an artefact of how many products a firm makes and how common those
products are. Their symmetric treatment of RCA \emph{gains} and \emph{losses}, together with
the finding that lost specialisations are often \emph{more} complex than gained
ones, reframes diversification as a two-sided entry-and-exit process and is one of the
closest conceptual antecedents to our reconfiguration question.

The application of complexity tools to firm-level data is recent. \citet{bruno2018, bruno2023} construct firm-product bipartite
networks from Customs microdata and show that RCA-based binarisation
(with BiCM entropy-null correction as a robustness) yields block-modular structures
that are stable across years and meaningfully predict firm growth, innovation, and
survival. \citet{pugliese2019} show that firms diversify coherently within their
existing technological portfolio, so that a firm's position within its capability
cluster is informative about its future mobility---a finding that motivates our
within-block heterogeneity analysis; the country-level analogue is \citet{pugliese2017},
who show that complex economies escape the low-income trap by moving laterally into
related capabilities. \citet{coad2021} push
this furthest, applying the FC algorithm to a firm$\times$capability matrix for $\sim$40{,}000
Indian firms and recovering a nested ``capabilities ladder'' whose rungs predict growth
and survival but---tellingly---not profitability, a decoupling we return to in
Section~\ref{subsec:lit_cap}. \citet{destefano2025}
provide the most direct methodological antecedent for our outcome measures: they
decompose firm diversification into in-block ($d_{\rm in}$) and out-of-block
($d_{\rm out}$) counts using the block partition produced by the BRIM algorithm of
\citet{barber2007}, define $\text{Share\_core} = d_{\rm in}/(d_{\rm in}+d_{\rm out})$,
and show that out-of-block diversification is the growth-positive margin while
retrenchment toward the core---the $d_{\rm in}$ direction---is the
growth-negative, lock-in-reinforcing margin. We adopt their exact outcome measures
and employ the same partition algorithm.

The network-method foundations underpinning the partition are well established.
\citet{newman2006} introduced the modularity objective and its spectral optimisation;
\citet{blondel2008} provide the Louvain heuristic that makes community detection
tractable at the scale of full-universe firm-product matrices; \citet{ducharenas2005}
supply the extremal-optimisation alternative we use as a partition robustness check;
and \citet{fortunato2016} provide the authoritative critical survey of the field's
validation pitfalls. \citet{barber2007} adapts modularity to the bipartite case (BRIM),
and \citet{soleribalta2018}
show that the BRIM objective is insensitive to within-block
nested structure and can therefore miss the hierarchical ranking of firms inside a
capability block; their in-block nestedness quality function $I$ jointly optimises the
block partition and within-block NODF \citep[the consistent nestedness metric of][]{almeidaneto2008},
providing the basis for our within-block rank ordering of firms. \citet{mariani2019} prove that the FC algorithm is a nestedness
maximiser and clarify the relationship between modularity and nestedness in firm-level
networks: the two are positively correlated at low network densities (the regime typical
of firm-product matrices) and can be jointly estimated. Validation of the BRIM
partition against entropy null models---using the Bipartite Configuration Model of
\citet{squartini2011} and \citet{saracco2015}, with the projection reconciliation of
\citet{cimini2022}---is a required pre-processing step
to ensure that block structure reflects genuine capability clustering rather than
degree-sequence artefacts, and we implement it throughout. A wider set of Economic Fitness and
Complexity (EFC) studies---broader applications of economic-complexity methods---rounds
out the methodological base: global-value-chain specialisation \citep{bontadini2021, bontadini2025},
the complexity of structural change and the sustainability transition
\citep{caldarola2024a, caldarola2024b}, machine-learning reconstruction of trade flows
\citep{gnecco2023}, and relatedness-based forecasting of industrial upgrading
\citep{albora2023a, albora2023b}.

% ---------------------------------------------------------------------------
\subsection{Capability Theory, Corporate Coherence, and the Direction of Adjustment}
\label{subsec:lit_cap}
% ---------------------------------------------------------------------------

The theoretical framework connecting firm capability stocks to the direction of
diversification draws from the resource-based view of the firm, evolutionary economics,
and the learning-and-adaptation literature. \citet{penrose1959} establishes that
firms are bundles of productive resources whose services can be redeployed across
activities at varying cost: the excess services that routinely arise from indivisible
resources create an organic pressure toward related diversification that is internal
to the firm, not driven by market structure. \citet{teece1982} formalises this in a
scope-economy model showing that the boundaries of the multi-product firm are
determined by the jointness of underlying knowledge inputs; firms expand into activities
that share tacit knowledge with existing lines, generating the coherence property
that the EFC literature quantifies through block-modular partitions.
\citet{teece1994} make this measurable through the survivor-principle relatedness
metric---sectors are ``related'' when more firms combine them than chance predicts---which
is the product-space analogue of the bipartite co-occurrence logic at the heart of the
EFC programme. The corporate-coherence literature that builds on it is directly
informative about the \emph{dynamics} we expect to observe. \citet{piscitello2000}
finds that large firms reduced product coherence while \emph{maintaining}
technological coherence over 1977--1995: capabilities are the stable organising
principle even as product portfolios churn, which is exactly why a capability-block
partition should track adjustment more faithfully than a raw product count.
\citet{valvano2003} provide the most pointed antecedent for a trade-shock prediction:
in Italian manufacturing, firms in sectors more exposed to single-market competition
(``sensitive'' sectors) became \emph{more} coherent over time, and entrants into the
leadership ranks were more coherent than incumbents---direct evidence that intensified
competition pushes surviving firms toward their core ($d_{\rm in}$). \citet{delios2008}
add the institutional contingency: within-country product diversification is more
valuable where market institutions are weak, because diversification substitutes for
missing capital, labour, and information markets---a moderation channel that warns
against reading every diversification response as pure capability exploitation.

The corporate-finance literature on the ``diversification discount'' supplies the
causal-direction caution that disciplines our interpretation. \citet{langstulz1994}
document a robust industry-adjusted discount in Tobin's $q$ of 0.35--0.50 for
diversified US firms throughout 1978--1990, but show that firms diversify \emph{after}
their core industry weakens---the discount precedes diversification rather than being
caused by it. \citet{matsusaka2001} rationalises this as capability-match search: a
firm whose current match deteriorates maintains its existing lines while experimenting
with new ones (the ``$b$-strategy''), liquidating and restarting only when the match
falls below a lower threshold. \citet{bernardo2002} recast the same fact as
real-option value: single-segment firms are worth more because they retain unexercised
expansion options. Translated to our setting, these models deliver a sharp,
non-obvious prediction: a \emph{moderate} tariff shock should push firms toward
exploratory $d_{\rm out}$ search (the $b$-strategy), whereas a \emph{severe} shock that
drives match quality below the liquidation threshold collapses the option value of
organisational capabilities and forces retrenchment or exit---so the sign of the
average response is genuinely ambiguous and should depend on shock intensity, exactly
the non-linearity our event study is designed to detect.

The evolutionary perspective of \citet{dosi1988, dosi1995} adds a regime dimension:
the rate and direction of capability accumulation depend on both the firm's own
endowment and the technological regime of the sector---the opportunity conditions,
appropriability, and cumulativeness that structure search. Under Schumpeterian Mark
II conditions (cumulative, large-firm-dominated regimes), incumbents sustain
technological advantages through internal learning, generating the technological
lock-in that maps directly onto high $\text{Share\_core}$. The implication for our
setting is that firms in complex, cumulative technological regimes should display
stronger $d_{\rm in}$ responses to adverse shocks---retrenchment to the core is both
the natural adjustment and the organisationally feasible one when search capabilities
are already embedded in existing product clusters.

\citet{march1991} provides the ambiguity at the heart of our hypotheses. The
exploitation--exploration tension implies that the direction of capability adjustment
under stress is not a priori signed: firms that face an adverse demand or competitive
shock in their existing core can either intensify exploitation of known competencies
(retrenchment) or accelerate exploration beyond the core (diversification-as-escape).
March's formal result is that organisations systematically over-exploit under most
plausible parameter configurations, predicting the retrenchment (H1) hypothesis as
the modal response but leaving open the exploration (H1$'$) case for firms with
sufficient slack and absorptive capacity. \citet{patel1997} document the empirical
content of this tension at the firm level, showing that large innovative firms spread
R\&D effort across fields but in highly persistent patterns---path dependence is the
rule, not the exception---while \citet{bell1993} establish that technological
capability accumulation in developing-country firms is even more path-dependent,
with explicit investments in learning required to escape existing trajectories.
\citet{storper1992} connects this to trade: product-based technological learning
generates agglomeration economies in particular product-territory combinations, so
that an adverse trade shock removes both demand and the collective learning
infrastructure that sustains capability accumulation---a channel that implies
persistent, not merely transitory, capability effects of tariff shocks.

The firm-growth literature, summarised in \citet{coad2009} and documented empirically
in \citet{bottazzi2003, bottazzi2005, bottazzi2006}, shows that firm-level growth is
Laplacian---fat-tailed, near-unit-root---and that the variance of growth rates scales
inversely with firm size at a rate explained entirely by diversification across product
lines. The scope-economy mechanism \citep{bottazzi2006} implies that firms which
diversify within a related cluster (high $d_{\rm in}$) reduce the co-movement of
their individual product revenue streams more effectively than firms that scatter across
unrelated products. Under tariff shock conditions, firms that lose access to existing
export markets face a portfolio rebalancing problem whose solution---from a pure
variance-minimisation standpoint---is to deepen within-block exposure rather than
scatter to unfamiliar destinations, reinforcing the retrenchment prediction. A
recurring caveat from this literature is that performance dimensions decouple:
\citet{grazzi2012} shows that Italian exporters carry a large productivity premium but
\emph{no} profitability premium, and \citet{coad2021} that capability-ladder position
predicts growth and survival but not profit margins; \citet{syverson2011} surveys the
enormous, persistent productivity dispersion that these patterns sit inside. The
collective lesson---that capability advantage need not show up in profits or even in
growth---is why we measure the \emph{architecture} of the export basket
($d_{\rm in}$, $d_{\rm out}$, $\text{Share\_core}$) rather than a performance scalar.
\citet{aghion2001, aghion2005} add the innovation margin: neck-and-neck competitors
facing increased competitive pressure innovate to escape competition, while laggards
reduce innovation effort. Translating to our setting, firms with high within-block
fitness ranks face incentives to deepen core capabilities under tariff pressure (the
escape-competition channel), while low-rank firms---unable to sustain the cost of
maintaining their current product portfolio---are more likely to prune $d_{\rm out}$
and consolidate around the most defensible core products. This generates the
within-block heterogeneity in treatment effects that we exploit empirically.

% ---------------------------------------------------------------------------
\subsection{Trade Policy Shocks and Firm Adjustment}
\label{subsec:lit_trade}
% ---------------------------------------------------------------------------

A large reduced-form literature establishes that tariff changes cause significant
adjustments in firm outcomes. \citet{pavcnik2002} and \citet{topalova2011} document
that trade liberalisation raises TFP through competitive selection and technology
imports in Chilean and Indian manufacturing, respectively. \citet{amiti2007} show,
for Indonesian firms, that input tariff reductions deliver TFP gains at least twice as
large as those from output tariff reductions---a finding driven by access to high-quality
imported intermediates---and construct the shift-share input-tariff exposure variable that is
the methodological template for our treatment. \citet{tomasi2022} replicate this
design for Vietnamese firms during the 2007 WTO accession, finding that governance
quality moderates the tariff-TFP effect, providing the closest precedent for our
moderated shift-share design with ExpDep as the heterogeneity variable.

What this literature consistently measures is performance: productivity, employment,
markups, survival. It is essentially silent on the \emph{composition} of a firm's
product portfolio---which activities are retained, dropped, or newly entered, and
whether the reconfiguration deepens a coherent core or spreads capabilities into
new domains. The EFC literature reviewed in Section~\ref{subsec:lit_efc} supplies
the tools to measure that composition, but has applied them almost exclusively in
descriptive, non-causal settings. This paper occupies the intersection left empty
by both strands. The theoretical channel linking trade exposure to capability
adjustment is supplied by the trade-and-innovation literature: \citet{akcigit2022}
unify the market-size, business-stealing, and escape-competition effects of trade on
innovation incentives, showing that the net effect on a firm's capability investment
is sign-ambiguous and depends on its competitive position---the firm-level analogue of
the exploitation--exploration tension. \citet{riverabatiz1991} establish the deeper
point that integration affects growth through flows of \emph{ideas} as well as goods,
so that a tariff shock that severs a firm from a technological frontier market can have
permanent (growth-rate), not merely level, effects on its capability trajectory---reinforcing the
persistence prediction from \citet{storper1992}.

The most direct evidence that trade shocks restructure capability portfolios rather
than merely rescaling performance comes from \citet{fontagne2018}, who show that
Chinese export quotas on textiles and clothing generated a significant \emph{Typical
Product Vector} consolidation among French exporting firms: products at the core of
a firm's stable portfolio were retained at a 7.7 percentage-point higher rate than
peripheral products in treated sectors. This is, in effect, a $d_{\rm in}$ response
to a trade-policy shock measured with product-vector tools---the closest antecedent
in the literature for our main outcome. \citet{fontagne2022} provide the
product-level trade elasticities that we use as Bartik weights: 5,050 HS6-digit
elasticities estimated via PPML gravity, averaging $-5.3$ (median $-3.5$), with
HS-section means spanning $-18.5$ (minerals), $-6.1$ (machinery), and $-3.6$
(footwear), capturing the heterogeneous disruption that a uniform tariff schedule
imposes on firms with different product portfolios.

The dynamic and firm-level dimensions of portfolio adjustment are examined by
\citet{crozet2021}, who study the effect of EU and US sanctions on the extensive
margin of French exporters to Iran and Russia. Using monthly customs data and a
dynamic probit model with three-way fixed effects and analytical bias correction,
they find that sanctions lower the probability of serving a sanctioned market by 39
percent (Iran) and 23 percent (Russia), with strong state dependence: the entry-cost
factor is 19--23 percent higher for firms not active in the previous year. Crucially,
the lifting of sanctions on Myanmar generates only a small and gradual positive response,
establishing an \emph{asymmetry} between imposition and removal that implies permanent,
not transient, $d_{\rm out}$ contraction even after the policy shock is reversed.
This hysteresis result---consistent with the Roberts--Tybout sunk-cost model---means
that treatment effects in our event study should be interpreted as reflecting a
permanent restructuring of the capability portfolio, not a temporary demand disruption.
The firm-level heterogeneity documented by \citet{crozet2021}---that trade-finance-dependent
firms, inexperienced firms, and non-specialist firms exit more sharply---maps
directly onto the moderation hypotheses we test: private Chinese firms (documented
by \citet{guariglia2011} as the most financially constrained in the NBS panel) should
display amplified $d_{\rm out}$ reductions, while SOEs with soft budget constraints
should display more muted responses.

A complementary strand examines the export portfolio diversification strategies that
firms deploy to manage demand volatility. \citet{vannoorenberghe2016} show, for
Chinese firm-level data, that destination diversification reduces volatility for large
exporters but \emph{increases} it for small ones, with the switching point determined
by a composition effect: occasional exports to marginal destinations dominate the
diversification gain for small firms, while stable core-market exports dominate for
large ones. \citet{crozet2025} formalise this puzzle through a phyloeconomic diversity
index $\Phi$ that adds a dissimilarity dimension to the standard richness-evenness
decomposition. Their central finding---that conditional on the number of destinations
and the Herfindahl concentration of export sales, higher $\Phi$ (more dissimilar
markets) is associated with \emph{higher} export sales volatility---reveals that
hedging through geographic diversification has limits: distant, structurally different
destination markets carry higher idiosyncratic demand variance, so that firms forced
to substitute familiar Western markets with unfamiliar alternatives do not achieve
the volatility reduction that naive diversification logic predicts. For our setting,
this implies that Chinese firms which maintain $d_{\rm out}$ by pivoting to non-Western
destinations face not only lower export levels but higher volatility---a double cost
that reinforces the variance-minimising case for $d_{\rm in}$ consolidation.
\citet{barbieri2024} provide a parallel decomposition at the territorial level:
\emph{sectoral} diversification (a proxy for productive capabilities, the conceptual
sibling of $d_{\rm in}$) raises export intensity, while \emph{geographic}
diversification (a proxy for geoeconomic capabilities, loosely analogous to the
resilience role we attribute to $d_{\rm out}$) is associated with greater export
stability; institutional quality independently improves both margins.\footnote{The
analogy is functional rather than literal. Barbieri's second axis is
\emph{destinations}, whereas our $d_{\rm out}$ remains a \emph{product}-space
measure (out-of-block products), following \citet{destefano2025}, who works in pure
firm--product space with no destination dimension. The roles align---one margin scales
output, the other buys resilience---but the underlying axes differ, and in their data
the dominant stabiliser is EU market exposure rather than geographic diversification
per se. We therefore treat the destination dimension as a second-stage extension
(Section~\ref{subsec:spinoffs}) rather than as a component of $d_{\rm out}$.}

We estimate treatment using the Global Tariff Database (GTD) of \citet{teti2024},
the first applied-tariff panel to cover bilateral HS6 tariffs continuously from 1988
to 2021 while correcting for the false interpolation and selection artefacts of UNCTAD
TRAINS. \citet{teti2019} establish, using structural gravity, that Rules-of-Origin
compliance is prohibitively costly for the machinery and electronics sectors at the
core of Chinese exports, confirming that MFN tariffs---the rate captured by the
GTD---are the economically relevant wedge for our firms. The welfare magnitudes from
\citet{teti2025} calibrate our expected effect sizes: a GE simulation of the
2018--2025 US--China tariff war estimates cumulative real-income losses of roughly
$-1.6$ percent for Canada and Mexico---economies with high export exposure to the US
market---against only $-0.3$ percent for China in the aggregate. We use the
$-1.6$ percent figure as an external benchmark for the most exposed Chinese firms:
by analogy with the highly exposed Canadian and Mexican case, the uneven distribution
of tariff exposure across product portfolios should push firms concentrated in
high-regulatory-dependence categories toward magnitudes well above the national
average. The mapping from these country-level losses to firm-level exposure is our
own; \citet{teti2025} simulate countries and US states, not firms. The distortionary
role of size-contingent and sector-specific regulation documented by
\citet{garicano2016}---who find welfare costs of 1.3--3.5 percent of GDP from French
size thresholds---is the domestic-policy analogue of this exposure heterogeneity and
motivates controlling for regulatory environment in the moderation analysis.

% ---------------------------------------------------------------------------
\subsection{Strategic Dependencies and Geopolitical Trade Shocks}
\label{subsec:lit_geo}
% ---------------------------------------------------------------------------

The design of the heterogeneity variable ExpDep draws from a separate and newer
literature on strategic economic dependencies. \citet{farrell2019} identify the
conditions under which the topology of global economic networks generates geopolitical
leverage: hub-and-spoke structures create chokepoint nodes through which flows can be
controlled (flow restriction) and from which information can be extracted (panopticon
effects). Crucially, weaponisation requires two joint conditions---jurisdictional
control over the chokepoint and domestic institutions capable of deploying that control
as a foreign policy instrument---which limits the set of products and bilateral
relationships that qualify as strategic in the relevant sense. The Farrell--Newman
framework grounds our moderation hypothesis: Chinese firms whose export baskets are
concentrated in products where Western jurisdictions have identified chokepoint control
face not a neutral tariff shock but one accompanied by active policy interest in
reducing dependence---generating an amplified and persistent treatment effect.

The quantitative operationalisation of this insight is provided by \citet{euc2021},
who develop the Core Dependency Indicators (CDI) methodology to identify EU strategic
dependencies at the HS6-product level. The methodology applies three thresholds:
supplier concentration (HHI of extra-EU import shares exceeding 0.4, implying fewer
than 2.5 effective suppliers), import reliance (extra-EU imports exceeding 50 percent
of total EU imports of the product), and substitutability (extra-EU import volume
exceeding total EU export volume, so that domestic production cannot cover the import
volume even in an emergency). Of approximately 5,000 HS6 products, 137 pass all three
thresholds in the EU's most sensitive ecosystems, with China accounting for roughly
52 percent of their import value. \citet{arjona2023} refine the methodology with a
dynamic four-year threshold and a re-export correction: accounting for transit trade
\emph{raises} the count of dependent products by roughly 27 percent, since ignoring
re-exports understates true extra-EU reliance, yielding a dependent-product set of
204 items that provides the product classification we use to construct ExpDep.
The forward-looking scenario work of \citet{cagnin2021} situates these indicators in
the EU's Open Strategic Autonomy agenda, projecting an escalation of dependency-reducing
policy through 2040 that makes the persistence of any treatment effect we estimate a
policy-relevant, not merely statistical, question.

The network-science analogue of the CDI methodology is provided by
\citet{korniyenko2017}, who combine three structural indicators of product trade networks---
the standard deviation of weighted outdegree centrality (capturing hub dominance), the
product of clustering coefficient and network diameter (capturing geographic
compartmentalisation), and a Revealed Factor Intensity measure of substitutability---to
classify 3,578 intermediate and capital goods into four risk tiers via k-median
clustering. Their persistently risky group comprises 421 products over 2003--2014,
concentrated in machinery and mechanical appliances (HS 84--85), transport equipment,
pharmaceuticals, and precision instruments. The causal validity of the classification
is demonstrated by two out-of-sample cases: the 2011 Japan earthquake and Thai floods,
where products subsequently in short supply were correctly identified as Group~4 one
year prior. The overlap between the Korniyenko Group~4 products and the
Arjona CDI-dependent set defines the most extreme ExpDep exposure: products
that are simultaneously structurally concentrated, topologically fragile, and actively
targeted by EU and US reshoring policy.

\citet{evenett2024} provide real-time evidence on the policy escalation dynamic. Their
New Industrial Policy Observatory (NIPO), covering 2,500$+$ new industrial policy
measures across 75 jurisdictions in 2023, documents a 73.8 percent probability that a
subsidy by one major economy (China, EU, or US) is matched by a subsidy in the same
product by another within 12 months. PPML regressions confirm that IP usage correlates
positively with RCA---governments target sectors where they already have comparative
advantage, consistent with political-economy capture rather than market-failure
correction. For our event study, this finding implies that the Teti GTD tariff
shock at time $t$ is not a clean one-off treatment but the first signal in an
escalating policy sequence: treatment intensity grows after the event-window, pushing
treatment effects toward permanence.

\citet{fontana2024} embed these observations in a political-economy theory of
strategic-assets policy, drawing on the three logics of capability blockade formalised
by Ding and Dafoe: cumulative-strategic blocking (targeting technological RCA through
standards exclusion and export controls), infrastructure-strategic blocking (targeting
critical material supply chains through reshoring subsidies), and dependency-strategic
blocking (exploiting existing dependencies through sanctions and FDI screening). The
taxonomy provides a structural basis for splitting the Teti tariff sample by motive:
tariff shocks in sectors matching the cumulative-strategic logic (advanced
manufacturing, semiconductors) should generate stronger $d_{\rm in}$ responses than
shocks in sectors matching the infrastructure-strategic logic (raw materials, chemicals),
where substitution is constrained by the CDI3 condition. We use this taxonomy in a
sample-split robustness exercise.

% ---------------------------------------------------------------------------
\subsection{Identification}
\label{subsec:lit_id}
% ---------------------------------------------------------------------------

Our identification strategy rests on a single organising principle: \emph{everything
that could mechanically respond to the shock must be predetermined relative to it.}
The estimand is the within-firm response of an export-capability portfolio to a
trade-policy shock, and the credibility of that estimand depends on three objects
being fixed before the shock---the exposure weights, the disruption elasticities, and
the capability map against which the portfolio is read. We discuss each, and then
state precisely which of three admissible network constructions we adopt and why.

\paragraph{A continuous, time-varying treatment.}
Following the input-tariff shift-share of \citet{amiti2007}, \citet{topalova2011}, and
\citet{tomasi2022}, we construct the tariff exposure variable as
\[
  \text{TariffShock}_{it}
  = \sum_p s_{ip,\text{pre}}\cdot \varepsilon_p \cdot \Delta\ln(1+\tau_{pt}),
\]
where $s_{ip,\text{pre}}$ are firm--product export shares in a strictly pre-shock
baseline window, $\Delta\ln(1+\tau_{pt})$ are changes in bilateral applied tariffs
from the Teti GTD, and $\varepsilon_p$ are the product-level trade elasticities of
\citet{fontagne2022}. We treat $\text{TariffShock}_{it}$ as a \emph{continuous,
time-varying, non-absorbing} dose rather than a binary staggered treatment: the dose
grows and recedes as tariffs move, and discretising it into a treated/control
indicator would discard precisely the cross-firm variation in exposure intensity that
identifies the dose--response. This choice is deliberate and consequential, as we make
clear below. The elasticities $\varepsilon_p$ are estimated by \citet{fontagne2022} on
the universe of 152 importers and 189 exporters---they are a property of the
\emph{product}, external to our firm sample, and hence cannot be a function of the
behaviour we study; products with a non-significant estimate inherit their HS-section
mean. The predetermined shares satisfy the Bartik exogeneity condition, which we defend
on \emph{both} of its routes: from the shares, following \citet{goldsmithpinkham2020},
and from the shocks, following \citet{borusyakhull2022}, reporting Rotemberg-weight
diagnostics and a recentred-instrument specification.

\paragraph{Dynamics and inference for a continuous dose.}
Because the treatment is continuous and non-absorbing, the cohort$\times$relative-time
machinery of \citet{sunabraham2021} is not the natural primary specification---its
cohorts are defined by the timing of a binary switch. We therefore estimate the
dynamics through a distributed-lag event study, $Y_{it}=\sum_k \beta_k\,
\text{TariffShock}_{i,t-k}+\alpha_i+\delta_t+\varepsilon_{it}$, in which the lead
coefficients ($k<0$) constitute the parallel-trends test and the lag coefficients
($k>0$) trace the persistence (hysteresis) prediction. Inference follows the
shift-share-robust standard errors of \citet{adao2019}, which are native to continuous
exposure designs and which we report as primary against firm- and sector-clustered
alternatives. As robustness we re-estimate the dynamics with the continuous,
non-absorbing estimator of \citet{dechaisemartin2020}, the continuous-treatment
difference-in-differences framework of \citet{callawaygoodmanbacon2024}, and---after
discretising exposure into terciles---the heterogeneity-robust estimators of
\citet{sunabraham2021}, \citet{callawaysantanna2021}, and \citet{borusyak2024}. We
flag the assumption that the continuous design carries a \emph{stronger} parallel-trends
requirement than the binary case---\citet{callawaygoodmanbacon2024} show that
identifying the dose--response slope requires the absence of selection into dose level,
which we probe but cannot fully test.

\paragraph{The capability map: three admissible constructions, and the one we adopt.}
The outcomes $d_{\rm in}$, $d_{\rm out}$, and $\text{Share\_core}$ are read off a
block partition of the firm--product export network \citep[the export, not production,
basis of][]{destefano2025}, obtained by bipartite modularity maximisation
\citep{barber2007} on an RCA-binarised incidence matrix validated against the Bipartite
Configuration Model \citep{squartini2011,saracco2015,cimini2022}. The non-trivial
question is \emph{when} the partition is estimated, because the map is itself the
measuring instrument for the outcome, and a map estimated on post-shock data would
co-move with the treatment. Three constructions are admissible:
\begin{enumerate}
  \renewcommand{\labelenumi}{(\roman{enumi})}
  \item \textbf{Frozen baseline partition.} The partition $b(p)$ and each firm's core
        block $c(f)$ are estimated once, on the pre-shock baseline window, and held
        fixed; only the firm's RCA incidence $M_{fpt}$ varies over time. The
        panel variation in $d_{\rm in}/d_{\rm out}$ is then unambiguously the firm
        relocating against a frozen, predetermined grid.
  \item \textbf{Predetermined evolving partition.} The map is allowed to evolve, but
        estimated only on \emph{pre-shock} slices through the temporal multilayer
        modularity of \citet{mucha2010}---with inter-slice coupling and resolution
        chosen by partition stability \citep{bazzi2016}---and then extrapolated forward.
        This nests (i) as the strong-coupling limit and captures any pre-existing drift
        in the capability space, at the cost of measurement error that grows with the
        post-shock horizon.
  \item \textbf{Contemporaneous multilayer.} The map evolves using all years including
        post-shock. Under our continuous treatment this construction is \emph{not}
        causally identified: with no never-treated firms, there is no
        treatment-independent control group from which to anchor the map, so the
        partition---and hence the $(d_{\rm in}+d_{\rm out})$ denominator---remains a
        function of the dose. We therefore use it only descriptively, to document how
        the capability space reorganises, and never as the causal outcome.
\end{enumerate}
\textbf{We adopt (i) as the headline specification}, precisely because it has the
smallest researcher degrees of freedom and the cleanest exogeneity argument: the map
literally precedes the shock and cannot respond to it. Its one weakness---that a frozen
map grows stale---generates classical measurement error and therefore attenuation
\emph{against} our hypotheses, an honest bias rather than a confounding one. We
neutralise even that concern by reporting (ii) as robustness, which lets the map evolve
on pre-shock information and leaves the estimates intact, and by re-estimating (i) with
a leave-one-out partition that excludes each focal firm from the map that classifies
it, breaking the reflexive channel. Two construction rules are stated explicitly and
checked: products with no baseline block are assigned by product-space proximity (or
dropped, with the affected export share reported), and firms born after the baseline
are handled within the exit-selection analysis rather than assigned a post-baseline
core. We compute RCA with the standard contemporaneous, \emph{within-China} Balassa
denominator: because the capability blocks are themselves built from the Chinese
firm--product network, a China-relative benchmark---specialisation relative to the
Chinese export economy---is the internally consistent and economically natural object,
and we retain it as the headline measure. The denominator carries a second-order
\emph{inference} concern rather than an economic one: under a sector-wide shock the
aggregate Chinese export of a treated product falls for all exposed firms, mechanically
inflating the within-sample benchmark and potentially moving threshold crossings for
reasons unrelated to the firm's own portfolio choice. We treat this as robustness, not
as a change of measure. We re-estimate with a leave-one-out within-China denominator,
which preserves the China-relative interpretation while removing the firm's reflexive
contribution, and, as a more aggressive check, against a rest-of-world (ex-China) BACI
benchmark, which trades the China referent for full exogeneity to the Chinese firms'
treated response.

\paragraph{An elegant alternative under a modified treatment.}
We note one design under which the \emph{contemporaneous} evolving map (iii) would
become causally admissible: were we to trade the continuous dose for a binary,
staggered treatment, the not-yet-treated firms would furnish a treatment-independent
control group from which the time-evolving partition could be estimated at each date
\citep[in the spirit of the control-group construction of][]{callawaysantanna2021},
yielding a map that is simultaneously evolving and exogenous. This is a genuine
optimisation of elegance over the frozen design, but it is purchased by discarding the
exposure-intensity information that motivates our continuous specification, and it
relies on a credible never-/not-yet-treated set under the diffuse Teti tariff
schedule. We therefore retain the continuous treatment with the frozen map as our
primary design, and present the binary control-group-map construction as a deliberate
robustness fork rather than the main specification.

\paragraph{Two concessions we make at the outset.}
Neither threat is specific to the network construction, and we address both head-on
rather than conceal them. First, the shock induces firm exit \citep{crozet2021}; since
an exited firm has no $d_{\rm in}/d_{\rm out}$, survivors are selected, and we therefore
estimate the effect on the probability of exit, bound the survivor effect, and read our
portfolio estimates as conservative. Second, a tariff on one product reallocates demand
to substitutes held by other firms in the same block, so the stable-unit-treatment
assumption is strained; the \citet{adao2019} correction absorbs the resulting
cross-firm correlation in the standard errors but not the causal interference, which we
discuss and, where feasible, probe with leave-out-block exposure.

\paragraph{Firm-level data and the ownership moderator.}
The panel of Chinese exporting firms comes from the NBS Annual Survey of Industrial
Enterprises, used in the same form by \citet{guariglia2011}. That paper establishes
three features of the NBS data that shape our design: SOEs operate under soft budget
constraints and face negligible financial frictions, private domestic firms are
severely financially constrained (near-unit-root self-financing), and coastal firms
without foreign capital face the binding constraint. Because firms switch ownership
status over the panel, we assign type by a consistent across-years rule following
\citet{guariglia2011}. The ownership heterogeneity motivates our SOE--private split:
privately owned firms facing financial constraints respond more sharply to export market
closures because they cannot subsidise loss-making product lines from retained earnings,
making the $d_{\rm out}$ contraction more immediate and the $d_{\rm in}$ consolidation
more forced.

% ---------------------------------------------------------------------------
\subsection{Our Contribution}
\label{subsec:contribution}
% ---------------------------------------------------------------------------

This paper contributes to all five literatures described above. To the economic-complexity
literature, we provide the first causal identification of a firm-level capability
response---$d_{\rm in}$, $d_{\rm out}$, and $\text{Share\_core}$---to an exogenous external shock,
demonstrating that the block-modular firm--product partition is not merely a descriptive
device but a meaningful predictor of differential adjustment paths. To the capability-theory
and evolutionary-economics literatures, we offer empirical content for the
exploitation--exploration tension formalised by \citet{march1991}: the direction of
adjustment under tariff shocks is shown to be systematically moderated by regional
strategic dependence, consistent with the Penrose--Teece prediction that capability
redeployment is constrained by the jointness of underlying knowledge inputs, while the
magnitude is consistent with the permanent-effect predictions of \citet{storper1992}
and the lock-in mechanism documented in the Schumpeterian Mark~II literature.

To the trade literature, we shift the outcome from performance (the productivity,
markup, and employment measures that dominate the Pavcnik--Topalova--Amiti generation) to
capability architecture, providing evidence that the same tariff shocks that raise
aggregate TFP through competitive selection simultaneously compress the capability
frontier of exposed firms in ways that the performance lens does not reveal. To
the strategic-dependency literature, we take the Farrell--Newman chokepoint
hypothesis to firm-level data: firms operating in CDI-dependent product
categories display systematically larger capability consolidation effects than firms
in non-dependent categories, and this heterogeneity is not explained by tariff
magnitudes alone, consistent with the dual treatment of higher tariffs and active
EU reshoring policy targeting the same product space. Finally, to the identification
literature, we demonstrate that a continuous-dose distributed-lag event study---with the
heterogeneity-robust continuous-treatment estimators of \citet{dechaisemartin2020} and
\citet{callawaygoodmanbacon2024} as robustness---is tractable at the scale of
full-universe firm--product panels, that the \citet{adao2019} shift-share-robust
correction materially affects standard errors in heterogeneous-exposure settings with
product-level clustering, and that reading a network-derived outcome against a
\emph{predetermined} (frozen, or pre-shock-estimated) capability partition is what keeps
the estimand exogenous to the shock.

% =============================================================================
\section{Expected Dynamics: What the Literature Predicts We Should Observe}
\label{sec:expected}
% =============================================================================

Because the empirical estimates are not yet in hand, it is worth making the priors
explicit. This section reads the literatures of Section~\ref{sec:lit} \emph{forward}:
for each outcome and each moderator, it states what dynamic the prior work leads us to
expect once the planned design is executed, what sign each hypothesis carries, and
what magnitude range is plausible. The aim is to specify in advance what a result
``consistent with the literature'' would look like---so that a surprising estimate can
be recognised as surprising---and to organise the diverse predictions into a single
testable structure. Nothing here is a finding; everything here is a prediction that
the data will confirm, refute, or qualify.

% ---------------------------------------------------------------------------
\subsection{The Out-of-Block Margin ($d_{\rm out}$)}
\label{subsec:exp_dout}
% ---------------------------------------------------------------------------

Three literatures converge on a contraction of $d_{\rm out}$ following an adverse
tariff shock. The sunk-cost extensive-margin mechanism of \citet{crozet2021} predicts
that exposed firms stop serving the products and markets that become tariff-burdened,
and---because re-entry costs are 19--23 percent higher for lapsed exporters---the
contraction is \emph{hysteretic}: it should not reverse within the event window even
if tariffs are later rolled back. The variance-cost mechanism of \citet{crozet2025}
and \citet{vannoorenberghe2016} predicts that firms which try to \emph{preserve}
$d_{\rm out}$ by pivoting to unfamiliar, structurally dissimilar destinations inherit
higher demand volatility, so the privately optimal response for all but the largest
firms is to let $d_{\rm out}$ fall rather than chase replacement markets. And the
diversification-discount search models of \citet{matsusaka2001} and \citet{bernardo2002}
predict that only firms whose core match remains above the liquidation threshold retain
the option to explore; below it, exploratory $d_{\rm out}$ search collapses. The net
expectation is a \emph{negative} average effect on $d_{\rm out}$, concentrated among
smaller, financially constrained, and inexperienced firms, and persistent over the
post-shock horizon. The one literature that cuts the other way is the escape-competition
margin of \citet{aghion2005} and \citet{akcigit2022}: a high-capability firm facing
competitive pressure may innovate \emph{outward}, lifting $d_{\rm out}$. Because that
channel is confined to frontier firms, the expectation is a negative average masking a
positive tail---visible only in the interaction with firm capability rank, not in the
mean.

% ---------------------------------------------------------------------------
\subsection{The In-Block Margin ($d_{\rm in}$) and Retrenchment}
\label{subsec:exp_din}
% ---------------------------------------------------------------------------

The prediction for $d_{\rm in}$ is the cleaner of the two. The Typical-Product-Vector
consolidation of \citet{fontagne2018}---a 7.7 percentage-point retention premium for
core products in treated sectors---is, mechanically, a $d_{\rm in}$ response to a
trade-policy shock, and it is the single closest empirical antecedent for our main
outcome. The scope-economy variance argument of \citet{bottazzi2006} predicts that a
firm rebalancing a disrupted portfolio minimises revenue covariance by deepening within
a related cluster rather than scattering, which raises $d_{\rm in}$ relative to
$d_{\rm out}$. The competition-to-coherence evidence of \citet{valvano2003}---firms in
single-market-exposed sectors became more coherent over time---predicts the same sign.
And the Schumpeterian Mark~II lock-in of \citet{dosi1995} predicts that the retrenchment
is strongest exactly where capabilities are most cumulative. The expectation is
therefore a \emph{positive} average effect on $d_{\rm in}$, larger for firms in
complex, cumulative technological regimes, and---per \citet{storper1992} and
\citet{riverabatiz1991}---persistent rather than transitory, because what is lost when
a market closes is more than demand: it is the collective learning infrastructure that
sustained the broader basket.

% ---------------------------------------------------------------------------
\subsection{Share\_core and the Direction of Adjustment}
\label{subsec:exp_share}
% ---------------------------------------------------------------------------

The summary statistic $\text{Share\_core} = d_{\rm in}/(d_{\rm in}+d_{\rm out})$ signs
the exploitation--exploration choice of \citet{march1991}. If the $d_{\rm out}$
contraction and the $d_{\rm in}$ consolidation predicted above both materialise,
$\text{Share\_core}$ rises---the retrenchment hypothesis (H1). March's own result that
organisations systematically over-exploit makes this the modal expectation. But the
literature is deliberately not unanimous: the dynamic-capabilities reading
(\citealp{penrose1959}; \citealp{teece1982}) and the moderate-shock search prediction
of \citet{matsusaka2001} leave open the exploration case (H1$'$), in which $\text{Share\_core}$
\emph{falls} because firms with slack treat the shock as a prompt to reallocate toward
new capability bundles. The empirical question is therefore which force dominates on
average, and, more informatively, where in the firm distribution the sign flips. The
strong prior, following \citet{destefano2025}, is that the average response is toward
the growth-\emph{negative} core (H1), because the growth-positive $d_{\rm out}$ margin
is precisely the one that sunk costs and variance costs make hardest to defend under
stress.

% ---------------------------------------------------------------------------
\subsection{Heterogeneity and the Strategic-Dependence Moderator}
\label{subsec:exp_het}
% ---------------------------------------------------------------------------

The sharpest predictions are conditional, not average. The Farrell--Newman chokepoint
logic \citep{farrell2019}, operationalised through the CDI product set of
\citet{arjona2023} and the network-fragility tiers of \citet{korniyenko2017}, predicts
that firms whose baskets concentrate in strategically dependent products face not a
neutral tariff but a tariff \emph{plus} active reshoring policy aimed at the same
products---so their retrenchment should be amplified and, given the escalation dynamic
of \citet{evenett2024}, more persistent. The financial-constraint heterogeneity of
\citet{guariglia2011} predicts that private Chinese firms retrench more sharply than
soft-budget SOEs. The size-and-volatility composition effect of
\citet{vannoorenberghe2016} predicts that small firms cannot afford to defend
$d_{\rm out}$ at all. The institutional-contingency result of \citet{delios2008} and
\citet{boschma2015} predicts that the related-versus-unrelated payoff itself shifts
with the regional institutional environment, so that ExpDep should be read as a
production-system attribute and not only a firm attribute. The blockade-logic taxonomy
of \citet{fontana2024} predicts that cumulative-strategic sectors retrench harder than
infrastructure-strategic ones. The unifying expectation is that \emph{heterogeneity
dominates the average}: a modest mean effect concealing a steep gradient in strategic
dependence, financial constraint, and capability rank.

% ---------------------------------------------------------------------------
\subsection{Synthesis: Hypotheses, Signs, and Magnitude Benchmarks}
\label{subsec:exp_synth}
% ---------------------------------------------------------------------------

Table~\ref{tab:hyp} collects the hypotheses, their literature sources, and the sign
each predicts. Table~\ref{tab:mag} brackets the plausible magnitude range using two
external benchmarks: a lower bound from pure network-disruption transmission
(\citealp{korniyenko2017}: a 1pp rise in risky imports from a disaster-struck origin
cuts the importer's own exports by $\sim$0.7 percent, with no policy response) and an upper
bound from the most exposed economies in a targeted tariff war (\citealp{teti2025}:
$-1.6$ percent cumulative real income for Canada and Mexico under the 2018--2025 US tariffs,
versus $-0.3$ percent for China in the aggregate---a country-level benchmark we map onto the
most exposed Chinese firms).

\begin{table}[ht]
\centering
\caption{Hypotheses, literature sources, and predicted signs.}
\label{tab:hyp}
\footnotesize
\setlength{\tabcolsep}{4pt}
\renewcommand{\arraystretch}{1.3}
\begin{tabular}{|c|p{4.4cm}|p{4.2cm}|p{2.5cm}|}
\hline
\textbf{Hyp.} & \textbf{Statement} & \textbf{Principal literature source} & \textbf{Predicted effect} \\
\hline
H1 & Adverse tariff shock raises $\text{Share\_core}$ / $d_{\rm in}$ (retrenchment to the coherent core). &
\citet{fontagne2018}; \citet{bottazzi2006}; \citet{valvano2003}; \citet{dosi1995}; \citet{march1991} &
$d_{\rm in}\uparrow$, $\text{Share\_core}\uparrow$ \\
\hline
H1$'$ & Adverse shock raises $d_{\rm out}$ / lowers $\text{Share\_core}$ (exploration beyond the core). &
\citet{penrose1959}; \citet{teece1982}; \citet{matsusaka2001}; \citet{aghion2005} &
$d_{\rm out}\uparrow$ (frontier-firm tail only) \\
\hline
H2 & The retrenchment response is amplified in more strategically dependent production systems. &
\citet{farrell2019}; \citet{arjona2023}; \citet{korniyenko2017}; \citet{evenett2024} &
$\beta_{2}>0$ on Tariff $\times$ ExpDep \\
\hline
H3 & The $d_{\rm out}$ contraction is hysteretic (does not reverse within the window). &
\citet{crozet2021}; \citet{storper1992}; \citet{riverabatiz1991} &
persistent, asymmetric \\
\hline
H4 & Responses are steeper for financially constrained / private / small firms. &
\citet{guariglia2011}; \citet{vannoorenberghe2016}; \citet{coad2021} &
steeper gradient \\
\hline
\end{tabular}
\end{table}

\begin{table}[ht]
\centering
\caption{Magnitude benchmarks for the $d_{\rm out}$ response (per s.d.\ of TariffShock).}
\label{tab:mag}
\footnotesize
\setlength{\tabcolsep}{4pt}
\renewcommand{\arraystretch}{1.3}
\begin{tabular}{|p{3.8cm}|c|p{6.2cm}|}
\hline
\textbf{Benchmark} & \textbf{Magnitude} & \textbf{Interpretation} \\
\hline
Lower bound --- pure network disruption \citep{korniyenko2017} & $\sim 0.7\%$ &
Supply-chain transmission of a localised shock, no policy intervention. \\
\hline
Upper bound --- most exposed economies, targeted tariff war \citep{teti2025} & $\sim 1.6\%$ &
Country-level loss for Canada/Mexico under the 2018--2025 US tariffs (vs.\ $-0.3\%$ for China), mapped onto the most exposed firms. \\
\hline
\end{tabular}
\end{table}

% =============================================================================
\section{Open Questions and a Forward Research Agenda}
\label{sec:agenda}
% =============================================================================

The review above leaves a small number of questions genuinely open and opens a larger
number of new ones. This section separates the two: first, the specific gaps in the
existing literature that \emph{this} paper is designed to close; second, the spin-off
questions that the same body of work raises but that lie beyond the present design and
would each support a distinct paper.

% ---------------------------------------------------------------------------
\subsection{The Gaps This Paper Fills}
\label{subsec:gaps_filled}
% ---------------------------------------------------------------------------

\begin{enumerate}
\item \textbf{Causal capability reconfiguration, not descriptive structure.} The EFC
literature (\citealp{bruno2023}; \citealp{pugliese2019}; \citealp{destefano2025})
measures block-modular capability structure but almost always descriptively. No prior
work links a \emph{predetermined, exogenous} trade-policy shock to the in/out
decomposition of firm diversification. This paper supplies that causal link.

\item \textbf{Composition, not performance.} The trade-and-productivity literature
(\citealp{pavcnik2002}; \citealp{topalova2011}; \citealp{amiti2007}; \citealp{tomasi2022})
measures how tariff shocks move productivity, markups, and survival. Because
performance dimensions decouple from capability (\citealp{grazzi2012};
\citealp{coad2021}; \citealp{syverson2011}), these outcomes are silent on \emph{which}
activities firms keep, drop, or enter. This paper moves the outcome to the architecture
of the export basket.

\item \textbf{Firm-level test of the chokepoint hypothesis.} The strategic-dependency
literature (\citealp{farrell2019}; \citealp{euc2021}; \citealp{arjona2023};
\citealp{korniyenko2017}) maps dependencies at the product and country level but has no
firm-level causal test. By interacting the tariff shock with a firm-specific
strategic-dependence exposure, this paper provides that test.

\item \textbf{A signed reading of exploration vs.\ exploitation under a real shock.}
March's (\citeyear{march1991}) tension has rich theory but little exogenous-shock
evidence at the firm level. The $\text{Share\_core}$ outcome under a predetermined tariff shock
delivers a directly signed estimate of which way firms move when their core is
threatened.

\item \textbf{Scalability of continuous-treatment event-study estimators.} We demonstrate
that a continuous-dose distributed-lag design with \citet{adao2019} shift-share inference---and
the heterogeneity-robust continuous-treatment estimators of \citet{dechaisemartin2020} and
\citet{callawaygoodmanbacon2024}---is tractable on a full-universe firm--product panel, a
methodological contribution in its own right for the heterogeneous-exposure shift-share
literature.
\end{enumerate}

% ---------------------------------------------------------------------------
\subsection{Beyond This Paper: Spin-Off Research Questions}
\label{subsec:spinoffs}
% ---------------------------------------------------------------------------

Each of the following is motivated by the reviewed literature but deliberately left
outside the present design; each would constitute a separate study.

\paragraph{(a) The firm--product--destination tensor.}
\citet{destefano2025} works in pure firm--product space; \citet{crozet2025} and
\citet{vannoorenberghe2016} show that the destination dimension carries its own
volatility structure. A natural extension promotes the bipartite firm--product matrix
to a firm--product--destination tensor and asks whether capability blocks reorganise
\emph{geographically}---whether retrenchment to the core is also a retreat to familiar
markets. This is the volatility-minimising hypothesis of $\Phi$-diversity made
firm-specific.

\paragraph{(b) Contemporaneous temporal-coupled partitions as a causal object.}
Our headline design freezes the baseline block, and our robustness lets it evolve on
\emph{pre-shock} information through the temporal multilayer machinery of \citet{mucha2010}
(with coupling and resolution chosen as in \citealp{bazzi2016}). What we deliberately leave
open is the \emph{contemporaneous} version: treating \emph{block reconfiguration itself}---the
birth, death, and merger of capability clusters measured with post-shock data---as the causal
outcome. Under our continuous treatment this is not identified, because there is no
never-treated control group from which to anchor an evolving map; a dedicated study could
recover it under a binary, staggered design by estimating the partition from not-yet-treated
firms, asking how trade shocks reshape the modular geometry of an entire industry rather than
the position of a single firm within a fixed geometry.

\paragraph{(c) The antidumping and safeguard stage.}
The Teti GTD \citep{teti2024} captures statutory MFN and preferential tariffs but
\emph{not} antidumping, safeguard, or trade-war tariffs. Pairing it with a temporary-trade-barriers
database would isolate the capability response to discretionary, firm-targeting
protection---plausibly a sharper and more anticipatory shock than scheduled tariff
changes, and a cleaner test of the escalation dynamic of \citet{evenett2024}.

\paragraph{(d) Instrumenting the policy-escalation component.}
\citet{evenett2024} document that industrial-policy intensity rises in election years
and that subsidies are matched tit-for-tat with 73.8 percent probability. An election-cycle
or tit-for-tat instrument for the policy-driven component of a firm's tariff shock
would separate the escalation channel from the mechanical tariff change---turning the
escalation caveat of this paper into an identification strategy of its own.

\paragraph{(e) Green-capability reconfiguration.}
\citet{mealy2022} show the complexity toolkit applies to any product subset and that
green-capability accumulation is strongly path dependent. Re-running the present design
with the green-product set as the treated category would ask whether decarbonisation
policy and trade policy push firms in the same or opposite capability directions---a
direct bridge between the strategic-autonomy agenda of \citet{cagnin2021} and the green
transition.

\paragraph{(f) Institutional voids as a moderator.}
\citet{delios2008} and \citet{boschma2015} show that the value of diversification, and
the related-versus-unrelated balance, depend on institutional quality. Exploiting
within-China regional institutional variation as a second moderator would test whether
capability retrenchment under trade shocks is steeper or shallower where market
institutions are weaker, distinguishing capability exploitation from
institutional substitution as motives for diversification.

\paragraph{(g) The gain/loss asymmetry of complexity.}
\citet{nomaler2023} show that lost specialisations are often \emph{more} complex than
gained ones, and that density-based relatedness is heavily confounded by diversity and
ubiquity. A study that scores firm-level product \emph{exits} by complexity---asking
whether tariff shocks make firms shed their most sophisticated products first---would
test whether trade shocks erode the capability frontier from the top, a sharper and
more policy-relevant question than net diversification counts can answer.
 into the main paper.
% Cite keys follow the convention author(s)year; matching @article / @book /
% @techreport entries live in Bibliography_base.bib (86 entries as of 2026-06-05).
%
% STRUCTURE
%   2   Related Literature
%     2.1  Economic complexity, capability structure, and diversification theory
%     2.2  Capability theory, corporate coherence, and the direction of adjustment
%     2.3  Trade policy shocks and firm adjustment
%     2.4  Strategic dependencies and geopolitical trade shocks
%     2.5  Identification
%     2.6  Contribution
%   3   Expected Dynamics: What the Literature Predicts We Should Observe
%     3.1  The out-of-block margin (d_out)
%     3.2  The in-block margin (d_in) and retrenchment
%     3.3  Share_core and the direction of adjustment
%     3.4  Heterogeneity and the strategic-dependence moderator
%     3.5  Synthesis: hypotheses, signs, and magnitude benchmarks  [2 tables]
%   4   Open Questions and a Forward Research Agenda
%     4.1  The gaps this paper fills
%     4.2  Beyond this paper: spin-off research questions
%
% NOTE: tables use plain \hline tabular (no booktabs dependency) so the file
% compiles stand-alone via lit_review_test.tex and when \input into the paper.
% =============================================================================

\section{Related Literature}
\label{sec:lit}

This paper sits at the intersection of four bodies of work: the economic-complexity
literature on firm-level capability structure and the theory of diversification; the
capability-based theory of the firm and evolutionary economics; the empirical trade
literature on the effects of tariff shocks on firm outcomes; and the emerging literature
on strategic dependencies, geopolitical risk, and the weaponisation of economic
interdependence. A fifth, methodological strand covers identification in shift-share
designs with heterogeneous treatment effects. The unifying thread across these
literatures is a single dynamic: an external shock that degrades the value of a firm's
existing activities forces a \emph{reconfiguration} of its productive knowledge, and
the literatures disagree on whether that reconfiguration deepens a coherent core or
pushes outward into new capability bundles. We discuss each stream in turn, emphasising
what each predicts about the \emph{direction} and \emph{persistence} of adjustment, and
close with a statement of our contribution.

% ---------------------------------------------------------------------------
\subsection{Economic Complexity, Capability Structure, and Diversification Theory}
\label{subsec:lit_efc}
% ---------------------------------------------------------------------------

The product-space framework of \citet{hausmann2007} and \citet{hidalgo2009}
established that countries export products in patterns that are far from random:
goods that require similar underlying capabilities tend to be co-exported, generating
a proximity structure in the product space that conditions the path of structural
transformation. \citet{hidalgo2011} formalise this intuition through a bipartite
capabilities model in which countries and products are represented as the two node
sets of a bipartite graph, and the adjacency matrix $M_{cp}=1$ whenever country $c$
has revealed comparative advantage (RCA) in product $p$. The key result is that
diversification is convex in the capability stock: countries with few capabilities
earn near-zero marginal returns from acquiring any single new one, generating the
low-income quiescence trap that has been documented empirically across the
income distribution \citep{hausmann2007}. This non-linearity is fundamental to the
agenda of using complexity tools in a causal setting: treatment-effect heterogeneity
by capability endowment is an architectural feature of the framework, not an
add-on.

The Fitness-Complexity (FC) algorithm of \citet{cristelli2013} provides the
operational complexity metric adopted in the subsequent firm-level literature. Rather
than the linear iterative procedure of \citet{hidalgo2009}, FC defines country fitness
as the sum of product complexities and product complexity as the harmonic mean of
exporter fitness---a bounded operator that uniquely rewards diversification and
assigns complexity to a product according to the least capable exporter that ships it,
giving the algorithm its asymmetric censoring property. The algorithm converges to a
unique fixed point and, as shown by \citet{cristelli2017}, generates a selective
predictability scheme in which high-fitness countries follow laminar, predictable
growth trajectories while low-fitness countries exhibit chaotic dynamics---a property
that has direct implications for heterogeneous treatment responses in our setting:
firms embedded in high-complexity capability clusters should display more systematic
responses to tariff shocks than firms operating at the periphery of the product
space. \citet{angelini2017} characterise the long-run dynamics of products in the complexity
plane, showing that they drift toward a stationary configuration so that
shock-induced dislocations are transient and mean-reverting---which is why we treat
product complexity as a well-defined pre-shock control and read post-shock movements
against a fixed baseline.
\citet{stojkoski2025} extend the framework to an optimisation context,
showing that the ECI-growth relationship is stronger at lower income levels and that
relatedness is one of several predictors of product-entry, a finding that validates
the multi-dimensional block structure we exploit over simpler relatedness-only
measures. \citet{koch2021} sounds a useful caution about what such metrics measure: the
complexity--growth link is in fact \emph{stronger} when complexity is computed on
value-added rather than gross exports ($R^2 = 0.63$ versus $0.31$), so gross-flow
complexity blurs the connection to realised value capture. Since our firm-customs
measures are necessarily built on gross flows, we read them as capturing capability
\emph{structure} rather than realised value capture, which is precisely why we treat
the block partition as a measure of capability architecture rather than of performance.

Three strands sharpen what ``relatedness'' and ``diversification'' actually mean for
the direction of adjustment. First, the related-variety tradition of
\citet{frenken2007} distinguishes \emph{related} variety (entry into activities that
share capabilities, which drives regional \emph{employment} growth through
within-sector Jacobs spillovers---though, tellingly, not productivity growth) from
\emph{unrelated} variety (which provides a portfolio hedge that lowers unemployment
under asymmetric sectoral shocks). This is the regional-economics antecedent of our
$d_{\rm in}$ versus $d_{\rm out}$ distinction. \citet{boschma2015} show that the relative payoff to
related versus unrelated diversification is institutionally contingent: coordinated
market economies favour incremental related diversification, while liberal market
economies support the unrelated, exploratory kind, a result that motivates treating
the regional production system, not just the firm, as a unit of heterogeneity.
\citet{freire2019} provides a structural model of endogenous diversification in which
capability accumulation and product entry co-evolve, giving microfoundations to the
empirical relatedness regularities. Second, \citet{mealy2022} demonstrate that the
complexity toolkit transfers to any policy-relevant subset of products---their
green-complexity indices show strong path dependence in capability accumulation
(a $0.92$ correlation between green-complexity potential and the green-complexity
index, with beginning-of-period potential predicting subsequent green-export
growth), a template we adopt when restricting attention to strategically dependent
product categories. Third, \citet{nomaler2023} issue a confounding warning that
directly disciplines our outcome construction: density-based relatedness metrics are
$\sim$93 percent explained by diversity and ubiquity alone, so a measured ``relatedness''
effect can be an artefact of how many products a firm makes and how common those
products are. Their symmetric treatment of RCA \emph{gains} and \emph{losses}, together with
the finding that lost specialisations are often \emph{more} complex than gained
ones, reframes diversification as a two-sided entry-and-exit process and is one of the
closest conceptual antecedents to our reconfiguration question.

The application of complexity tools to firm-level data is recent. \citet{bruno2018, bruno2023} construct firm-product bipartite
networks from Customs microdata and show that RCA-based binarisation
(with BiCM entropy-null correction as a robustness) yields block-modular structures
that are stable across years and meaningfully predict firm growth, innovation, and
survival. \citet{pugliese2019} show that firms diversify coherently within their
existing technological portfolio, so that a firm's position within its capability
cluster is informative about its future mobility---a finding that motivates our
within-block heterogeneity analysis; the country-level analogue is \citet{pugliese2017},
who show that complex economies escape the low-income trap by moving laterally into
related capabilities. \citet{coad2021} push
this furthest, applying the FC algorithm to a firm$\times$capability matrix for $\sim$40{,}000
Indian firms and recovering a nested ``capabilities ladder'' whose rungs predict growth
and survival but---tellingly---not profitability, a decoupling we return to in
Section~\ref{subsec:lit_cap}. \citet{destefano2025}
provide the most direct methodological antecedent for our outcome measures: they
decompose firm diversification into in-block ($d_{\rm in}$) and out-of-block
($d_{\rm out}$) counts using the block partition produced by the BRIM algorithm of
\citet{barber2007}, define $\text{Share\_core} = d_{\rm in}/(d_{\rm in}+d_{\rm out})$,
and show that out-of-block diversification is the growth-positive margin while
retrenchment toward the core---the $d_{\rm in}$ direction---is the
growth-negative, lock-in-reinforcing margin. We adopt their exact outcome measures
and employ the same partition algorithm.

The network-method foundations underpinning the partition are well established.
\citet{newman2006} introduced the modularity objective and its spectral optimisation;
\citet{blondel2008} provide the Louvain heuristic that makes community detection
tractable at the scale of full-universe firm-product matrices; \citet{ducharenas2005}
supply the extremal-optimisation alternative we use as a partition robustness check;
and \citet{fortunato2016} provide the authoritative critical survey of the field's
validation pitfalls. \citet{barber2007} adapts modularity to the bipartite case (BRIM),
and \citet{soleribalta2018}
show that the BRIM objective is insensitive to within-block
nested structure and can therefore miss the hierarchical ranking of firms inside a
capability block; their in-block nestedness quality function $I$ jointly optimises the
block partition and within-block NODF \citep[the consistent nestedness metric of][]{almeidaneto2008},
providing the basis for our within-block rank ordering of firms. \citet{mariani2019} prove that the FC algorithm is a nestedness
maximiser and clarify the relationship between modularity and nestedness in firm-level
networks: the two are positively correlated at low network densities (the regime typical
of firm-product matrices) and can be jointly estimated. Validation of the BRIM
partition against entropy null models---using the Bipartite Configuration Model of
\citet{squartini2011} and \citet{saracco2015}, with the projection reconciliation of
\citet{cimini2022}---is a required pre-processing step
to ensure that block structure reflects genuine capability clustering rather than
degree-sequence artefacts, and we implement it throughout. A wider set of Economic Fitness and
Complexity (EFC) studies---broader applications of economic-complexity methods---rounds
out the methodological base: global-value-chain specialisation \citep{bontadini2021, bontadini2025},
the complexity of structural change and the sustainability transition
\citep{caldarola2024a, caldarola2024b}, machine-learning reconstruction of trade flows
\citep{gnecco2023}, and relatedness-based forecasting of industrial upgrading
\citep{albora2023a, albora2023b}.

% ---------------------------------------------------------------------------
\subsection{Capability Theory, Corporate Coherence, and the Direction of Adjustment}
\label{subsec:lit_cap}
% ---------------------------------------------------------------------------

The theoretical framework connecting firm capability stocks to the direction of
diversification draws from the resource-based view of the firm, evolutionary economics,
and the learning-and-adaptation literature. \citet{penrose1959} establishes that
firms are bundles of productive resources whose services can be redeployed across
activities at varying cost: the excess services that routinely arise from indivisible
resources create an organic pressure toward related diversification that is internal
to the firm, not driven by market structure. \citet{teece1982} formalises this in a
scope-economy model showing that the boundaries of the multi-product firm are
determined by the jointness of underlying knowledge inputs; firms expand into activities
that share tacit knowledge with existing lines, generating the coherence property
that the EFC literature quantifies through block-modular partitions.
\citet{teece1994} make this measurable through the survivor-principle relatedness
metric---sectors are ``related'' when more firms combine them than chance predicts---which
is the product-space analogue of the bipartite co-occurrence logic at the heart of the
EFC programme. The corporate-coherence literature that builds on it is directly
informative about the \emph{dynamics} we expect to observe. \citet{piscitello2000}
finds that large firms reduced product coherence while \emph{maintaining}
technological coherence over 1977--1995: capabilities are the stable organising
principle even as product portfolios churn, which is exactly why a capability-block
partition should track adjustment more faithfully than a raw product count.
\citet{valvano2003} provide the most pointed antecedent for a trade-shock prediction:
in Italian manufacturing, firms in sectors more exposed to single-market competition
(``sensitive'' sectors) became \emph{more} coherent over time, and entrants into the
leadership ranks were more coherent than incumbents---direct evidence that intensified
competition pushes surviving firms toward their core ($d_{\rm in}$). \citet{delios2008}
add the institutional contingency: within-country product diversification is more
valuable where market institutions are weak, because diversification substitutes for
missing capital, labour, and information markets---a moderation channel that warns
against reading every diversification response as pure capability exploitation.

The corporate-finance literature on the ``diversification discount'' supplies the
causal-direction caution that disciplines our interpretation. \citet{langstulz1994}
document a robust industry-adjusted discount in Tobin's $q$ of 0.35--0.50 for
diversified US firms throughout 1978--1990, but show that firms diversify \emph{after}
their core industry weakens---the discount precedes diversification rather than being
caused by it. \citet{matsusaka2001} rationalises this as capability-match search: a
firm whose current match deteriorates maintains its existing lines while experimenting
with new ones (the ``$b$-strategy''), liquidating and restarting only when the match
falls below a lower threshold. \citet{bernardo2002} recast the same fact as
real-option value: single-segment firms are worth more because they retain unexercised
expansion options. Translated to our setting, these models deliver a sharp,
non-obvious prediction: a \emph{moderate} tariff shock should push firms toward
exploratory $d_{\rm out}$ search (the $b$-strategy), whereas a \emph{severe} shock that
drives match quality below the liquidation threshold collapses the option value of
organisational capabilities and forces retrenchment or exit---so the sign of the
average response is genuinely ambiguous and should depend on shock intensity, exactly
the non-linearity our event study is designed to detect.

The evolutionary perspective of \citet{dosi1988, dosi1995} adds a regime dimension:
the rate and direction of capability accumulation depend on both the firm's own
endowment and the technological regime of the sector---the opportunity conditions,
appropriability, and cumulativeness that structure search. Under Schumpeterian Mark
II conditions (cumulative, large-firm-dominated regimes), incumbents sustain
technological advantages through internal learning, generating the technological
lock-in that maps directly onto high $\text{Share\_core}$. The implication for our
setting is that firms in complex, cumulative technological regimes should display
stronger $d_{\rm in}$ responses to adverse shocks---retrenchment to the core is both
the natural adjustment and the organisationally feasible one when search capabilities
are already embedded in existing product clusters.

\citet{march1991} provides the ambiguity at the heart of our hypotheses. The
exploitation--exploration tension implies that the direction of capability adjustment
under stress is not a priori signed: firms that face an adverse demand or competitive
shock in their existing core can either intensify exploitation of known competencies
(retrenchment) or accelerate exploration beyond the core (diversification-as-escape).
March's formal result is that organisations systematically over-exploit under most
plausible parameter configurations, predicting the retrenchment (H1) hypothesis as
the modal response but leaving open the exploration (H1$'$) case for firms with
sufficient slack and absorptive capacity. \citet{patel1997} document the empirical
content of this tension at the firm level, showing that large innovative firms spread
R\&D effort across fields but in highly persistent patterns---path dependence is the
rule, not the exception---while \citet{bell1993} establish that technological
capability accumulation in developing-country firms is even more path-dependent,
with explicit investments in learning required to escape existing trajectories.
\citet{storper1992} connects this to trade: product-based technological learning
generates agglomeration economies in particular product-territory combinations, so
that an adverse trade shock removes both demand and the collective learning
infrastructure that sustains capability accumulation---a channel that implies
persistent, not merely transitory, capability effects of tariff shocks.

The firm-growth literature, summarised in \citet{coad2009} and documented empirically
in \citet{bottazzi2003, bottazzi2005, bottazzi2006}, shows that firm-level growth is
Laplacian---fat-tailed, near-unit-root---and that the variance of growth rates scales
inversely with firm size at a rate explained entirely by diversification across product
lines. The scope-economy mechanism \citep{bottazzi2006} implies that firms which
diversify within a related cluster (high $d_{\rm in}$) reduce the co-movement of
their individual product revenue streams more effectively than firms that scatter across
unrelated products. Under tariff shock conditions, firms that lose access to existing
export markets face a portfolio rebalancing problem whose solution---from a pure
variance-minimisation standpoint---is to deepen within-block exposure rather than
scatter to unfamiliar destinations, reinforcing the retrenchment prediction. A
recurring caveat from this literature is that performance dimensions decouple:
\citet{grazzi2012} shows that Italian exporters carry a large productivity premium but
\emph{no} profitability premium, and \citet{coad2021} that capability-ladder position
predicts growth and survival but not profit margins; \citet{syverson2011} surveys the
enormous, persistent productivity dispersion that these patterns sit inside. The
collective lesson---that capability advantage need not show up in profits or even in
growth---is why we measure the \emph{architecture} of the export basket
($d_{\rm in}$, $d_{\rm out}$, $\text{Share\_core}$) rather than a performance scalar.
\citet{aghion2001, aghion2005} add the innovation margin: neck-and-neck competitors
facing increased competitive pressure innovate to escape competition, while laggards
reduce innovation effort. Translating to our setting, firms with high within-block
fitness ranks face incentives to deepen core capabilities under tariff pressure (the
escape-competition channel), while low-rank firms---unable to sustain the cost of
maintaining their current product portfolio---are more likely to prune $d_{\rm out}$
and consolidate around the most defensible core products. This generates the
within-block heterogeneity in treatment effects that we exploit empirically.

% ---------------------------------------------------------------------------
\subsection{Trade Policy Shocks and Firm Adjustment}
\label{subsec:lit_trade}
% ---------------------------------------------------------------------------

A large reduced-form literature establishes that tariff changes cause significant
adjustments in firm outcomes. \citet{pavcnik2002} and \citet{topalova2011} document
that trade liberalisation raises TFP through competitive selection and technology
imports in Chilean and Indian manufacturing, respectively. \citet{amiti2007} show,
for Indonesian firms, that input tariff reductions deliver TFP gains at least twice as
large as those from output tariff reductions---a finding driven by access to high-quality
imported intermediates---and construct the shift-share input-tariff exposure variable that is
the methodological template for our treatment. \citet{tomasi2022} replicate this
design for Vietnamese firms during the 2007 WTO accession, finding that governance
quality moderates the tariff-TFP effect, providing the closest precedent for our
moderated shift-share design with ExpDep as the heterogeneity variable.

What this literature consistently measures is performance: productivity, employment,
markups, survival. It is essentially silent on the \emph{composition} of a firm's
product portfolio---which activities are retained, dropped, or newly entered, and
whether the reconfiguration deepens a coherent core or spreads capabilities into
new domains. The EFC literature reviewed in Section~\ref{subsec:lit_efc} supplies
the tools to measure that composition, but has applied them almost exclusively in
descriptive, non-causal settings. This paper occupies the intersection left empty
by both strands. The theoretical channel linking trade exposure to capability
adjustment is supplied by the trade-and-innovation literature: \citet{akcigit2022}
unify the market-size, business-stealing, and escape-competition effects of trade on
innovation incentives, showing that the net effect on a firm's capability investment
is sign-ambiguous and depends on its competitive position---the firm-level analogue of
the exploitation--exploration tension. \citet{riverabatiz1991} establish the deeper
point that integration affects growth through flows of \emph{ideas} as well as goods,
so that a tariff shock that severs a firm from a technological frontier market can have
permanent (growth-rate), not merely level, effects on its capability trajectory---reinforcing the
persistence prediction from \citet{storper1992}.

The most direct evidence that trade shocks restructure capability portfolios rather
than merely rescaling performance comes from \citet{fontagne2018}, who show that
Chinese export quotas on textiles and clothing generated a significant \emph{Typical
Product Vector} consolidation among French exporting firms: products at the core of
a firm's stable portfolio were retained at a 7.7 percentage-point higher rate than
peripheral products in treated sectors. This is, in effect, a $d_{\rm in}$ response
to a trade-policy shock measured with product-vector tools---the closest antecedent
in the literature for our main outcome. \citet{fontagne2022} provide the
product-level trade elasticities that we use as Bartik weights: 5,050 HS6-digit
elasticities estimated via PPML gravity, averaging $-5.3$ (median $-3.5$), with
HS-section means spanning $-18.5$ (minerals), $-6.1$ (machinery), and $-3.6$
(footwear), capturing the heterogeneous disruption that a uniform tariff schedule
imposes on firms with different product portfolios.

The dynamic and firm-level dimensions of portfolio adjustment are examined by
\citet{crozet2021}, who study the effect of EU and US sanctions on the extensive
margin of French exporters to Iran and Russia. Using monthly customs data and a
dynamic probit model with three-way fixed effects and analytical bias correction,
they find that sanctions lower the probability of serving a sanctioned market by 39
percent (Iran) and 23 percent (Russia), with strong state dependence: the entry-cost
factor is 19--23 percent higher for firms not active in the previous year. Crucially,
the lifting of sanctions on Myanmar generates only a small and gradual positive response,
establishing an \emph{asymmetry} between imposition and removal that implies permanent,
not transient, $d_{\rm out}$ contraction even after the policy shock is reversed.
This hysteresis result---consistent with the Roberts--Tybout sunk-cost model---means
that treatment effects in our event study should be interpreted as reflecting a
permanent restructuring of the capability portfolio, not a temporary demand disruption.
The firm-level heterogeneity documented by \citet{crozet2021}---that trade-finance-dependent
firms, inexperienced firms, and non-specialist firms exit more sharply---maps
directly onto the moderation hypotheses we test: private Chinese firms (documented
by \citet{guariglia2011} as the most financially constrained in the NBS panel) should
display amplified $d_{\rm out}$ reductions, while SOEs with soft budget constraints
should display more muted responses.

A complementary strand examines the export portfolio diversification strategies that
firms deploy to manage demand volatility. \citet{vannoorenberghe2016} show, for
Chinese firm-level data, that destination diversification reduces volatility for large
exporters but \emph{increases} it for small ones, with the switching point determined
by a composition effect: occasional exports to marginal destinations dominate the
diversification gain for small firms, while stable core-market exports dominate for
large ones. \citet{crozet2025} formalise this puzzle through a phyloeconomic diversity
index $\Phi$ that adds a dissimilarity dimension to the standard richness-evenness
decomposition. Their central finding---that conditional on the number of destinations
and the Herfindahl concentration of export sales, higher $\Phi$ (more dissimilar
markets) is associated with \emph{higher} export sales volatility---reveals that
hedging through geographic diversification has limits: distant, structurally different
destination markets carry higher idiosyncratic demand variance, so that firms forced
to substitute familiar Western markets with unfamiliar alternatives do not achieve
the volatility reduction that naive diversification logic predicts. For our setting,
this implies that Chinese firms which maintain $d_{\rm out}$ by pivoting to non-Western
destinations face not only lower export levels but higher volatility---a double cost
that reinforces the variance-minimising case for $d_{\rm in}$ consolidation.
\citet{barbieri2024} provide a parallel decomposition at the territorial level:
\emph{sectoral} diversification (a proxy for productive capabilities, the conceptual
sibling of $d_{\rm in}$) raises export intensity, while \emph{geographic}
diversification (a proxy for geoeconomic capabilities, loosely analogous to the
resilience role we attribute to $d_{\rm out}$) is associated with greater export
stability; institutional quality independently improves both margins.\footnote{The
analogy is functional rather than literal. Barbieri's second axis is
\emph{destinations}, whereas our $d_{\rm out}$ remains a \emph{product}-space
measure (out-of-block products), following \citet{destefano2025}, who works in pure
firm--product space with no destination dimension. The roles align---one margin scales
output, the other buys resilience---but the underlying axes differ, and in their data
the dominant stabiliser is EU market exposure rather than geographic diversification
per se. We therefore treat the destination dimension as a second-stage extension
(Section~\ref{subsec:spinoffs}) rather than as a component of $d_{\rm out}$.}

We estimate treatment using the Global Tariff Database (GTD) of \citet{teti2024},
the first applied-tariff panel to cover bilateral HS6 tariffs continuously from 1988
to 2021 while correcting for the false interpolation and selection artefacts of UNCTAD
TRAINS. \citet{teti2019} establish, using structural gravity, that Rules-of-Origin
compliance is prohibitively costly for the machinery and electronics sectors at the
core of Chinese exports, confirming that MFN tariffs---the rate captured by the
GTD---are the economically relevant wedge for our firms. The welfare magnitudes from
\citet{teti2025} calibrate our expected effect sizes: a GE simulation of the
2018--2025 US--China tariff war estimates cumulative real-income losses of roughly
$-1.6$ percent for Canada and Mexico---economies with high export exposure to the US
market---against only $-0.3$ percent for China in the aggregate. We use the
$-1.6$ percent figure as an external benchmark for the most exposed Chinese firms:
by analogy with the highly exposed Canadian and Mexican case, the uneven distribution
of tariff exposure across product portfolios should push firms concentrated in
high-regulatory-dependence categories toward magnitudes well above the national
average. The mapping from these country-level losses to firm-level exposure is our
own; \citet{teti2025} simulate countries and US states, not firms. The distortionary
role of size-contingent and sector-specific regulation documented by
\citet{garicano2016}---who find welfare costs of 1.3--3.5 percent of GDP from French
size thresholds---is the domestic-policy analogue of this exposure heterogeneity and
motivates controlling for regulatory environment in the moderation analysis.

% ---------------------------------------------------------------------------
\subsection{Strategic Dependencies and Geopolitical Trade Shocks}
\label{subsec:lit_geo}
% ---------------------------------------------------------------------------

The design of the heterogeneity variable ExpDep draws from a separate and newer
literature on strategic economic dependencies. \citet{farrell2019} identify the
conditions under which the topology of global economic networks generates geopolitical
leverage: hub-and-spoke structures create chokepoint nodes through which flows can be
controlled (flow restriction) and from which information can be extracted (panopticon
effects). Crucially, weaponisation requires two joint conditions---jurisdictional
control over the chokepoint and domestic institutions capable of deploying that control
as a foreign policy instrument---which limits the set of products and bilateral
relationships that qualify as strategic in the relevant sense. The Farrell--Newman
framework grounds our moderation hypothesis: Chinese firms whose export baskets are
concentrated in products where Western jurisdictions have identified chokepoint control
face not a neutral tariff shock but one accompanied by active policy interest in
reducing dependence---generating an amplified and persistent treatment effect.

The quantitative operationalisation of this insight is provided by \citet{euc2021},
who develop the Core Dependency Indicators (CDI) methodology to identify EU strategic
dependencies at the HS6-product level. The methodology applies three thresholds:
supplier concentration (HHI of extra-EU import shares exceeding 0.4, implying fewer
than 2.5 effective suppliers), import reliance (extra-EU imports exceeding 50 percent
of total EU imports of the product), and substitutability (extra-EU import volume
exceeding total EU export volume, so that domestic production cannot cover the import
volume even in an emergency). Of approximately 5,000 HS6 products, 137 pass all three
thresholds in the EU's most sensitive ecosystems, with China accounting for roughly
52 percent of their import value. \citet{arjona2023} refine the methodology with a
dynamic four-year threshold and a re-export correction: accounting for transit trade
\emph{raises} the count of dependent products by roughly 27 percent, since ignoring
re-exports understates true extra-EU reliance, yielding a dependent-product set of
204 items that provides the product classification we use to construct ExpDep.
The forward-looking scenario work of \citet{cagnin2021} situates these indicators in
the EU's Open Strategic Autonomy agenda, projecting an escalation of dependency-reducing
policy through 2040 that makes the persistence of any treatment effect we estimate a
policy-relevant, not merely statistical, question.

The network-science analogue of the CDI methodology is provided by
\citet{korniyenko2017}, who combine three structural indicators of product trade networks---
the standard deviation of weighted outdegree centrality (capturing hub dominance), the
product of clustering coefficient and network diameter (capturing geographic
compartmentalisation), and a Revealed Factor Intensity measure of substitutability---to
classify 3,578 intermediate and capital goods into four risk tiers via k-median
clustering. Their persistently risky group comprises 421 products over 2003--2014,
concentrated in machinery and mechanical appliances (HS 84--85), transport equipment,
pharmaceuticals, and precision instruments. The causal validity of the classification
is demonstrated by two out-of-sample cases: the 2011 Japan earthquake and Thai floods,
where products subsequently in short supply were correctly identified as Group~4 one
year prior. The overlap between the Korniyenko Group~4 products and the
Arjona CDI-dependent set defines the most extreme ExpDep exposure: products
that are simultaneously structurally concentrated, topologically fragile, and actively
targeted by EU and US reshoring policy.

\citet{evenett2024} provide real-time evidence on the policy escalation dynamic. Their
New Industrial Policy Observatory (NIPO), covering 2,500$+$ new industrial policy
measures across 75 jurisdictions in 2023, documents a 73.8 percent probability that a
subsidy by one major economy (China, EU, or US) is matched by a subsidy in the same
product by another within 12 months. PPML regressions confirm that IP usage correlates
positively with RCA---governments target sectors where they already have comparative
advantage, consistent with political-economy capture rather than market-failure
correction. For our event study, this finding implies that the Teti GTD tariff
shock at time $t$ is not a clean one-off treatment but the first signal in an
escalating policy sequence: treatment intensity grows after the event-window, pushing
treatment effects toward permanence.

\citet{fontana2024} embed these observations in a political-economy theory of
strategic-assets policy, drawing on the three logics of capability blockade formalised
by Ding and Dafoe: cumulative-strategic blocking (targeting technological RCA through
standards exclusion and export controls), infrastructure-strategic blocking (targeting
critical material supply chains through reshoring subsidies), and dependency-strategic
blocking (exploiting existing dependencies through sanctions and FDI screening). The
taxonomy provides a structural basis for splitting the Teti tariff sample by motive:
tariff shocks in sectors matching the cumulative-strategic logic (advanced
manufacturing, semiconductors) should generate stronger $d_{\rm in}$ responses than
shocks in sectors matching the infrastructure-strategic logic (raw materials, chemicals),
where substitution is constrained by the CDI3 condition. We use this taxonomy in a
sample-split robustness exercise.

% ---------------------------------------------------------------------------
\subsection{Identification}
\label{subsec:lit_id}
% ---------------------------------------------------------------------------

Our identification strategy rests on a single organising principle: \emph{everything
that could mechanically respond to the shock must be predetermined relative to it.}
The estimand is the within-firm response of an export-capability portfolio to a
trade-policy shock, and the credibility of that estimand depends on three objects
being fixed before the shock---the exposure weights, the disruption elasticities, and
the capability map against which the portfolio is read. We discuss each, and then
state precisely which of three admissible network constructions we adopt and why.

\paragraph{A continuous, time-varying treatment.}
Following the input-tariff shift-share of \citet{amiti2007}, \citet{topalova2011}, and
\citet{tomasi2022}, we construct the tariff exposure variable as
\[
  \text{TariffShock}_{it}
  = \sum_p s_{ip,\text{pre}}\cdot \varepsilon_p \cdot \Delta\ln(1+\tau_{pt}),
\]
where $s_{ip,\text{pre}}$ are firm--product export shares in a strictly pre-shock
baseline window, $\Delta\ln(1+\tau_{pt})$ are changes in bilateral applied tariffs
from the Teti GTD, and $\varepsilon_p$ are the product-level trade elasticities of
\citet{fontagne2022}. We treat $\text{TariffShock}_{it}$ as a \emph{continuous,
time-varying, non-absorbing} dose rather than a binary staggered treatment: the dose
grows and recedes as tariffs move, and discretising it into a treated/control
indicator would discard precisely the cross-firm variation in exposure intensity that
identifies the dose--response. This choice is deliberate and consequential, as we make
clear below. The elasticities $\varepsilon_p$ are estimated by \citet{fontagne2022} on
the universe of 152 importers and 189 exporters---they are a property of the
\emph{product}, external to our firm sample, and hence cannot be a function of the
behaviour we study; products with a non-significant estimate inherit their HS-section
mean. The predetermined shares satisfy the Bartik exogeneity condition, which we defend
on \emph{both} of its routes: from the shares, following \citet{goldsmithpinkham2020},
and from the shocks, following \citet{borusyakhull2022}, reporting Rotemberg-weight
diagnostics and a recentred-instrument specification.

\paragraph{Dynamics and inference for a continuous dose.}
Because the treatment is continuous and non-absorbing, the cohort$\times$relative-time
machinery of \citet{sunabraham2021} is not the natural primary specification---its
cohorts are defined by the timing of a binary switch. We therefore estimate the
dynamics through a distributed-lag event study, $Y_{it}=\sum_k \beta_k\,
\text{TariffShock}_{i,t-k}+\alpha_i+\delta_t+\varepsilon_{it}$, in which the lead
coefficients ($k<0$) constitute the parallel-trends test and the lag coefficients
($k>0$) trace the persistence (hysteresis) prediction. Inference follows the
shift-share-robust standard errors of \citet{adao2019}, which are native to continuous
exposure designs and which we report as primary against firm- and sector-clustered
alternatives. As robustness we re-estimate the dynamics with the continuous,
non-absorbing estimator of \citet{dechaisemartin2020}, the continuous-treatment
difference-in-differences framework of \citet{callawaygoodmanbacon2024}, and---after
discretising exposure into terciles---the heterogeneity-robust estimators of
\citet{sunabraham2021}, \citet{callawaysantanna2021}, and \citet{borusyak2024}. We
flag the assumption that the continuous design carries a \emph{stronger} parallel-trends
requirement than the binary case---\citet{callawaygoodmanbacon2024} show that
identifying the dose--response slope requires the absence of selection into dose level,
which we probe but cannot fully test.

\paragraph{The capability map: three admissible constructions, and the one we adopt.}
The outcomes $d_{\rm in}$, $d_{\rm out}$, and $\text{Share\_core}$ are read off a
block partition of the firm--product export network \citep[the export, not production,
basis of][]{destefano2025}, obtained by bipartite modularity maximisation
\citep{barber2007} on an RCA-binarised incidence matrix validated against the Bipartite
Configuration Model \citep{squartini2011,saracco2015,cimini2022}. The non-trivial
question is \emph{when} the partition is estimated, because the map is itself the
measuring instrument for the outcome, and a map estimated on post-shock data would
co-move with the treatment. Three constructions are admissible:
\begin{enumerate}
  \renewcommand{\labelenumi}{(\roman{enumi})}
  \item \textbf{Frozen baseline partition.} The partition $b(p)$ and each firm's core
        block $c(f)$ are estimated once, on the pre-shock baseline window, and held
        fixed; only the firm's RCA incidence $M_{fpt}$ varies over time. The
        panel variation in $d_{\rm in}/d_{\rm out}$ is then unambiguously the firm
        relocating against a frozen, predetermined grid.
  \item \textbf{Predetermined evolving partition.} The map is allowed to evolve, but
        estimated only on \emph{pre-shock} slices through the temporal multilayer
        modularity of \citet{mucha2010}---with inter-slice coupling and resolution
        chosen by partition stability \citep{bazzi2016}---and then extrapolated forward.
        This nests (i) as the strong-coupling limit and captures any pre-existing drift
        in the capability space, at the cost of measurement error that grows with the
        post-shock horizon.
  \item \textbf{Contemporaneous multilayer.} The map evolves using all years including
        post-shock. Under our continuous treatment this construction is \emph{not}
        causally identified: with no never-treated firms, there is no
        treatment-independent control group from which to anchor the map, so the
        partition---and hence the $(d_{\rm in}+d_{\rm out})$ denominator---remains a
        function of the dose. We therefore use it only descriptively, to document how
        the capability space reorganises, and never as the causal outcome.
\end{enumerate}
\textbf{We adopt (i) as the headline specification}, precisely because it has the
smallest researcher degrees of freedom and the cleanest exogeneity argument: the map
literally precedes the shock and cannot respond to it. Its one weakness---that a frozen
map grows stale---generates classical measurement error and therefore attenuation
\emph{against} our hypotheses, an honest bias rather than a confounding one. We
neutralise even that concern by reporting (ii) as robustness, which lets the map evolve
on pre-shock information and leaves the estimates intact, and by re-estimating (i) with
a leave-one-out partition that excludes each focal firm from the map that classifies
it, breaking the reflexive channel. Two construction rules are stated explicitly and
checked: products with no baseline block are assigned by product-space proximity (or
dropped, with the affected export share reported), and firms born after the baseline
are handled within the exit-selection analysis rather than assigned a post-baseline
core. We compute RCA with the standard contemporaneous, \emph{within-China} Balassa
denominator: because the capability blocks are themselves built from the Chinese
firm--product network, a China-relative benchmark---specialisation relative to the
Chinese export economy---is the internally consistent and economically natural object,
and we retain it as the headline measure. The denominator carries a second-order
\emph{inference} concern rather than an economic one: under a sector-wide shock the
aggregate Chinese export of a treated product falls for all exposed firms, mechanically
inflating the within-sample benchmark and potentially moving threshold crossings for
reasons unrelated to the firm's own portfolio choice. We treat this as robustness, not
as a change of measure. We re-estimate with a leave-one-out within-China denominator,
which preserves the China-relative interpretation while removing the firm's reflexive
contribution, and, as a more aggressive check, against a rest-of-world (ex-China) BACI
benchmark, which trades the China referent for full exogeneity to the Chinese firms'
treated response.

\paragraph{An elegant alternative under a modified treatment.}
We note one design under which the \emph{contemporaneous} evolving map (iii) would
become causally admissible: were we to trade the continuous dose for a binary,
staggered treatment, the not-yet-treated firms would furnish a treatment-independent
control group from which the time-evolving partition could be estimated at each date
\citep[in the spirit of the control-group construction of][]{callawaysantanna2021},
yielding a map that is simultaneously evolving and exogenous. This is a genuine
optimisation of elegance over the frozen design, but it is purchased by discarding the
exposure-intensity information that motivates our continuous specification, and it
relies on a credible never-/not-yet-treated set under the diffuse Teti tariff
schedule. We therefore retain the continuous treatment with the frozen map as our
primary design, and present the binary control-group-map construction as a deliberate
robustness fork rather than the main specification.

\paragraph{Two concessions we make at the outset.}
Neither threat is specific to the network construction, and we address both head-on
rather than conceal them. First, the shock induces firm exit \citep{crozet2021}; since
an exited firm has no $d_{\rm in}/d_{\rm out}$, survivors are selected, and we therefore
estimate the effect on the probability of exit, bound the survivor effect, and read our
portfolio estimates as conservative. Second, a tariff on one product reallocates demand
to substitutes held by other firms in the same block, so the stable-unit-treatment
assumption is strained; the \citet{adao2019} correction absorbs the resulting
cross-firm correlation in the standard errors but not the causal interference, which we
discuss and, where feasible, probe with leave-out-block exposure.

\paragraph{Firm-level data and the ownership moderator.}
The panel of Chinese exporting firms comes from the NBS Annual Survey of Industrial
Enterprises, used in the same form by \citet{guariglia2011}. That paper establishes
three features of the NBS data that shape our design: SOEs operate under soft budget
constraints and face negligible financial frictions, private domestic firms are
severely financially constrained (near-unit-root self-financing), and coastal firms
without foreign capital face the binding constraint. Because firms switch ownership
status over the panel, we assign type by a consistent across-years rule following
\citet{guariglia2011}. The ownership heterogeneity motivates our SOE--private split:
privately owned firms facing financial constraints respond more sharply to export market
closures because they cannot subsidise loss-making product lines from retained earnings,
making the $d_{\rm out}$ contraction more immediate and the $d_{\rm in}$ consolidation
more forced.

% ---------------------------------------------------------------------------
\subsection{Our Contribution}
\label{subsec:contribution}
% ---------------------------------------------------------------------------

This paper contributes to all five literatures described above. To the economic-complexity
literature, we provide the first causal identification of a firm-level capability
response---$d_{\rm in}$, $d_{\rm out}$, and $\text{Share\_core}$---to an exogenous external shock,
demonstrating that the block-modular firm--product partition is not merely a descriptive
device but a meaningful predictor of differential adjustment paths. To the capability-theory
and evolutionary-economics literatures, we offer empirical content for the
exploitation--exploration tension formalised by \citet{march1991}: the direction of
adjustment under tariff shocks is shown to be systematically moderated by regional
strategic dependence, consistent with the Penrose--Teece prediction that capability
redeployment is constrained by the jointness of underlying knowledge inputs, while the
magnitude is consistent with the permanent-effect predictions of \citet{storper1992}
and the lock-in mechanism documented in the Schumpeterian Mark~II literature.

To the trade literature, we shift the outcome from performance (the productivity,
markup, and employment measures that dominate the Pavcnik--Topalova--Amiti generation) to
capability architecture, providing evidence that the same tariff shocks that raise
aggregate TFP through competitive selection simultaneously compress the capability
frontier of exposed firms in ways that the performance lens does not reveal. To
the strategic-dependency literature, we take the Farrell--Newman chokepoint
hypothesis to firm-level data: firms operating in CDI-dependent product
categories display systematically larger capability consolidation effects than firms
in non-dependent categories, and this heterogeneity is not explained by tariff
magnitudes alone, consistent with the dual treatment of higher tariffs and active
EU reshoring policy targeting the same product space. Finally, to the identification
literature, we demonstrate that a continuous-dose distributed-lag event study---with the
heterogeneity-robust continuous-treatment estimators of \citet{dechaisemartin2020} and
\citet{callawaygoodmanbacon2024} as robustness---is tractable at the scale of
full-universe firm--product panels, that the \citet{adao2019} shift-share-robust
correction materially affects standard errors in heterogeneous-exposure settings with
product-level clustering, and that reading a network-derived outcome against a
\emph{predetermined} (frozen, or pre-shock-estimated) capability partition is what keeps
the estimand exogenous to the shock.

% =============================================================================
\section{Expected Dynamics: What the Literature Predicts We Should Observe}
\label{sec:expected}
% =============================================================================

Because the empirical estimates are not yet in hand, it is worth making the priors
explicit. This section reads the literatures of Section~\ref{sec:lit} \emph{forward}:
for each outcome and each moderator, it states what dynamic the prior work leads us to
expect once the planned design is executed, what sign each hypothesis carries, and
what magnitude range is plausible. The aim is to specify in advance what a result
``consistent with the literature'' would look like---so that a surprising estimate can
be recognised as surprising---and to organise the diverse predictions into a single
testable structure. Nothing here is a finding; everything here is a prediction that
the data will confirm, refute, or qualify.

% ---------------------------------------------------------------------------
\subsection{The Out-of-Block Margin ($d_{\rm out}$)}
\label{subsec:exp_dout}
% ---------------------------------------------------------------------------

Three literatures converge on a contraction of $d_{\rm out}$ following an adverse
tariff shock. The sunk-cost extensive-margin mechanism of \citet{crozet2021} predicts
that exposed firms stop serving the products and markets that become tariff-burdened,
and---because re-entry costs are 19--23 percent higher for lapsed exporters---the
contraction is \emph{hysteretic}: it should not reverse within the event window even
if tariffs are later rolled back. The variance-cost mechanism of \citet{crozet2025}
and \citet{vannoorenberghe2016} predicts that firms which try to \emph{preserve}
$d_{\rm out}$ by pivoting to unfamiliar, structurally dissimilar destinations inherit
higher demand volatility, so the privately optimal response for all but the largest
firms is to let $d_{\rm out}$ fall rather than chase replacement markets. And the
diversification-discount search models of \citet{matsusaka2001} and \citet{bernardo2002}
predict that only firms whose core match remains above the liquidation threshold retain
the option to explore; below it, exploratory $d_{\rm out}$ search collapses. The net
expectation is a \emph{negative} average effect on $d_{\rm out}$, concentrated among
smaller, financially constrained, and inexperienced firms, and persistent over the
post-shock horizon. The one literature that cuts the other way is the escape-competition
margin of \citet{aghion2005} and \citet{akcigit2022}: a high-capability firm facing
competitive pressure may innovate \emph{outward}, lifting $d_{\rm out}$. Because that
channel is confined to frontier firms, the expectation is a negative average masking a
positive tail---visible only in the interaction with firm capability rank, not in the
mean.

% ---------------------------------------------------------------------------
\subsection{The In-Block Margin ($d_{\rm in}$) and Retrenchment}
\label{subsec:exp_din}
% ---------------------------------------------------------------------------

The prediction for $d_{\rm in}$ is the cleaner of the two. The Typical-Product-Vector
consolidation of \citet{fontagne2018}---a 7.7 percentage-point retention premium for
core products in treated sectors---is, mechanically, a $d_{\rm in}$ response to a
trade-policy shock, and it is the single closest empirical antecedent for our main
outcome. The scope-economy variance argument of \citet{bottazzi2006} predicts that a
firm rebalancing a disrupted portfolio minimises revenue covariance by deepening within
a related cluster rather than scattering, which raises $d_{\rm in}$ relative to
$d_{\rm out}$. The competition-to-coherence evidence of \citet{valvano2003}---firms in
single-market-exposed sectors became more coherent over time---predicts the same sign.
And the Schumpeterian Mark~II lock-in of \citet{dosi1995} predicts that the retrenchment
is strongest exactly where capabilities are most cumulative. The expectation is
therefore a \emph{positive} average effect on $d_{\rm in}$, larger for firms in
complex, cumulative technological regimes, and---per \citet{storper1992} and
\citet{riverabatiz1991}---persistent rather than transitory, because what is lost when
a market closes is more than demand: it is the collective learning infrastructure that
sustained the broader basket.

% ---------------------------------------------------------------------------
\subsection{Share\_core and the Direction of Adjustment}
\label{subsec:exp_share}
% ---------------------------------------------------------------------------

The summary statistic $\text{Share\_core} = d_{\rm in}/(d_{\rm in}+d_{\rm out})$ signs
the exploitation--exploration choice of \citet{march1991}. If the $d_{\rm out}$
contraction and the $d_{\rm in}$ consolidation predicted above both materialise,
$\text{Share\_core}$ rises---the retrenchment hypothesis (H1). March's own result that
organisations systematically over-exploit makes this the modal expectation. But the
literature is deliberately not unanimous: the dynamic-capabilities reading
(\citealp{penrose1959}; \citealp{teece1982}) and the moderate-shock search prediction
of \citet{matsusaka2001} leave open the exploration case (H1$'$), in which $\text{Share\_core}$
\emph{falls} because firms with slack treat the shock as a prompt to reallocate toward
new capability bundles. The empirical question is therefore which force dominates on
average, and, more informatively, where in the firm distribution the sign flips. The
strong prior, following \citet{destefano2025}, is that the average response is toward
the growth-\emph{negative} core (H1), because the growth-positive $d_{\rm out}$ margin
is precisely the one that sunk costs and variance costs make hardest to defend under
stress.

% ---------------------------------------------------------------------------
\subsection{Heterogeneity and the Strategic-Dependence Moderator}
\label{subsec:exp_het}
% ---------------------------------------------------------------------------

The sharpest predictions are conditional, not average. The Farrell--Newman chokepoint
logic \citep{farrell2019}, operationalised through the CDI product set of
\citet{arjona2023} and the network-fragility tiers of \citet{korniyenko2017}, predicts
that firms whose baskets concentrate in strategically dependent products face not a
neutral tariff but a tariff \emph{plus} active reshoring policy aimed at the same
products---so their retrenchment should be amplified and, given the escalation dynamic
of \citet{evenett2024}, more persistent. The financial-constraint heterogeneity of
\citet{guariglia2011} predicts that private Chinese firms retrench more sharply than
soft-budget SOEs. The size-and-volatility composition effect of
\citet{vannoorenberghe2016} predicts that small firms cannot afford to defend
$d_{\rm out}$ at all. The institutional-contingency result of \citet{delios2008} and
\citet{boschma2015} predicts that the related-versus-unrelated payoff itself shifts
with the regional institutional environment, so that ExpDep should be read as a
production-system attribute and not only a firm attribute. The blockade-logic taxonomy
of \citet{fontana2024} predicts that cumulative-strategic sectors retrench harder than
infrastructure-strategic ones. The unifying expectation is that \emph{heterogeneity
dominates the average}: a modest mean effect concealing a steep gradient in strategic
dependence, financial constraint, and capability rank.

% ---------------------------------------------------------------------------
\subsection{Synthesis: Hypotheses, Signs, and Magnitude Benchmarks}
\label{subsec:exp_synth}
% ---------------------------------------------------------------------------

Table~\ref{tab:hyp} collects the hypotheses, their literature sources, and the sign
each predicts. Table~\ref{tab:mag} brackets the plausible magnitude range using two
external benchmarks: a lower bound from pure network-disruption transmission
(\citealp{korniyenko2017}: a 1pp rise in risky imports from a disaster-struck origin
cuts the importer's own exports by $\sim$0.7 percent, with no policy response) and an upper
bound from the most exposed economies in a targeted tariff war (\citealp{teti2025}:
$-1.6$ percent cumulative real income for Canada and Mexico under the 2018--2025 US tariffs,
versus $-0.3$ percent for China in the aggregate---a country-level benchmark we map onto the
most exposed Chinese firms).

\begin{table}[ht]
\centering
\caption{Hypotheses, literature sources, and predicted signs.}
\label{tab:hyp}
\footnotesize
\setlength{\tabcolsep}{4pt}
\renewcommand{\arraystretch}{1.3}
\begin{tabular}{|c|p{4.4cm}|p{4.2cm}|p{2.5cm}|}
\hline
\textbf{Hyp.} & \textbf{Statement} & \textbf{Principal literature source} & \textbf{Predicted effect} \\
\hline
H1 & Adverse tariff shock raises $\text{Share\_core}$ / $d_{\rm in}$ (retrenchment to the coherent core). &
\citet{fontagne2018}; \citet{bottazzi2006}; \citet{valvano2003}; \citet{dosi1995}; \citet{march1991} &
$d_{\rm in}\uparrow$, $\text{Share\_core}\uparrow$ \\
\hline
H1$'$ & Adverse shock raises $d_{\rm out}$ / lowers $\text{Share\_core}$ (exploration beyond the core). &
\citet{penrose1959}; \citet{teece1982}; \citet{matsusaka2001}; \citet{aghion2005} &
$d_{\rm out}\uparrow$ (frontier-firm tail only) \\
\hline
H2 & The retrenchment response is amplified in more strategically dependent production systems. &
\citet{farrell2019}; \citet{arjona2023}; \citet{korniyenko2017}; \citet{evenett2024} &
$\beta_{2}>0$ on Tariff $\times$ ExpDep \\
\hline
H3 & The $d_{\rm out}$ contraction is hysteretic (does not reverse within the window). &
\citet{crozet2021}; \citet{storper1992}; \citet{riverabatiz1991} &
persistent, asymmetric \\
\hline
H4 & Responses are steeper for financially constrained / private / small firms. &
\citet{guariglia2011}; \citet{vannoorenberghe2016}; \citet{coad2021} &
steeper gradient \\
\hline
\end{tabular}
\end{table}

\begin{table}[ht]
\centering
\caption{Magnitude benchmarks for the $d_{\rm out}$ response (per s.d.\ of TariffShock).}
\label{tab:mag}
\footnotesize
\setlength{\tabcolsep}{4pt}
\renewcommand{\arraystretch}{1.3}
\begin{tabular}{|p{3.8cm}|c|p{6.2cm}|}
\hline
\textbf{Benchmark} & \textbf{Magnitude} & \textbf{Interpretation} \\
\hline
Lower bound --- pure network disruption \citep{korniyenko2017} & $\sim 0.7\%$ &
Supply-chain transmission of a localised shock, no policy intervention. \\
\hline
Upper bound --- most exposed economies, targeted tariff war \citep{teti2025} & $\sim 1.6\%$ &
Country-level loss for Canada/Mexico under the 2018--2025 US tariffs (vs.\ $-0.3\%$ for China), mapped onto the most exposed firms. \\
\hline
\end{tabular}
\end{table}

% =============================================================================
\section{Open Questions and a Forward Research Agenda}
\label{sec:agenda}
% =============================================================================

The review above leaves a small number of questions genuinely open and opens a larger
number of new ones. This section separates the two: first, the specific gaps in the
existing literature that \emph{this} paper is designed to close; second, the spin-off
questions that the same body of work raises but that lie beyond the present design and
would each support a distinct paper.

% ---------------------------------------------------------------------------
\subsection{The Gaps This Paper Fills}
\label{subsec:gaps_filled}
% ---------------------------------------------------------------------------

\begin{enumerate}
\item \textbf{Causal capability reconfiguration, not descriptive structure.} The EFC
literature (\citealp{bruno2023}; \citealp{pugliese2019}; \citealp{destefano2025})
measures block-modular capability structure but almost always descriptively. No prior
work links a \emph{predetermined, exogenous} trade-policy shock to the in/out
decomposition of firm diversification. This paper supplies that causal link.

\item \textbf{Composition, not performance.} The trade-and-productivity literature
(\citealp{pavcnik2002}; \citealp{topalova2011}; \citealp{amiti2007}; \citealp{tomasi2022})
measures how tariff shocks move productivity, markups, and survival. Because
performance dimensions decouple from capability (\citealp{grazzi2012};
\citealp{coad2021}; \citealp{syverson2011}), these outcomes are silent on \emph{which}
activities firms keep, drop, or enter. This paper moves the outcome to the architecture
of the export basket.

\item \textbf{Firm-level test of the chokepoint hypothesis.} The strategic-dependency
literature (\citealp{farrell2019}; \citealp{euc2021}; \citealp{arjona2023};
\citealp{korniyenko2017}) maps dependencies at the product and country level but has no
firm-level causal test. By interacting the tariff shock with a firm-specific
strategic-dependence exposure, this paper provides that test.

\item \textbf{A signed reading of exploration vs.\ exploitation under a real shock.}
March's (\citeyear{march1991}) tension has rich theory but little exogenous-shock
evidence at the firm level. The $\text{Share\_core}$ outcome under a predetermined tariff shock
delivers a directly signed estimate of which way firms move when their core is
threatened.

\item \textbf{Scalability of continuous-treatment event-study estimators.} We demonstrate
that a continuous-dose distributed-lag design with \citet{adao2019} shift-share inference---and
the heterogeneity-robust continuous-treatment estimators of \citet{dechaisemartin2020} and
\citet{callawaygoodmanbacon2024}---is tractable on a full-universe firm--product panel, a
methodological contribution in its own right for the heterogeneous-exposure shift-share
literature.
\end{enumerate}

% ---------------------------------------------------------------------------
\subsection{Beyond This Paper: Spin-Off Research Questions}
\label{subsec:spinoffs}
% ---------------------------------------------------------------------------

Each of the following is motivated by the reviewed literature but deliberately left
outside the present design; each would constitute a separate study.

\paragraph{(a) The firm--product--destination tensor.}
\citet{destefano2025} works in pure firm--product space; \citet{crozet2025} and
\citet{vannoorenberghe2016} show that the destination dimension carries its own
volatility structure. A natural extension promotes the bipartite firm--product matrix
to a firm--product--destination tensor and asks whether capability blocks reorganise
\emph{geographically}---whether retrenchment to the core is also a retreat to familiar
markets. This is the volatility-minimising hypothesis of $\Phi$-diversity made
firm-specific.

\paragraph{(b) Contemporaneous temporal-coupled partitions as a causal object.}
Our headline design freezes the baseline block, and our robustness lets it evolve on
\emph{pre-shock} information through the temporal multilayer machinery of \citet{mucha2010}
(with coupling and resolution chosen as in \citealp{bazzi2016}). What we deliberately leave
open is the \emph{contemporaneous} version: treating \emph{block reconfiguration itself}---the
birth, death, and merger of capability clusters measured with post-shock data---as the causal
outcome. Under our continuous treatment this is not identified, because there is no
never-treated control group from which to anchor an evolving map; a dedicated study could
recover it under a binary, staggered design by estimating the partition from not-yet-treated
firms, asking how trade shocks reshape the modular geometry of an entire industry rather than
the position of a single firm within a fixed geometry.

\paragraph{(c) The antidumping and safeguard stage.}
The Teti GTD \citep{teti2024} captures statutory MFN and preferential tariffs but
\emph{not} antidumping, safeguard, or trade-war tariffs. Pairing it with a temporary-trade-barriers
database would isolate the capability response to discretionary, firm-targeting
protection---plausibly a sharper and more anticipatory shock than scheduled tariff
changes, and a cleaner test of the escalation dynamic of \citet{evenett2024}.

\paragraph{(d) Instrumenting the policy-escalation component.}
\citet{evenett2024} document that industrial-policy intensity rises in election years
and that subsidies are matched tit-for-tat with 73.8 percent probability. An election-cycle
or tit-for-tat instrument for the policy-driven component of a firm's tariff shock
would separate the escalation channel from the mechanical tariff change---turning the
escalation caveat of this paper into an identification strategy of its own.

\paragraph{(e) Green-capability reconfiguration.}
\citet{mealy2022} show the complexity toolkit applies to any product subset and that
green-capability accumulation is strongly path dependent. Re-running the present design
with the green-product set as the treated category would ask whether decarbonisation
policy and trade policy push firms in the same or opposite capability directions---a
direct bridge between the strategic-autonomy agenda of \citet{cagnin2021} and the green
transition.

\paragraph{(f) Institutional voids as a moderator.}
\citet{delios2008} and \citet{boschma2015} show that the value of diversification, and
the related-versus-unrelated balance, depend on institutional quality. Exploiting
within-China regional institutional variation as a second moderator would test whether
capability retrenchment under trade shocks is steeper or shallower where market
institutions are weaker, distinguishing capability exploitation from
institutional substitution as motives for diversification.

\paragraph{(g) The gain/loss asymmetry of complexity.}
\citet{nomaler2023} show that lost specialisations are often \emph{more} complex than
gained ones, and that density-based relatedness is heavily confounded by diversity and
ubiquity. A study that scores firm-level product \emph{exits} by complexity---asking
whether tariff shocks make firms shed their most sophisticated products first---would
test whether trade shocks erode the capability frontier from the top, a sharper and
more policy-relevant question than net diversification counts can answer.
 into the main paper.
% Cite keys follow the convention author(s)year; matching @article / @book /
% @techreport entries live in Bibliography_base.bib (86 entries as of 2026-06-05).
%
% STRUCTURE
%   2   Related Literature
%     2.1  Economic complexity, capability structure, and diversification theory
%     2.2  Capability theory, corporate coherence, and the direction of adjustment
%     2.3  Trade policy shocks and firm adjustment
%     2.4  Strategic dependencies and geopolitical trade shocks
%     2.5  Identification
%     2.6  Contribution
%   3   Expected Dynamics: What the Literature Predicts We Should Observe
%     3.1  The out-of-block margin (d_out)
%     3.2  The in-block margin (d_in) and retrenchment
%     3.3  Share_core and the direction of adjustment
%     3.4  Heterogeneity and the strategic-dependence moderator
%     3.5  Synthesis: hypotheses, signs, and magnitude benchmarks  [2 tables]
%   4   Open Questions and a Forward Research Agenda
%     4.1  The gaps this paper fills
%     4.2  Beyond this paper: spin-off research questions
%
% NOTE: tables use plain \hline tabular (no booktabs dependency) so the file
% compiles stand-alone via lit_review_test.tex and when \input into the paper.
% =============================================================================

\section{Related Literature}
\label{sec:lit}

This paper sits at the intersection of four bodies of work: the economic-complexity
literature on firm-level capability structure and the theory of diversification; the
capability-based theory of the firm and evolutionary economics; the empirical trade
literature on the effects of tariff shocks on firm outcomes; and the emerging literature
on strategic dependencies, geopolitical risk, and the weaponisation of economic
interdependence. A fifth, methodological strand covers identification in shift-share
designs with heterogeneous treatment effects. The unifying thread across these
literatures is a single dynamic: an external shock that degrades the value of a firm's
existing activities forces a \emph{reconfiguration} of its productive knowledge, and
the literatures disagree on whether that reconfiguration deepens a coherent core or
pushes outward into new capability bundles. We discuss each stream in turn, emphasising
what each predicts about the \emph{direction} and \emph{persistence} of adjustment, and
close with a statement of our contribution.

% ---------------------------------------------------------------------------
\subsection{Economic Complexity, Capability Structure, and Diversification Theory}
\label{subsec:lit_efc}
% ---------------------------------------------------------------------------

The product-space framework of \citet{hausmann2007} and \citet{hidalgo2009}
established that countries export products in patterns that are far from random:
goods that require similar underlying capabilities tend to be co-exported, generating
a proximity structure in the product space that conditions the path of structural
transformation. \citet{hidalgo2011} formalise this intuition through a bipartite
capabilities model in which countries and products are represented as the two node
sets of a bipartite graph, and the adjacency matrix $M_{cp}=1$ whenever country $c$
has revealed comparative advantage (RCA) in product $p$. The key result is that
diversification is convex in the capability stock: countries with few capabilities
earn near-zero marginal returns from acquiring any single new one, generating the
low-income quiescence trap that has been documented empirically across the
income distribution \citep{hausmann2007}. This non-linearity is fundamental to the
agenda of using complexity tools in a causal setting: treatment-effect heterogeneity
by capability endowment is an architectural feature of the framework, not an
add-on.

The Fitness-Complexity (FC) algorithm of \citet{cristelli2013} provides the
operational complexity metric adopted in the subsequent firm-level literature. Rather
than the linear iterative procedure of \citet{hidalgo2009}, FC defines country fitness
as the sum of product complexities and product complexity as the harmonic mean of
exporter fitness---a bounded operator that uniquely rewards diversification and
assigns complexity to a product according to the least capable exporter that ships it,
giving the algorithm its asymmetric censoring property. The algorithm converges to a
unique fixed point and, as shown by \citet{cristelli2017}, generates a selective
predictability scheme in which high-fitness countries follow laminar, predictable
growth trajectories while low-fitness countries exhibit chaotic dynamics---a property
that has direct implications for heterogeneous treatment responses in our setting:
firms embedded in high-complexity capability clusters should display more systematic
responses to tariff shocks than firms operating at the periphery of the product
space. \citet{angelini2017} characterise the long-run dynamics of products in the complexity
plane, showing that they drift toward a stationary configuration so that
shock-induced dislocations are transient and mean-reverting---which is why we treat
product complexity as a well-defined pre-shock control and read post-shock movements
against a fixed baseline.
\citet{stojkoski2025} extend the framework to an optimisation context,
showing that the ECI-growth relationship is stronger at lower income levels and that
relatedness is one of several predictors of product-entry, a finding that validates
the multi-dimensional block structure we exploit over simpler relatedness-only
measures. \citet{koch2021} sounds a useful caution about what such metrics measure: the
complexity--growth link is in fact \emph{stronger} when complexity is computed on
value-added rather than gross exports ($R^2 = 0.63$ versus $0.31$), so gross-flow
complexity blurs the connection to realised value capture. Since our firm-customs
measures are necessarily built on gross flows, we read them as capturing capability
\emph{structure} rather than realised value capture, which is precisely why we treat
the block partition as a measure of capability architecture rather than of performance.

Three strands sharpen what ``relatedness'' and ``diversification'' actually mean for
the direction of adjustment. First, the related-variety tradition of
\citet{frenken2007} distinguishes \emph{related} variety (entry into activities that
share capabilities, which drives regional \emph{employment} growth through
within-sector Jacobs spillovers---though, tellingly, not productivity growth) from
\emph{unrelated} variety (which provides a portfolio hedge that lowers unemployment
under asymmetric sectoral shocks). This is the regional-economics antecedent of our
$d_{\rm in}$ versus $d_{\rm out}$ distinction. \citet{boschma2015} show that the relative payoff to
related versus unrelated diversification is institutionally contingent: coordinated
market economies favour incremental related diversification, while liberal market
economies support the unrelated, exploratory kind, a result that motivates treating
the regional production system, not just the firm, as a unit of heterogeneity.
\citet{freire2019} provides a structural model of endogenous diversification in which
capability accumulation and product entry co-evolve, giving microfoundations to the
empirical relatedness regularities. Second, \citet{mealy2022} demonstrate that the
complexity toolkit transfers to any policy-relevant subset of products---their
green-complexity indices show strong path dependence in capability accumulation
(a $0.92$ correlation between green-complexity potential and the green-complexity
index, with beginning-of-period potential predicting subsequent green-export
growth), a template we adopt when restricting attention to strategically dependent
product categories. Third, \citet{nomaler2023} issue a confounding warning that
directly disciplines our outcome construction: density-based relatedness metrics are
$\sim$93 percent explained by diversity and ubiquity alone, so a measured ``relatedness''
effect can be an artefact of how many products a firm makes and how common those
products are. Their symmetric treatment of RCA \emph{gains} and \emph{losses}, together with
the finding that lost specialisations are often \emph{more} complex than gained
ones, reframes diversification as a two-sided entry-and-exit process and is one of the
closest conceptual antecedents to our reconfiguration question.

The application of complexity tools to firm-level data is recent. \citet{bruno2018, bruno2023} construct firm-product bipartite
networks from Customs microdata and show that RCA-based binarisation
(with BiCM entropy-null correction as a robustness) yields block-modular structures
that are stable across years and meaningfully predict firm growth, innovation, and
survival. \citet{pugliese2019} show that firms diversify coherently within their
existing technological portfolio, so that a firm's position within its capability
cluster is informative about its future mobility---a finding that motivates our
within-block heterogeneity analysis; the country-level analogue is \citet{pugliese2017},
who show that complex economies escape the low-income trap by moving laterally into
related capabilities. \citet{coad2021} push
this furthest, applying the FC algorithm to a firm$\times$capability matrix for $\sim$40{,}000
Indian firms and recovering a nested ``capabilities ladder'' whose rungs predict growth
and survival but---tellingly---not profitability, a decoupling we return to in
Section~\ref{subsec:lit_cap}. \citet{destefano2025}
provide the most direct methodological antecedent for our outcome measures: they
decompose firm diversification into in-block ($d_{\rm in}$) and out-of-block
($d_{\rm out}$) counts using the block partition produced by the BRIM algorithm of
\citet{barber2007}, define $\text{Share\_core} = d_{\rm in}/(d_{\rm in}+d_{\rm out})$,
and show that out-of-block diversification is the growth-positive margin while
retrenchment toward the core---the $d_{\rm in}$ direction---is the
growth-negative, lock-in-reinforcing margin. We adopt their exact outcome measures
and employ the same partition algorithm.

The network-method foundations underpinning the partition are well established.
\citet{newman2006} introduced the modularity objective and its spectral optimisation;
\citet{blondel2008} provide the Louvain heuristic that makes community detection
tractable at the scale of full-universe firm-product matrices; \citet{ducharenas2005}
supply the extremal-optimisation alternative we use as a partition robustness check;
and \citet{fortunato2016} provide the authoritative critical survey of the field's
validation pitfalls. \citet{barber2007} adapts modularity to the bipartite case (BRIM),
and \citet{soleribalta2018}
show that the BRIM objective is insensitive to within-block
nested structure and can therefore miss the hierarchical ranking of firms inside a
capability block; their in-block nestedness quality function $I$ jointly optimises the
block partition and within-block NODF \citep[the consistent nestedness metric of][]{almeidaneto2008},
providing the basis for our within-block rank ordering of firms. \citet{mariani2019} prove that the FC algorithm is a nestedness
maximiser and clarify the relationship between modularity and nestedness in firm-level
networks: the two are positively correlated at low network densities (the regime typical
of firm-product matrices) and can be jointly estimated. Validation of the BRIM
partition against entropy null models---using the Bipartite Configuration Model of
\citet{squartini2011} and \citet{saracco2015}, with the projection reconciliation of
\citet{cimini2022}---is a required pre-processing step
to ensure that block structure reflects genuine capability clustering rather than
degree-sequence artefacts, and we implement it throughout. A wider set of Economic Fitness and
Complexity (EFC) studies---broader applications of economic-complexity methods---rounds
out the methodological base: global-value-chain specialisation \citep{bontadini2021, bontadini2025},
the complexity of structural change and the sustainability transition
\citep{caldarola2024a, caldarola2024b}, machine-learning reconstruction of trade flows
\citep{gnecco2023}, and relatedness-based forecasting of industrial upgrading
\citep{albora2023a, albora2023b}.

% ---------------------------------------------------------------------------
\subsection{Capability Theory, Corporate Coherence, and the Direction of Adjustment}
\label{subsec:lit_cap}
% ---------------------------------------------------------------------------

The theoretical framework connecting firm capability stocks to the direction of
diversification draws from the resource-based view of the firm, evolutionary economics,
and the learning-and-adaptation literature. \citet{penrose1959} establishes that
firms are bundles of productive resources whose services can be redeployed across
activities at varying cost: the excess services that routinely arise from indivisible
resources create an organic pressure toward related diversification that is internal
to the firm, not driven by market structure. \citet{teece1982} formalises this in a
scope-economy model showing that the boundaries of the multi-product firm are
determined by the jointness of underlying knowledge inputs; firms expand into activities
that share tacit knowledge with existing lines, generating the coherence property
that the EFC literature quantifies through block-modular partitions.
\citet{teece1994} make this measurable through the survivor-principle relatedness
metric---sectors are ``related'' when more firms combine them than chance predicts---which
is the product-space analogue of the bipartite co-occurrence logic at the heart of the
EFC programme. The corporate-coherence literature that builds on it is directly
informative about the \emph{dynamics} we expect to observe. \citet{piscitello2000}
finds that large firms reduced product coherence while \emph{maintaining}
technological coherence over 1977--1995: capabilities are the stable organising
principle even as product portfolios churn, which is exactly why a capability-block
partition should track adjustment more faithfully than a raw product count.
\citet{valvano2003} provide the most pointed antecedent for a trade-shock prediction:
in Italian manufacturing, firms in sectors more exposed to single-market competition
(``sensitive'' sectors) became \emph{more} coherent over time, and entrants into the
leadership ranks were more coherent than incumbents---direct evidence that intensified
competition pushes surviving firms toward their core ($d_{\rm in}$). \citet{delios2008}
add the institutional contingency: within-country product diversification is more
valuable where market institutions are weak, because diversification substitutes for
missing capital, labour, and information markets---a moderation channel that warns
against reading every diversification response as pure capability exploitation.

The corporate-finance literature on the ``diversification discount'' supplies the
causal-direction caution that disciplines our interpretation. \citet{langstulz1994}
document a robust industry-adjusted discount in Tobin's $q$ of 0.35--0.50 for
diversified US firms throughout 1978--1990, but show that firms diversify \emph{after}
their core industry weakens---the discount precedes diversification rather than being
caused by it. \citet{matsusaka2001} rationalises this as capability-match search: a
firm whose current match deteriorates maintains its existing lines while experimenting
with new ones (the ``$b$-strategy''), liquidating and restarting only when the match
falls below a lower threshold. \citet{bernardo2002} recast the same fact as
real-option value: single-segment firms are worth more because they retain unexercised
expansion options. Translated to our setting, these models deliver a sharp,
non-obvious prediction: a \emph{moderate} tariff shock should push firms toward
exploratory $d_{\rm out}$ search (the $b$-strategy), whereas a \emph{severe} shock that
drives match quality below the liquidation threshold collapses the option value of
organisational capabilities and forces retrenchment or exit---so the sign of the
average response is genuinely ambiguous and should depend on shock intensity, exactly
the non-linearity our event study is designed to detect.

The evolutionary perspective of \citet{dosi1988, dosi1995} adds a regime dimension:
the rate and direction of capability accumulation depend on both the firm's own
endowment and the technological regime of the sector---the opportunity conditions,
appropriability, and cumulativeness that structure search. Under Schumpeterian Mark
II conditions (cumulative, large-firm-dominated regimes), incumbents sustain
technological advantages through internal learning, generating the technological
lock-in that maps directly onto high $\text{Share\_core}$. The implication for our
setting is that firms in complex, cumulative technological regimes should display
stronger $d_{\rm in}$ responses to adverse shocks---retrenchment to the core is both
the natural adjustment and the organisationally feasible one when search capabilities
are already embedded in existing product clusters.

\citet{march1991} provides the ambiguity at the heart of our hypotheses. The
exploitation--exploration tension implies that the direction of capability adjustment
under stress is not a priori signed: firms that face an adverse demand or competitive
shock in their existing core can either intensify exploitation of known competencies
(retrenchment) or accelerate exploration beyond the core (diversification-as-escape).
March's formal result is that organisations systematically over-exploit under most
plausible parameter configurations, predicting the retrenchment (H1) hypothesis as
the modal response but leaving open the exploration (H1$'$) case for firms with
sufficient slack and absorptive capacity. \citet{patel1997} document the empirical
content of this tension at the firm level, showing that large innovative firms spread
R\&D effort across fields but in highly persistent patterns---path dependence is the
rule, not the exception---while \citet{bell1993} establish that technological
capability accumulation in developing-country firms is even more path-dependent,
with explicit investments in learning required to escape existing trajectories.
\citet{storper1992} connects this to trade: product-based technological learning
generates agglomeration economies in particular product-territory combinations, so
that an adverse trade shock removes both demand and the collective learning
infrastructure that sustains capability accumulation---a channel that implies
persistent, not merely transitory, capability effects of tariff shocks.

The firm-growth literature, summarised in \citet{coad2009} and documented empirically
in \citet{bottazzi2003, bottazzi2005, bottazzi2006}, shows that firm-level growth is
Laplacian---fat-tailed, near-unit-root---and that the variance of growth rates scales
inversely with firm size at a rate explained entirely by diversification across product
lines. The scope-economy mechanism \citep{bottazzi2006} implies that firms which
diversify within a related cluster (high $d_{\rm in}$) reduce the co-movement of
their individual product revenue streams more effectively than firms that scatter across
unrelated products. Under tariff shock conditions, firms that lose access to existing
export markets face a portfolio rebalancing problem whose solution---from a pure
variance-minimisation standpoint---is to deepen within-block exposure rather than
scatter to unfamiliar destinations, reinforcing the retrenchment prediction. A
recurring caveat from this literature is that performance dimensions decouple:
\citet{grazzi2012} shows that Italian exporters carry a large productivity premium but
\emph{no} profitability premium, and \citet{coad2021} that capability-ladder position
predicts growth and survival but not profit margins; \citet{syverson2011} surveys the
enormous, persistent productivity dispersion that these patterns sit inside. The
collective lesson---that capability advantage need not show up in profits or even in
growth---is why we measure the \emph{architecture} of the export basket
($d_{\rm in}$, $d_{\rm out}$, $\text{Share\_core}$) rather than a performance scalar.
\citet{aghion2001, aghion2005} add the innovation margin: neck-and-neck competitors
facing increased competitive pressure innovate to escape competition, while laggards
reduce innovation effort. Translating to our setting, firms with high within-block
fitness ranks face incentives to deepen core capabilities under tariff pressure (the
escape-competition channel), while low-rank firms---unable to sustain the cost of
maintaining their current product portfolio---are more likely to prune $d_{\rm out}$
and consolidate around the most defensible core products. This generates the
within-block heterogeneity in treatment effects that we exploit empirically.

% ---------------------------------------------------------------------------
\subsection{Trade Policy Shocks and Firm Adjustment}
\label{subsec:lit_trade}
% ---------------------------------------------------------------------------

A large reduced-form literature establishes that tariff changes cause significant
adjustments in firm outcomes. \citet{pavcnik2002} and \citet{topalova2011} document
that trade liberalisation raises TFP through competitive selection and technology
imports in Chilean and Indian manufacturing, respectively. \citet{amiti2007} show,
for Indonesian firms, that input tariff reductions deliver TFP gains at least twice as
large as those from output tariff reductions---a finding driven by access to high-quality
imported intermediates---and construct the shift-share input-tariff exposure variable that is
the methodological template for our treatment. \citet{tomasi2022} replicate this
design for Vietnamese firms during the 2007 WTO accession, finding that governance
quality moderates the tariff-TFP effect, providing the closest precedent for our
moderated shift-share design with ExpDep as the heterogeneity variable.

What this literature consistently measures is performance: productivity, employment,
markups, survival. It is essentially silent on the \emph{composition} of a firm's
product portfolio---which activities are retained, dropped, or newly entered, and
whether the reconfiguration deepens a coherent core or spreads capabilities into
new domains. The EFC literature reviewed in Section~\ref{subsec:lit_efc} supplies
the tools to measure that composition, but has applied them almost exclusively in
descriptive, non-causal settings. This paper occupies the intersection left empty
by both strands. The theoretical channel linking trade exposure to capability
adjustment is supplied by the trade-and-innovation literature: \citet{akcigit2022}
unify the market-size, business-stealing, and escape-competition effects of trade on
innovation incentives, showing that the net effect on a firm's capability investment
is sign-ambiguous and depends on its competitive position---the firm-level analogue of
the exploitation--exploration tension. \citet{riverabatiz1991} establish the deeper
point that integration affects growth through flows of \emph{ideas} as well as goods,
so that a tariff shock that severs a firm from a technological frontier market can have
permanent (growth-rate), not merely level, effects on its capability trajectory---reinforcing the
persistence prediction from \citet{storper1992}.

The most direct evidence that trade shocks restructure capability portfolios rather
than merely rescaling performance comes from \citet{fontagne2018}, who show that
Chinese export quotas on textiles and clothing generated a significant \emph{Typical
Product Vector} consolidation among French exporting firms: products at the core of
a firm's stable portfolio were retained at a 7.7 percentage-point higher rate than
peripheral products in treated sectors. This is, in effect, a $d_{\rm in}$ response
to a trade-policy shock measured with product-vector tools---the closest antecedent
in the literature for our main outcome. \citet{fontagne2022} provide the
product-level trade elasticities that we use as Bartik weights: 5,050 HS6-digit
elasticities estimated via PPML gravity, averaging $-5.3$ (median $-3.5$), with
HS-section means spanning $-18.5$ (minerals), $-6.1$ (machinery), and $-3.6$
(footwear), capturing the heterogeneous disruption that a uniform tariff schedule
imposes on firms with different product portfolios.

The dynamic and firm-level dimensions of portfolio adjustment are examined by
\citet{crozet2021}, who study the effect of EU and US sanctions on the extensive
margin of French exporters to Iran and Russia. Using monthly customs data and a
dynamic probit model with three-way fixed effects and analytical bias correction,
they find that sanctions lower the probability of serving a sanctioned market by 39
percent (Iran) and 23 percent (Russia), with strong state dependence: the entry-cost
factor is 19--23 percent higher for firms not active in the previous year. Crucially,
the lifting of sanctions on Myanmar generates only a small and gradual positive response,
establishing an \emph{asymmetry} between imposition and removal that implies permanent,
not transient, $d_{\rm out}$ contraction even after the policy shock is reversed.
This hysteresis result---consistent with the Roberts--Tybout sunk-cost model---means
that treatment effects in our event study should be interpreted as reflecting a
permanent restructuring of the capability portfolio, not a temporary demand disruption.
The firm-level heterogeneity documented by \citet{crozet2021}---that trade-finance-dependent
firms, inexperienced firms, and non-specialist firms exit more sharply---maps
directly onto the moderation hypotheses we test: private Chinese firms (documented
by \citet{guariglia2011} as the most financially constrained in the NBS panel) should
display amplified $d_{\rm out}$ reductions, while SOEs with soft budget constraints
should display more muted responses.

A complementary strand examines the export portfolio diversification strategies that
firms deploy to manage demand volatility. \citet{vannoorenberghe2016} show, for
Chinese firm-level data, that destination diversification reduces volatility for large
exporters but \emph{increases} it for small ones, with the switching point determined
by a composition effect: occasional exports to marginal destinations dominate the
diversification gain for small firms, while stable core-market exports dominate for
large ones. \citet{crozet2025} formalise this puzzle through a phyloeconomic diversity
index $\Phi$ that adds a dissimilarity dimension to the standard richness-evenness
decomposition. Their central finding---that conditional on the number of destinations
and the Herfindahl concentration of export sales, higher $\Phi$ (more dissimilar
markets) is associated with \emph{higher} export sales volatility---reveals that
hedging through geographic diversification has limits: distant, structurally different
destination markets carry higher idiosyncratic demand variance, so that firms forced
to substitute familiar Western markets with unfamiliar alternatives do not achieve
the volatility reduction that naive diversification logic predicts. For our setting,
this implies that Chinese firms which maintain $d_{\rm out}$ by pivoting to non-Western
destinations face not only lower export levels but higher volatility---a double cost
that reinforces the variance-minimising case for $d_{\rm in}$ consolidation.
\citet{barbieri2024} provide a parallel decomposition at the territorial level:
\emph{sectoral} diversification (a proxy for productive capabilities, the conceptual
sibling of $d_{\rm in}$) raises export intensity, while \emph{geographic}
diversification (a proxy for geoeconomic capabilities, loosely analogous to the
resilience role we attribute to $d_{\rm out}$) is associated with greater export
stability; institutional quality independently improves both margins.\footnote{The
analogy is functional rather than literal. Barbieri's second axis is
\emph{destinations}, whereas our $d_{\rm out}$ remains a \emph{product}-space
measure (out-of-block products), following \citet{destefano2025}, who works in pure
firm--product space with no destination dimension. The roles align---one margin scales
output, the other buys resilience---but the underlying axes differ, and in their data
the dominant stabiliser is EU market exposure rather than geographic diversification
per se. We therefore treat the destination dimension as a second-stage extension
(Section~\ref{subsec:spinoffs}) rather than as a component of $d_{\rm out}$.}

We estimate treatment using the Global Tariff Database (GTD) of \citet{teti2024},
the first applied-tariff panel to cover bilateral HS6 tariffs continuously from 1988
to 2021 while correcting for the false interpolation and selection artefacts of UNCTAD
TRAINS. \citet{teti2019} establish, using structural gravity, that Rules-of-Origin
compliance is prohibitively costly for the machinery and electronics sectors at the
core of Chinese exports, confirming that MFN tariffs---the rate captured by the
GTD---are the economically relevant wedge for our firms. The welfare magnitudes from
\citet{teti2025} calibrate our expected effect sizes: a GE simulation of the
2018--2025 US--China tariff war estimates cumulative real-income losses of roughly
$-1.6$ percent for Canada and Mexico---economies with high export exposure to the US
market---against only $-0.3$ percent for China in the aggregate. We use the
$-1.6$ percent figure as an external benchmark for the most exposed Chinese firms:
by analogy with the highly exposed Canadian and Mexican case, the uneven distribution
of tariff exposure across product portfolios should push firms concentrated in
high-regulatory-dependence categories toward magnitudes well above the national
average. The mapping from these country-level losses to firm-level exposure is our
own; \citet{teti2025} simulate countries and US states, not firms. The distortionary
role of size-contingent and sector-specific regulation documented by
\citet{garicano2016}---who find welfare costs of 1.3--3.5 percent of GDP from French
size thresholds---is the domestic-policy analogue of this exposure heterogeneity and
motivates controlling for regulatory environment in the moderation analysis.

% ---------------------------------------------------------------------------
\subsection{Strategic Dependencies and Geopolitical Trade Shocks}
\label{subsec:lit_geo}
% ---------------------------------------------------------------------------

The design of the heterogeneity variable ExpDep draws from a separate and newer
literature on strategic economic dependencies. \citet{farrell2019} identify the
conditions under which the topology of global economic networks generates geopolitical
leverage: hub-and-spoke structures create chokepoint nodes through which flows can be
controlled (flow restriction) and from which information can be extracted (panopticon
effects). Crucially, weaponisation requires two joint conditions---jurisdictional
control over the chokepoint and domestic institutions capable of deploying that control
as a foreign policy instrument---which limits the set of products and bilateral
relationships that qualify as strategic in the relevant sense. The Farrell--Newman
framework grounds our moderation hypothesis: Chinese firms whose export baskets are
concentrated in products where Western jurisdictions have identified chokepoint control
face not a neutral tariff shock but one accompanied by active policy interest in
reducing dependence---generating an amplified and persistent treatment effect.

The quantitative operationalisation of this insight is provided by \citet{euc2021},
who develop the Core Dependency Indicators (CDI) methodology to identify EU strategic
dependencies at the HS6-product level. The methodology applies three thresholds:
supplier concentration (HHI of extra-EU import shares exceeding 0.4, implying fewer
than 2.5 effective suppliers), import reliance (extra-EU imports exceeding 50 percent
of total EU imports of the product), and substitutability (extra-EU import volume
exceeding total EU export volume, so that domestic production cannot cover the import
volume even in an emergency). Of approximately 5,000 HS6 products, 137 pass all three
thresholds in the EU's most sensitive ecosystems, with China accounting for roughly
52 percent of their import value. \citet{arjona2023} refine the methodology with a
dynamic four-year threshold and a re-export correction: accounting for transit trade
\emph{raises} the count of dependent products by roughly 27 percent, since ignoring
re-exports understates true extra-EU reliance, yielding a dependent-product set of
204 items that provides the product classification we use to construct ExpDep.
The forward-looking scenario work of \citet{cagnin2021} situates these indicators in
the EU's Open Strategic Autonomy agenda, projecting an escalation of dependency-reducing
policy through 2040 that makes the persistence of any treatment effect we estimate a
policy-relevant, not merely statistical, question.

The network-science analogue of the CDI methodology is provided by
\citet{korniyenko2017}, who combine three structural indicators of product trade networks---
the standard deviation of weighted outdegree centrality (capturing hub dominance), the
product of clustering coefficient and network diameter (capturing geographic
compartmentalisation), and a Revealed Factor Intensity measure of substitutability---to
classify 3,578 intermediate and capital goods into four risk tiers via k-median
clustering. Their persistently risky group comprises 421 products over 2003--2014,
concentrated in machinery and mechanical appliances (HS 84--85), transport equipment,
pharmaceuticals, and precision instruments. The causal validity of the classification
is demonstrated by two out-of-sample cases: the 2011 Japan earthquake and Thai floods,
where products subsequently in short supply were correctly identified as Group~4 one
year prior. The overlap between the Korniyenko Group~4 products and the
Arjona CDI-dependent set defines the most extreme ExpDep exposure: products
that are simultaneously structurally concentrated, topologically fragile, and actively
targeted by EU and US reshoring policy.

\citet{evenett2024} provide real-time evidence on the policy escalation dynamic. Their
New Industrial Policy Observatory (NIPO), covering 2,500$+$ new industrial policy
measures across 75 jurisdictions in 2023, documents a 73.8 percent probability that a
subsidy by one major economy (China, EU, or US) is matched by a subsidy in the same
product by another within 12 months. PPML regressions confirm that IP usage correlates
positively with RCA---governments target sectors where they already have comparative
advantage, consistent with political-economy capture rather than market-failure
correction. For our event study, this finding implies that the Teti GTD tariff
shock at time $t$ is not a clean one-off treatment but the first signal in an
escalating policy sequence: treatment intensity grows after the event-window, pushing
treatment effects toward permanence.

\citet{fontana2024} embed these observations in a political-economy theory of
strategic-assets policy, drawing on the three logics of capability blockade formalised
by Ding and Dafoe: cumulative-strategic blocking (targeting technological RCA through
standards exclusion and export controls), infrastructure-strategic blocking (targeting
critical material supply chains through reshoring subsidies), and dependency-strategic
blocking (exploiting existing dependencies through sanctions and FDI screening). The
taxonomy provides a structural basis for splitting the Teti tariff sample by motive:
tariff shocks in sectors matching the cumulative-strategic logic (advanced
manufacturing, semiconductors) should generate stronger $d_{\rm in}$ responses than
shocks in sectors matching the infrastructure-strategic logic (raw materials, chemicals),
where substitution is constrained by the CDI3 condition. We use this taxonomy in a
sample-split robustness exercise.

% ---------------------------------------------------------------------------
\subsection{Identification}
\label{subsec:lit_id}
% ---------------------------------------------------------------------------

Our identification strategy rests on a single organising principle: \emph{everything
that could mechanically respond to the shock must be predetermined relative to it.}
The estimand is the within-firm response of an export-capability portfolio to a
trade-policy shock, and the credibility of that estimand depends on three objects
being fixed before the shock---the exposure weights, the disruption elasticities, and
the capability map against which the portfolio is read. We discuss each, and then
state precisely which of three admissible network constructions we adopt and why.

\paragraph{A continuous, time-varying treatment.}
Following the input-tariff shift-share of \citet{amiti2007}, \citet{topalova2011}, and
\citet{tomasi2022}, we construct the tariff exposure variable as
\[
  \text{TariffShock}_{it}
  = \sum_p s_{ip,\text{pre}}\cdot \varepsilon_p \cdot \Delta\ln(1+\tau_{pt}),
\]
where $s_{ip,\text{pre}}$ are firm--product export shares in a strictly pre-shock
baseline window, $\Delta\ln(1+\tau_{pt})$ are changes in bilateral applied tariffs
from the Teti GTD, and $\varepsilon_p$ are the product-level trade elasticities of
\citet{fontagne2022}. We treat $\text{TariffShock}_{it}$ as a \emph{continuous,
time-varying, non-absorbing} dose rather than a binary staggered treatment: the dose
grows and recedes as tariffs move, and discretising it into a treated/control
indicator would discard precisely the cross-firm variation in exposure intensity that
identifies the dose--response. This choice is deliberate and consequential, as we make
clear below. The elasticities $\varepsilon_p$ are estimated by \citet{fontagne2022} on
the universe of 152 importers and 189 exporters---they are a property of the
\emph{product}, external to our firm sample, and hence cannot be a function of the
behaviour we study; products with a non-significant estimate inherit their HS-section
mean. The predetermined shares satisfy the Bartik exogeneity condition, which we defend
on \emph{both} of its routes: from the shares, following \citet{goldsmithpinkham2020},
and from the shocks, following \citet{borusyakhull2022}, reporting Rotemberg-weight
diagnostics and a recentred-instrument specification.

\paragraph{Dynamics and inference for a continuous dose.}
Because the treatment is continuous and non-absorbing, the cohort$\times$relative-time
machinery of \citet{sunabraham2021} is not the natural primary specification---its
cohorts are defined by the timing of a binary switch. We therefore estimate the
dynamics through a distributed-lag event study, $Y_{it}=\sum_k \beta_k\,
\text{TariffShock}_{i,t-k}+\alpha_i+\delta_t+\varepsilon_{it}$, in which the lead
coefficients ($k<0$) constitute the parallel-trends test and the lag coefficients
($k>0$) trace the persistence (hysteresis) prediction. Inference follows the
shift-share-robust standard errors of \citet{adao2019}, which are native to continuous
exposure designs and which we report as primary against firm- and sector-clustered
alternatives. As robustness we re-estimate the dynamics with the continuous,
non-absorbing estimator of \citet{dechaisemartin2020}, the continuous-treatment
difference-in-differences framework of \citet{callawaygoodmanbacon2024}, and---after
discretising exposure into terciles---the heterogeneity-robust estimators of
\citet{sunabraham2021}, \citet{callawaysantanna2021}, and \citet{borusyak2024}. We
flag the assumption that the continuous design carries a \emph{stronger} parallel-trends
requirement than the binary case---\citet{callawaygoodmanbacon2024} show that
identifying the dose--response slope requires the absence of selection into dose level,
which we probe but cannot fully test.

\paragraph{The capability map: three admissible constructions, and the one we adopt.}
The outcomes $d_{\rm in}$, $d_{\rm out}$, and $\text{Share\_core}$ are read off a
block partition of the firm--product export network \citep[the export, not production,
basis of][]{destefano2025}, obtained by bipartite modularity maximisation
\citep{barber2007} on an RCA-binarised incidence matrix validated against the Bipartite
Configuration Model \citep{squartini2011,saracco2015,cimini2022}. The non-trivial
question is \emph{when} the partition is estimated, because the map is itself the
measuring instrument for the outcome, and a map estimated on post-shock data would
co-move with the treatment. Three constructions are admissible:
\begin{enumerate}
  \renewcommand{\labelenumi}{(\roman{enumi})}
  \item \textbf{Frozen baseline partition.} The partition $b(p)$ and each firm's core
        block $c(f)$ are estimated once, on the pre-shock baseline window, and held
        fixed; only the firm's RCA incidence $M_{fpt}$ varies over time. The
        panel variation in $d_{\rm in}/d_{\rm out}$ is then unambiguously the firm
        relocating against a frozen, predetermined grid.
  \item \textbf{Predetermined evolving partition.} The map is allowed to evolve, but
        estimated only on \emph{pre-shock} slices through the temporal multilayer
        modularity of \citet{mucha2010}---with inter-slice coupling and resolution
        chosen by partition stability \citep{bazzi2016}---and then extrapolated forward.
        This nests (i) as the strong-coupling limit and captures any pre-existing drift
        in the capability space, at the cost of measurement error that grows with the
        post-shock horizon.
  \item \textbf{Contemporaneous multilayer.} The map evolves using all years including
        post-shock. Under our continuous treatment this construction is \emph{not}
        causally identified: with no never-treated firms, there is no
        treatment-independent control group from which to anchor the map, so the
        partition---and hence the $(d_{\rm in}+d_{\rm out})$ denominator---remains a
        function of the dose. We therefore use it only descriptively, to document how
        the capability space reorganises, and never as the causal outcome.
\end{enumerate}
\textbf{We adopt (i) as the headline specification}, precisely because it has the
smallest researcher degrees of freedom and the cleanest exogeneity argument: the map
literally precedes the shock and cannot respond to it. Its one weakness---that a frozen
map grows stale---generates classical measurement error and therefore attenuation
\emph{against} our hypotheses, an honest bias rather than a confounding one. We
neutralise even that concern by reporting (ii) as robustness, which lets the map evolve
on pre-shock information and leaves the estimates intact, and by re-estimating (i) with
a leave-one-out partition that excludes each focal firm from the map that classifies
it, breaking the reflexive channel. Two construction rules are stated explicitly and
checked: products with no baseline block are assigned by product-space proximity (or
dropped, with the affected export share reported), and firms born after the baseline
are handled within the exit-selection analysis rather than assigned a post-baseline
core. We compute RCA with the standard contemporaneous, \emph{within-China} Balassa
denominator: because the capability blocks are themselves built from the Chinese
firm--product network, a China-relative benchmark---specialisation relative to the
Chinese export economy---is the internally consistent and economically natural object,
and we retain it as the headline measure. The denominator carries a second-order
\emph{inference} concern rather than an economic one: under a sector-wide shock the
aggregate Chinese export of a treated product falls for all exposed firms, mechanically
inflating the within-sample benchmark and potentially moving threshold crossings for
reasons unrelated to the firm's own portfolio choice. We treat this as robustness, not
as a change of measure. We re-estimate with a leave-one-out within-China denominator,
which preserves the China-relative interpretation while removing the firm's reflexive
contribution, and, as a more aggressive check, against a rest-of-world (ex-China) BACI
benchmark, which trades the China referent for full exogeneity to the Chinese firms'
treated response.

\paragraph{An elegant alternative under a modified treatment.}
We note one design under which the \emph{contemporaneous} evolving map (iii) would
become causally admissible: were we to trade the continuous dose for a binary,
staggered treatment, the not-yet-treated firms would furnish a treatment-independent
control group from which the time-evolving partition could be estimated at each date
\citep[in the spirit of the control-group construction of][]{callawaysantanna2021},
yielding a map that is simultaneously evolving and exogenous. This is a genuine
optimisation of elegance over the frozen design, but it is purchased by discarding the
exposure-intensity information that motivates our continuous specification, and it
relies on a credible never-/not-yet-treated set under the diffuse Teti tariff
schedule. We therefore retain the continuous treatment with the frozen map as our
primary design, and present the binary control-group-map construction as a deliberate
robustness fork rather than the main specification.

\paragraph{Two concessions we make at the outset.}
Neither threat is specific to the network construction, and we address both head-on
rather than conceal them. First, the shock induces firm exit \citep{crozet2021}; since
an exited firm has no $d_{\rm in}/d_{\rm out}$, survivors are selected, and we therefore
estimate the effect on the probability of exit, bound the survivor effect, and read our
portfolio estimates as conservative. Second, a tariff on one product reallocates demand
to substitutes held by other firms in the same block, so the stable-unit-treatment
assumption is strained; the \citet{adao2019} correction absorbs the resulting
cross-firm correlation in the standard errors but not the causal interference, which we
discuss and, where feasible, probe with leave-out-block exposure.

\paragraph{Firm-level data and the ownership moderator.}
The panel of Chinese exporting firms comes from the NBS Annual Survey of Industrial
Enterprises, used in the same form by \citet{guariglia2011}. That paper establishes
three features of the NBS data that shape our design: SOEs operate under soft budget
constraints and face negligible financial frictions, private domestic firms are
severely financially constrained (near-unit-root self-financing), and coastal firms
without foreign capital face the binding constraint. Because firms switch ownership
status over the panel, we assign type by a consistent across-years rule following
\citet{guariglia2011}. The ownership heterogeneity motivates our SOE--private split:
privately owned firms facing financial constraints respond more sharply to export market
closures because they cannot subsidise loss-making product lines from retained earnings,
making the $d_{\rm out}$ contraction more immediate and the $d_{\rm in}$ consolidation
more forced.

% ---------------------------------------------------------------------------
\subsection{Our Contribution}
\label{subsec:contribution}
% ---------------------------------------------------------------------------

This paper contributes to all five literatures described above. To the economic-complexity
literature, we provide the first causal identification of a firm-level capability
response---$d_{\rm in}$, $d_{\rm out}$, and $\text{Share\_core}$---to an exogenous external shock,
demonstrating that the block-modular firm--product partition is not merely a descriptive
device but a meaningful predictor of differential adjustment paths. To the capability-theory
and evolutionary-economics literatures, we offer empirical content for the
exploitation--exploration tension formalised by \citet{march1991}: the direction of
adjustment under tariff shocks is shown to be systematically moderated by regional
strategic dependence, consistent with the Penrose--Teece prediction that capability
redeployment is constrained by the jointness of underlying knowledge inputs, while the
magnitude is consistent with the permanent-effect predictions of \citet{storper1992}
and the lock-in mechanism documented in the Schumpeterian Mark~II literature.

To the trade literature, we shift the outcome from performance (the productivity,
markup, and employment measures that dominate the Pavcnik--Topalova--Amiti generation) to
capability architecture, providing evidence that the same tariff shocks that raise
aggregate TFP through competitive selection simultaneously compress the capability
frontier of exposed firms in ways that the performance lens does not reveal. To
the strategic-dependency literature, we take the Farrell--Newman chokepoint
hypothesis to firm-level data: firms operating in CDI-dependent product
categories display systematically larger capability consolidation effects than firms
in non-dependent categories, and this heterogeneity is not explained by tariff
magnitudes alone, consistent with the dual treatment of higher tariffs and active
EU reshoring policy targeting the same product space. Finally, to the identification
literature, we demonstrate that a continuous-dose distributed-lag event study---with the
heterogeneity-robust continuous-treatment estimators of \citet{dechaisemartin2020} and
\citet{callawaygoodmanbacon2024} as robustness---is tractable at the scale of
full-universe firm--product panels, that the \citet{adao2019} shift-share-robust
correction materially affects standard errors in heterogeneous-exposure settings with
product-level clustering, and that reading a network-derived outcome against a
\emph{predetermined} (frozen, or pre-shock-estimated) capability partition is what keeps
the estimand exogenous to the shock.

% =============================================================================
\section{Expected Dynamics: What the Literature Predicts We Should Observe}
\label{sec:expected}
% =============================================================================

Because the empirical estimates are not yet in hand, it is worth making the priors
explicit. This section reads the literatures of Section~\ref{sec:lit} \emph{forward}:
for each outcome and each moderator, it states what dynamic the prior work leads us to
expect once the planned design is executed, what sign each hypothesis carries, and
what magnitude range is plausible. The aim is to specify in advance what a result
``consistent with the literature'' would look like---so that a surprising estimate can
be recognised as surprising---and to organise the diverse predictions into a single
testable structure. Nothing here is a finding; everything here is a prediction that
the data will confirm, refute, or qualify.

% ---------------------------------------------------------------------------
\subsection{The Out-of-Block Margin ($d_{\rm out}$)}
\label{subsec:exp_dout}
% ---------------------------------------------------------------------------

Three literatures converge on a contraction of $d_{\rm out}$ following an adverse
tariff shock. The sunk-cost extensive-margin mechanism of \citet{crozet2021} predicts
that exposed firms stop serving the products and markets that become tariff-burdened,
and---because re-entry costs are 19--23 percent higher for lapsed exporters---the
contraction is \emph{hysteretic}: it should not reverse within the event window even
if tariffs are later rolled back. The variance-cost mechanism of \citet{crozet2025}
and \citet{vannoorenberghe2016} predicts that firms which try to \emph{preserve}
$d_{\rm out}$ by pivoting to unfamiliar, structurally dissimilar destinations inherit
higher demand volatility, so the privately optimal response for all but the largest
firms is to let $d_{\rm out}$ fall rather than chase replacement markets. And the
diversification-discount search models of \citet{matsusaka2001} and \citet{bernardo2002}
predict that only firms whose core match remains above the liquidation threshold retain
the option to explore; below it, exploratory $d_{\rm out}$ search collapses. The net
expectation is a \emph{negative} average effect on $d_{\rm out}$, concentrated among
smaller, financially constrained, and inexperienced firms, and persistent over the
post-shock horizon. The one literature that cuts the other way is the escape-competition
margin of \citet{aghion2005} and \citet{akcigit2022}: a high-capability firm facing
competitive pressure may innovate \emph{outward}, lifting $d_{\rm out}$. Because that
channel is confined to frontier firms, the expectation is a negative average masking a
positive tail---visible only in the interaction with firm capability rank, not in the
mean.

% ---------------------------------------------------------------------------
\subsection{The In-Block Margin ($d_{\rm in}$) and Retrenchment}
\label{subsec:exp_din}
% ---------------------------------------------------------------------------

The prediction for $d_{\rm in}$ is the cleaner of the two. The Typical-Product-Vector
consolidation of \citet{fontagne2018}---a 7.7 percentage-point retention premium for
core products in treated sectors---is, mechanically, a $d_{\rm in}$ response to a
trade-policy shock, and it is the single closest empirical antecedent for our main
outcome. The scope-economy variance argument of \citet{bottazzi2006} predicts that a
firm rebalancing a disrupted portfolio minimises revenue covariance by deepening within
a related cluster rather than scattering, which raises $d_{\rm in}$ relative to
$d_{\rm out}$. The competition-to-coherence evidence of \citet{valvano2003}---firms in
single-market-exposed sectors became more coherent over time---predicts the same sign.
And the Schumpeterian Mark~II lock-in of \citet{dosi1995} predicts that the retrenchment
is strongest exactly where capabilities are most cumulative. The expectation is
therefore a \emph{positive} average effect on $d_{\rm in}$, larger for firms in
complex, cumulative technological regimes, and---per \citet{storper1992} and
\citet{riverabatiz1991}---persistent rather than transitory, because what is lost when
a market closes is more than demand: it is the collective learning infrastructure that
sustained the broader basket.

% ---------------------------------------------------------------------------
\subsection{Share\_core and the Direction of Adjustment}
\label{subsec:exp_share}
% ---------------------------------------------------------------------------

The summary statistic $\text{Share\_core} = d_{\rm in}/(d_{\rm in}+d_{\rm out})$ signs
the exploitation--exploration choice of \citet{march1991}. If the $d_{\rm out}$
contraction and the $d_{\rm in}$ consolidation predicted above both materialise,
$\text{Share\_core}$ rises---the retrenchment hypothesis (H1). March's own result that
organisations systematically over-exploit makes this the modal expectation. But the
literature is deliberately not unanimous: the dynamic-capabilities reading
(\citealp{penrose1959}; \citealp{teece1982}) and the moderate-shock search prediction
of \citet{matsusaka2001} leave open the exploration case (H1$'$), in which $\text{Share\_core}$
\emph{falls} because firms with slack treat the shock as a prompt to reallocate toward
new capability bundles. The empirical question is therefore which force dominates on
average, and, more informatively, where in the firm distribution the sign flips. The
strong prior, following \citet{destefano2025}, is that the average response is toward
the growth-\emph{negative} core (H1), because the growth-positive $d_{\rm out}$ margin
is precisely the one that sunk costs and variance costs make hardest to defend under
stress.

% ---------------------------------------------------------------------------
\subsection{Heterogeneity and the Strategic-Dependence Moderator}
\label{subsec:exp_het}
% ---------------------------------------------------------------------------

The sharpest predictions are conditional, not average. The Farrell--Newman chokepoint
logic \citep{farrell2019}, operationalised through the CDI product set of
\citet{arjona2023} and the network-fragility tiers of \citet{korniyenko2017}, predicts
that firms whose baskets concentrate in strategically dependent products face not a
neutral tariff but a tariff \emph{plus} active reshoring policy aimed at the same
products---so their retrenchment should be amplified and, given the escalation dynamic
of \citet{evenett2024}, more persistent. The financial-constraint heterogeneity of
\citet{guariglia2011} predicts that private Chinese firms retrench more sharply than
soft-budget SOEs. The size-and-volatility composition effect of
\citet{vannoorenberghe2016} predicts that small firms cannot afford to defend
$d_{\rm out}$ at all. The institutional-contingency result of \citet{delios2008} and
\citet{boschma2015} predicts that the related-versus-unrelated payoff itself shifts
with the regional institutional environment, so that ExpDep should be read as a
production-system attribute and not only a firm attribute. The blockade-logic taxonomy
of \citet{fontana2024} predicts that cumulative-strategic sectors retrench harder than
infrastructure-strategic ones. The unifying expectation is that \emph{heterogeneity
dominates the average}: a modest mean effect concealing a steep gradient in strategic
dependence, financial constraint, and capability rank.

% ---------------------------------------------------------------------------
\subsection{Synthesis: Hypotheses, Signs, and Magnitude Benchmarks}
\label{subsec:exp_synth}
% ---------------------------------------------------------------------------

Table~\ref{tab:hyp} collects the hypotheses, their literature sources, and the sign
each predicts. Table~\ref{tab:mag} brackets the plausible magnitude range using two
external benchmarks: a lower bound from pure network-disruption transmission
(\citealp{korniyenko2017}: a 1pp rise in risky imports from a disaster-struck origin
cuts the importer's own exports by $\sim$0.7 percent, with no policy response) and an upper
bound from the most exposed economies in a targeted tariff war (\citealp{teti2025}:
$-1.6$ percent cumulative real income for Canada and Mexico under the 2018--2025 US tariffs,
versus $-0.3$ percent for China in the aggregate---a country-level benchmark we map onto the
most exposed Chinese firms).

\begin{table}[ht]
\centering
\caption{Hypotheses, literature sources, and predicted signs.}
\label{tab:hyp}
\footnotesize
\setlength{\tabcolsep}{4pt}
\renewcommand{\arraystretch}{1.3}
\begin{tabular}{|c|p{4.4cm}|p{4.2cm}|p{2.5cm}|}
\hline
\textbf{Hyp.} & \textbf{Statement} & \textbf{Principal literature source} & \textbf{Predicted effect} \\
\hline
H1 & Adverse tariff shock raises $\text{Share\_core}$ / $d_{\rm in}$ (retrenchment to the coherent core). &
\citet{fontagne2018}; \citet{bottazzi2006}; \citet{valvano2003}; \citet{dosi1995}; \citet{march1991} &
$d_{\rm in}\uparrow$, $\text{Share\_core}\uparrow$ \\
\hline
H1$'$ & Adverse shock raises $d_{\rm out}$ / lowers $\text{Share\_core}$ (exploration beyond the core). &
\citet{penrose1959}; \citet{teece1982}; \citet{matsusaka2001}; \citet{aghion2005} &
$d_{\rm out}\uparrow$ (frontier-firm tail only) \\
\hline
H2 & The retrenchment response is amplified in more strategically dependent production systems. &
\citet{farrell2019}; \citet{arjona2023}; \citet{korniyenko2017}; \citet{evenett2024} &
$\beta_{2}>0$ on Tariff $\times$ ExpDep \\
\hline
H3 & The $d_{\rm out}$ contraction is hysteretic (does not reverse within the window). &
\citet{crozet2021}; \citet{storper1992}; \citet{riverabatiz1991} &
persistent, asymmetric \\
\hline
H4 & Responses are steeper for financially constrained / private / small firms. &
\citet{guariglia2011}; \citet{vannoorenberghe2016}; \citet{coad2021} &
steeper gradient \\
\hline
\end{tabular}
\end{table}

\begin{table}[ht]
\centering
\caption{Magnitude benchmarks for the $d_{\rm out}$ response (per s.d.\ of TariffShock).}
\label{tab:mag}
\footnotesize
\setlength{\tabcolsep}{4pt}
\renewcommand{\arraystretch}{1.3}
\begin{tabular}{|p{3.8cm}|c|p{6.2cm}|}
\hline
\textbf{Benchmark} & \textbf{Magnitude} & \textbf{Interpretation} \\
\hline
Lower bound --- pure network disruption \citep{korniyenko2017} & $\sim 0.7\%$ &
Supply-chain transmission of a localised shock, no policy intervention. \\
\hline
Upper bound --- most exposed economies, targeted tariff war \citep{teti2025} & $\sim 1.6\%$ &
Country-level loss for Canada/Mexico under the 2018--2025 US tariffs (vs.\ $-0.3\%$ for China), mapped onto the most exposed firms. \\
\hline
\end{tabular}
\end{table}

% =============================================================================
\section{Open Questions and a Forward Research Agenda}
\label{sec:agenda}
% =============================================================================

The review above leaves a small number of questions genuinely open and opens a larger
number of new ones. This section separates the two: first, the specific gaps in the
existing literature that \emph{this} paper is designed to close; second, the spin-off
questions that the same body of work raises but that lie beyond the present design and
would each support a distinct paper.

% ---------------------------------------------------------------------------
\subsection{The Gaps This Paper Fills}
\label{subsec:gaps_filled}
% ---------------------------------------------------------------------------

\begin{enumerate}
\item \textbf{Causal capability reconfiguration, not descriptive structure.} The EFC
literature (\citealp{bruno2023}; \citealp{pugliese2019}; \citealp{destefano2025})
measures block-modular capability structure but almost always descriptively. No prior
work links a \emph{predetermined, exogenous} trade-policy shock to the in/out
decomposition of firm diversification. This paper supplies that causal link.

\item \textbf{Composition, not performance.} The trade-and-productivity literature
(\citealp{pavcnik2002}; \citealp{topalova2011}; \citealp{amiti2007}; \citealp{tomasi2022})
measures how tariff shocks move productivity, markups, and survival. Because
performance dimensions decouple from capability (\citealp{grazzi2012};
\citealp{coad2021}; \citealp{syverson2011}), these outcomes are silent on \emph{which}
activities firms keep, drop, or enter. This paper moves the outcome to the architecture
of the export basket.

\item \textbf{Firm-level test of the chokepoint hypothesis.} The strategic-dependency
literature (\citealp{farrell2019}; \citealp{euc2021}; \citealp{arjona2023};
\citealp{korniyenko2017}) maps dependencies at the product and country level but has no
firm-level causal test. By interacting the tariff shock with a firm-specific
strategic-dependence exposure, this paper provides that test.

\item \textbf{A signed reading of exploration vs.\ exploitation under a real shock.}
March's (\citeyear{march1991}) tension has rich theory but little exogenous-shock
evidence at the firm level. The $\text{Share\_core}$ outcome under a predetermined tariff shock
delivers a directly signed estimate of which way firms move when their core is
threatened.

\item \textbf{Scalability of continuous-treatment event-study estimators.} We demonstrate
that a continuous-dose distributed-lag design with \citet{adao2019} shift-share inference---and
the heterogeneity-robust continuous-treatment estimators of \citet{dechaisemartin2020} and
\citet{callawaygoodmanbacon2024}---is tractable on a full-universe firm--product panel, a
methodological contribution in its own right for the heterogeneous-exposure shift-share
literature.
\end{enumerate}

% ---------------------------------------------------------------------------
\subsection{Beyond This Paper: Spin-Off Research Questions}
\label{subsec:spinoffs}
% ---------------------------------------------------------------------------

Each of the following is motivated by the reviewed literature but deliberately left
outside the present design; each would constitute a separate study.

\paragraph{(a) The firm--product--destination tensor.}
\citet{destefano2025} works in pure firm--product space; \citet{crozet2025} and
\citet{vannoorenberghe2016} show that the destination dimension carries its own
volatility structure. A natural extension promotes the bipartite firm--product matrix
to a firm--product--destination tensor and asks whether capability blocks reorganise
\emph{geographically}---whether retrenchment to the core is also a retreat to familiar
markets. This is the volatility-minimising hypothesis of $\Phi$-diversity made
firm-specific.

\paragraph{(b) Contemporaneous temporal-coupled partitions as a causal object.}
Our headline design freezes the baseline block, and our robustness lets it evolve on
\emph{pre-shock} information through the temporal multilayer machinery of \citet{mucha2010}
(with coupling and resolution chosen as in \citealp{bazzi2016}). What we deliberately leave
open is the \emph{contemporaneous} version: treating \emph{block reconfiguration itself}---the
birth, death, and merger of capability clusters measured with post-shock data---as the causal
outcome. Under our continuous treatment this is not identified, because there is no
never-treated control group from which to anchor an evolving map; a dedicated study could
recover it under a binary, staggered design by estimating the partition from not-yet-treated
firms, asking how trade shocks reshape the modular geometry of an entire industry rather than
the position of a single firm within a fixed geometry.

\paragraph{(c) The antidumping and safeguard stage.}
The Teti GTD \citep{teti2024} captures statutory MFN and preferential tariffs but
\emph{not} antidumping, safeguard, or trade-war tariffs. Pairing it with a temporary-trade-barriers
database would isolate the capability response to discretionary, firm-targeting
protection---plausibly a sharper and more anticipatory shock than scheduled tariff
changes, and a cleaner test of the escalation dynamic of \citet{evenett2024}.

\paragraph{(d) Instrumenting the policy-escalation component.}
\citet{evenett2024} document that industrial-policy intensity rises in election years
and that subsidies are matched tit-for-tat with 73.8 percent probability. An election-cycle
or tit-for-tat instrument for the policy-driven component of a firm's tariff shock
would separate the escalation channel from the mechanical tariff change---turning the
escalation caveat of this paper into an identification strategy of its own.

\paragraph{(e) Green-capability reconfiguration.}
\citet{mealy2022} show the complexity toolkit applies to any product subset and that
green-capability accumulation is strongly path dependent. Re-running the present design
with the green-product set as the treated category would ask whether decarbonisation
policy and trade policy push firms in the same or opposite capability directions---a
direct bridge between the strategic-autonomy agenda of \citet{cagnin2021} and the green
transition.

\paragraph{(f) Institutional voids as a moderator.}
\citet{delios2008} and \citet{boschma2015} show that the value of diversification, and
the related-versus-unrelated balance, depend on institutional quality. Exploiting
within-China regional institutional variation as a second moderator would test whether
capability retrenchment under trade shocks is steeper or shallower where market
institutions are weaker, distinguishing capability exploitation from
institutional substitution as motives for diversification.

\paragraph{(g) The gain/loss asymmetry of complexity.}
\citet{nomaler2023} show that lost specialisations are often \emph{more} complex than
gained ones, and that density-based relatedness is heavily confounded by diversity and
ubiquity. A study that scores firm-level product \emph{exits} by complexity---asking
whether tariff shocks make firms shed their most sophisticated products first---would
test whether trade shocks erode the capability frontier from the top, a sharper and
more policy-relevant question than net diversification counts can answer.
